\chapter{Capacity and quasicontinuity}
\index{capacity!BFJ}\index{quasicontinuity}
\index{negligible set}\index{pluripolar set!BFJ}
\index{convergence!in capacity}
\label{chap:bfj-capacity}

Weak convergence of measures tests only continuous functions, whereas the
sets naturally produced by psh functions are level sets of upper
semicontinuous ones.  Capacity is the device that bridges this gap.  It
declares a set small when \emph{every} normalized bounded Monge--Amp\`ere
measure assigns little mass to it, and it turns out that removing a set of
small capacity makes a bounded psh function continuous.  From that single
theorem---quasicontinuity---the maximum principle, the comparison principle
and the domination principle all follow, and with them the measure-theoretic
half of the variational method.

The chapter is a self-contained construction.  \Cref{sec:capacity-preview}
supplied the axiomatic calculus of an abstract capacity: monotonicity,
countable subadditivity, polar sets, quasicontinuous functions and their
stability.  Nothing there constructs a capacity, and nothing there is true of
a particular capacity beyond what the axioms give.  Here the axioms are
verified for one specific set function, the Monge--Amp\`ere capacity attached
to the bounded operator of \Cref{chap:bfj-ma}, and the analytic content is
supplied by two estimates: a Chern--Levine--Nirenberg bound
(\Cref{lem:ma-capacity-bound}) and a Cauchy--Schwarz bootstrap
(\Cref{lem:fundamental-capacity-estimate}).

\section{Orientation}
\label{sec:ch25-orientation}

\subsection{The gap that capacity fills}
\label{sec:ch25-gap}

\Cref{thm:bounded-ma-extension} produced, for every $n$-tuple of bounded
$\theta$-psh functions, a positive Radon measure of mass $V$, continuous
along decreasing nets when tested against continuous functions, and with
weighted monotone continuity for bounded psh weights.  Neither statement
by itself controls the mass of every Borel set.  Every statement one actually wants
concerns a set defined by an inequality between psh functions:
$\{\varphi>\psi\}$ in the maximum principle, $\{\varphi<\psi\}$ in the
comparison principle, $\{\varphi<-t\}$ in every energy estimate.  Such a
set need not be open or closed, and its indicator need not be continuous;
\Cref{rem:portmanteau-sharp} recorded that weak convergence says nothing
about the mass of a general Borel set.

There is a second, subtler gap.  A bounded $\theta$-psh function is upper
semicontinuous and is a decreasing limit of continuous model functions
(\Cref{thm:bfj-decreasing-regularization}), but it need not be continuous,
and the regularization theorem initially supplies a net.
The absence of a countably cofinal model system
(\Cref{ex:ch22-no-countable-cofinal}) prevents a simple diagonal model
argument.  The later capacity estimate supplies a sequence adapted to each
potential; it does not make the system of all models countably cofinal.

The capacity closes both gaps at once.  Write
\begin{equation}
\label{eq:ch25-admissible}
 \mathcal U_\theta
 :=\bigl\{u\in\PSH(X,\theta)\ :\ -1\leq u\leq0\bigr\}
\end{equation}
for the \emph{admissible family}.  This is a nonempty set---the constant
$0$ belongs to it, because $\theta$ is semipositive
(\Cref{prop:ch21-constants})---and it is convex and stable under finite
maxima.  Declaring
\[
              \Capa_\theta(E)=\sup_{u\in\mathcal U_\theta}\MA_\theta(u)(E)
\]
measures a set by the largest mass any normalized potential can put on it.
Three things then happen, and they are the three theorems of the chapter.

\begin{enumerate}
\item\label{itm:ch25-plan:qc}
 \emph{Quasicontinuity} (\Cref{thm:bounded-quasicontinuity}): every bounded
 $\theta$-psh function is continuous off an open set of arbitrarily small
 capacity, and the model regularization converges uniformly there.  The
 exceptional set is produced, not merely asserted to exist, and the
 approximants along which it is produced form a \emph{sequence}
 (\Cref{cor:psh-sequential-regularization})---a genuine improvement on
 \Cref{rem:ch22-decreasing-selection-improvement}, which explicitly declined
 to produce one.
\item\label{itm:ch25-plan:loc}
 \emph{Locality} (\Cref{thm:bounded-maximum-principle}): the bounded
 Monge--Amp\`ere measure of a maximum agrees with that of the larger factor
 on the set where the latter is larger.  This is the statement that the
 operator is a differential operator and not merely a weak limit.
\item\label{itm:ch25-plan:cmp}
 \emph{Comparison and domination}
 (\Cref{cor:bounded-comparison-principle},
 \Cref{prop:bounded-capacity-domination}): the sublevel set
 $\{\varphi<\psi\}$ carries at least as much $\MA_\theta(\varphi)$-mass
 as $\MA_\theta(\psi)$-mass, and its capacity is controlled by the former.
\end{enumerate}

\subsection{The five results, stated in advance}
\label{sec:ch25-payoff}

The reader who wants the statements before the constructions may read the
following five displays and then turn to \Cref{sec:ch25-worked}, where each
of them is exercised on explicit numbers.

\emph{(i) The capacity is monotone, countably subadditive, continuous along
monotone families in the precise sense of
\Cref{prop:ch25-choquet-axioms}, and it sees every nonempty open set.}
\Cref{prop:capacity-elementary},
\Cref{prop:ch25-choquet-axioms} and \Cref{lem:divisorial-positive-capacity}:
for Borel sets,
\[
 \begin{gathered}
 0\leq\Capa_\theta(E)\leq V,
 \qquad
 \Capa_\theta\Bigl(\bigcup_j E_j\Bigr)\leq\sum_j\Capa_\theta(E_j),\\
 \Capa_\theta(G)>0\quad\text{for every nonempty open set }G.
 \end{gathered}
\]

\emph{(ii) Bounded psh functions are quasicontinuous, uniformly along a
decreasing sequence of model functions.}
\Cref{thm:bounded-quasicontinuity}.

\emph{(iii) Polar sets are negligible, with an explicit rate.}
\Cref{prop:ch25-sublevel-capacity} and
\Cref{cor:ch25-polar-negligible}: for $u\in\PSH(X,\theta)$ with
$\sup u=0$,
\[
 \Capa_\theta\bigl(\{u<-t\}\bigr)\leq\frac{C_\theta+nV}{t},
 \qquad\text{hence}\qquad
 \Capa_\theta\bigl(\{u=-\infty\}\bigr)=0 ,
\]
where $C_\theta$ is the explicit constant of
\Cref{lem:ch25-normalized-bound}.  The converse implication, from negligible
to polar, is \emph{not} proved in this book.

\emph{(iv) Uniform convergence implies convergence in capacity; pointwise
convergence and convergence in capacity do not imply each other for the
whole sequence.}  \Cref{prop:decreasing-capacity-convergence},
\Cref{ex:ch25-capacity-not-uniform} and
\Cref{ex:ch25-pointwise-not-capacity}.

\emph{(v) The comparison principle and its capacity form.}
\Cref{cor:bounded-comparison-principle} and
\eqref{eq:capacity-comparison}: for bounded $\theta$-psh
$\varphi,\psi$,
\[
 \int_{\{\varphi<\psi\}}\MA_\theta(\psi)
 \leq\int_{\{\varphi<\psi\}}\MA_\theta(\varphi),
\]
and, when $\psi\leq0$ and $0<t<1$,
$\Capa_\theta\{\varphi<\psi\}
 \leq t^{-n}\MA_\theta(\varphi)\{\varphi<(1-t)\psi+t\}$.

\subsection{Standing hypotheses}
\label{sec:ch25-hypotheses}

The original BFJ signature \eqref{eq:ch20-standing} is in force throughout:
$K$ is complete and discretely valued with residue field of characteristic
zero, $X$ is smooth, connected and projective over $K$ of dimension
$n\geq1$, and models are normal, flat and projective over $R=K^\circ$.
In addition we fix once and for all a closed model $(1,1)$-form
\begin{equation}
\label{eq:ch25-reference-form}
 \theta\in\mathcal Z^{1,1}(X),
 \qquad\text{$\theta$ semipositive with ample de Rham class,}
 \qquad V:=(\{\theta\}^n)>0 .
\end{equation}
Semipositivity is used constantly and is not decorative: it is what makes the
constants $\theta$-psh (\Cref{prop:ch21-constants}), hence what makes the
admissible family \eqref{eq:ch25-admissible} nonempty, and it is what makes
every mixed measure in sight positive of mass $V$
(\Cref{prop:model-ma-basic}, \eqref{eq:mixed-total-mass}).  Ampleness of the
de Rham class is used through
\Cref{prop:ch20-positive-representative} whenever a strictly positive
representative is required, and through
\Cref{thm:bfj-skeletal-equicontinuity} in every estimate that quantifies over
all $\theta$-psh functions at once.

Following \cite[\S\S4--5]{BFJ15} one may normalize $V=1$; we do not, and
every constant below displays the power of $V$ that it carries.  The
verification that these powers are the ones forced by
\Cref{app:normalizations} is \Cref{sec:ch25-scope}.

\subsection{The bounded regime is a hypothesis, not a convenience}
\label{sec:ch25-bounded-regime}

The analytic estimates and the locality and comparison principles of this
chapter are statements about \emph{bounded} $\theta$-psh functions, and the
boundedness enters in two distinct ways which it is worth separating now.

First, the Monge--Amp\`ere operator itself is only defined for bounded
potentials at this stage of the book: \Cref{thm:bounded-ma-extension}
constructs $\MA_\theta(\varphi_1,\ldots,\varphi_n)$ for bounded
$\varphi_i$ and for no others.  The non-pluripolar product
\Cref{def:nonpluripolar-ma} and the finite-energy class $\cE^1$ are built
in \Cref{chap:bfj-energy}, after this chapter and using it.

Second, the boundedness enters quantitatively.  The constant
$M$ in \Cref{lem:ma-capacity-bound} and in
\Cref{lem:fundamental-capacity-estimate} is a bound for the potentials
involved, and it appears in the conclusion; a family of potentials with
unbounded oscillation is not controlled by any of these estimates, no matter
how well behaved its energies are.

Accordingly we adopt the following discipline.  The locality, comparison,
and moving-measure statements concern bounded potentials.  Unbounded
$\theta$-psh functions enter only through the sublevel estimate
\Cref{prop:ch25-sublevel-capacity}, the sequential regularization
\Cref{cor:psh-sequential-regularization}, and the truncation argument of
\Cref{sec:ch25-e1}.  In each case every Monge--Amp\`ere measure that occurs
still has bounded potentials as entries.  In particular no estimate proved
with a bound $M$ is later applied to a member of $\cE^1$ without first
truncating it.

\begin{warning}{Two capacities, two regimes}
It is tempting to define a capacity by taking the supremum over
\emph{all} $\theta$-psh $u$ with $\sup u=0$, or over all
$u\in\cE^1$ with bounded energy.  Neither variant is used here.  The first
is unavailable because $\MA_\theta(u)$ is not defined for unbounded $u$ at
this stage; the second is a genuinely different set
function, adapted to \Cref{chap:bfj-finite-energy-measures}, and none of the
estimates below is asserted for it.  Whenever the symbol $\Capa_\theta$
occurs in this book it means \eqref{eq:capacity-definition}, with the
admissible family \eqref{eq:ch25-admissible} of \emph{bounded} potentials.
\end{warning}

\subsection{What is imported}
\label{sec:ch25-imported}

Everything in this chapter is proved, with the following exceptions, each of
which is quoted with its hypotheses at the point of use.

\begin{enumerate}
\item\label{itm:ch25-import:ma}
 \emph{The bounded Monge--Amp\`ere operator and its weighted monotone
 continuity}, \Cref{thm:bounded-ma-extension}, together with
 multilinearity, integration by parts and the Cauchy--Schwarz inequality for
 differences of bounded psh functions,
 \Cref{cor:bounded-ma-multilinearity}.  These are the analytic input; the
 present chapter adds no new construction of measures.
\item\label{itm:ch25-import:model}
 \emph{The model maximum principle} \Cref{prop:model-maximum-principle},
 which is the combinatorial seed of
 \Cref{thm:bounded-maximum-principle}, and the model mixed measure
 \eqref{eq:model-mixed-measure} with its total mass
 \eqref{eq:mixed-total-mass}.
\item\label{itm:ch25-import:reg}
 \emph{Decreasing regularization} \Cref{thm:bfj-decreasing-regularization}:
 every $\theta$-psh function is the limit of a decreasing net of
 $\theta$-psh model functions.  This is the theorem whose proof needs
 residue characteristic zero, and it is used in every proof below that
 begins ``choose a decreasing model net''.
\item\label{itm:ch25-import:equi}
 \emph{Skeletal equicontinuity} \Cref{thm:bfj-skeletal-equicontinuity} and
 the identification \Cref{prop:ch21-sup-on-skeleton} of the supremum of a
 $\theta$-psh function with its supremum on a dual complex.  These give
 the only uniform bound in the chapter that does not come from an
 assumption, namely \Cref{lem:ch25-normalized-bound}.
\item\label{itm:ch25-import:tools}
 \emph{The abstract capacity calculus} of \Cref{sec:capacity-preview}: the
 $\sigma$-ideal property of polar sets
 (\Cref{prop:polar-sigma-ideal}), the stability of quasicontinuous
 functions (\Cref{prop:quasicontinuity-calculus}), and the sequential
 extraction (\Cref{prop:capacity-subsequence}).  The hypotheses of those
 statements---countable subadditivity, in each case---are verified here in
 \Cref{prop:capacity-elementary}.
\end{enumerate}

No result of \Cref{part:cld} is used.  In particular
\Cref{thm:cld-polar-negligibility}, the Chambert--Loir--Ducros negligibility
theorem, is a local statement about a different product and is not being
transported; the relation between the two is discussed in
\Cref{rem:ch25-cld-comparison}.

\subsection{Nets, sequences, and where the reduction is legitimate}
\label{sec:ch25-nets}

The space $X^{\An}$ is compact Hausdorff; no metrizability assumption is
made.  In the model system illustrated in
\Cref{ex:ch22-no-countable-cofinal}, no countably cofinal family is available
(\Cref{lem:ch21-no-cofinal-sequence}).  Convergence statements in this chapter
are therefore formulated for nets, in accordance with
\Cref{rem:net-vs-sequence-checklist}, with exactly three exceptions, each of
which is a theorem rather than an abbreviation.

\begin{enumerate}
\item\label{itm:ch25-seq:chebyshev}
 In \Cref{thm:bounded-quasicontinuity} a \emph{sequence} of indices
 $\alpha_1\leq\alpha_2\leq\cdots$ is extracted from a decreasing model
 net.  The extraction is legitimate because the quantity being made small,
 $\sup_{u\in\mathcal U_\theta}\int(\varphi_\alpha-\varphi)\MA_\theta(u)$,
 tends to $0$ along the net, so each of the countably many required
 conditions is satisfied on a cofinal subset; directedness then produces a
 single increasing sequence of indices.  No cofinal sequence in the model
 net is claimed or produced: the sequence $(\varphi_{\alpha_m})$ is
 cofinal for the purposes of the capacity estimate only.
\item\label{itm:ch25-seq:diag}
 In \Cref{cor:psh-sequential-regularization} the same device, run with a
 diagonal choice of truncation level, produces a decreasing
 \emph{sequence} of model functions converging pointwise to a given psh
 function.  This is a genuine strengthening of
 \Cref{thm:bfj-decreasing-regularization}, and it is available only because
 divisorial points have positive capacity
 (\Cref{lem:divisorial-positive-capacity}) and $\theta$-psh functions are
 determined by their values there (\Cref{prop:ch21-divisorial-order}).
\item\label{itm:ch25-seq:subseq}
 \Cref{prop:capacity-subsequence}, quoted from
 \Cref{chap:topological-tools}, extracts from a \emph{sequence} converging
 in capacity a subsequence converging outside a negligible set.  It is
 stated for sequences, and no subnet conclusion is inferred from it; see
 \Cref{rem:ch25-nets-in-capacity}.
\end{enumerate}

\subsection{Plan}
\label{sec:ch25-plan}

\Cref{sec:ch25-capacity} defines the capacity and proves its elementary
properties, including the Choquet axioms and the positivity of the capacity
of a nonempty open set.  \Cref{sec:ch25-estimates} proves the two analytic
estimates on which everything else rests.  \Cref{sec:ch25-negligible}
compares the three notions of smallness---negligible, polar, pluripolar---and
records precisely which implications are proved.
\Cref{sec:ch25-quasicontinuity} proves quasicontinuity and the sequential
regularization.  \Cref{sec:ch25-convergence} develops convergence in capacity
and separates it from the other convergences by examples.
\Cref{sec:ch25-maximum} and \Cref{sec:ch25-comparison} prove the maximum,
comparison and domination principles.  \Cref{sec:ch25-e1} passes from the
bounded theorem to arbitrary extended-real-valued psh functions by
truncation and records the sharper finite-energy tail available in the next
chapter.  \Cref{sec:ch25-curves} computes capacity on
a metric graph, and \Cref{sec:ch25-worked} carries five examples to explicit
numbers.  \Cref{sec:ch25-scope} is the scope table and
\Cref{sec:ch25-dictionary} the chapter dictionary.

\section{The Monge--Amp\`ere capacity}
\label{sec:ch25-capacity}

\subsection{Definition and normalizations}
\label{sec:ch25-definition}

\begin{definition}[Capacity]
\label{def:bfj-capacity}
For a Borel set $E\subset X^{\An}$, define
\begin{equation}
\label{eq:capacity-definition}
 \Capa_\theta(E)
 :=\sup\left\{\int_E\MA_\theta(u)\ \middle|\
       u\in\PSH(X,\theta),\ -1\leq u\leq0\right\}.
\end{equation}
An arbitrary set is assigned its outer capacity, the infimum of
$\Capa_\theta(G)$ over open $G\supset E$. A set of outer capacity zero
will be called \emph{negligible}.
\end{definition}

Three points of the definition deserve comment before anything is proved
with it.

\emph{The supremum is over the admissible family
\eqref{eq:ch25-admissible}, not over a larger class.}  Each
$\MA_\theta(u)$ with $u\in\mathcal U_\theta$ is a positive Radon measure
of total mass $V$ by \Cref{thm:bounded-ma-extension}, so the supremum is a
supremum of numbers in $[0,V]$ and is finite without any argument.  This is
the first and most elementary reason for restricting to bounded potentials.

\emph{The definition is invariant under translation of the admissible
family.}  If $a\in\mathbf R$ then $u\mapsto u+a$ is a bijection from
$\mathcal U_\theta$ onto $\{u\in\PSH(X,\theta):a-1\leq u\leq a\}$, and
$\MA_\theta(u+a)=\MA_\theta(u)$ because $\ddc$ kills constants
(\Cref{def:ch20-ddc}).  We use this freedom constantly, in the form recorded
in \Cref{lem:ch25-admissible-calculus}\,\cref{itm:ch25-adm:transl}, to compare
\eqref{eq:capacity-definition} with a supremum over potentials normalized by
$0\leq u\leq1$.

\emph{Outer regularity is imposed, not proved.}  The two-step definition
assigns to a Borel set $E$ both the number \eqref{eq:capacity-definition}
and the outer number $\inf\{\Capa_\theta(G):G\supseteq E\ \text{open}\}$.
These need not coincide: the supremum over $u$ and the infimum over $G$
occur in the wrong order, and no argument in this book interchanges them.  To
avoid ambiguity we write $\Capa_\theta^*$ for the outer capacity when the
distinction matters, and record in \Cref{lem:ch25-inner-outer} the two cases
in which the distinction does not matter: $E$ open, and $E$ of outer
capacity zero.  When an outer-capacity conclusion is needed below, we
construct open majorants with small capacity; an estimate for the Borel-set
capacity alone is not substituted for this step.  The reader should not read
``$\Capa_\theta(E)=0$'' for a general Borel $E$ as if outer regularity
had been established.

\begin{lemma}
\label{lem:ch25-admissible-calculus}
Let $\mathcal U_\theta$ be as in \eqref{eq:ch25-admissible}.
\begin{enumerate}
\item\label{itm:ch25-adm:nonempty}
 $\mathcal U_\theta$ is nonempty, convex, and stable under finite maxima.
\item\label{itm:ch25-adm:transl}
 For every $a\in\mathbf R$,
 $\{\MA_\theta(u):u\in\mathcal U_\theta\}
   =\{\MA_\theta(u):u\in\PSH(X,\theta),\ a-1\leq u\leq a\}$.
 In particular \eqref{eq:capacity-definition} is unchanged if the
 normalization $-1\leq u\leq0$ is replaced by $0\leq u\leq1$.
\item\label{itm:ch25-adm:scale}
 If $u\in\PSH(X,\theta)$ satisfies $-M\leq u\leq0$ with $M\geq1$, then
 $u/M\in\mathcal U_\theta$.
\item\label{itm:ch25-adm:model}
 $\mathcal U_\theta$ contains the $\theta$-psh model functions with
 values in $[-1,0]$, and every member of $\mathcal U_\theta$ is the limit
 of a decreasing net of $\theta$-psh model functions with values in
 $[-1,c_\alpha]$, where $c_\alpha\downarrow0$.
\end{enumerate}
\end{lemma}

\begin{proof}
\cref{itm:ch25-adm:nonempty}  The constant $0$ is $\theta$-psh because
$\theta$ is semipositive
(\Cref{prop:ch21-constants}\,\cref{itm:ch21-const:const}), and it satisfies
the bounds.  Convexity is
\Cref{thm:ch21-stability}\,\cref{itm:ch21-stab:sum} applied to a convex
combination, together with the fact that the constraint $-1\leq u\leq0$ is
convex; stability under finite maxima is
\Cref{prop:ch21-elementary-stability}, the constraint again being preserved.

\cref{itm:ch25-adm:transl}  Adding a constant preserves $\PSH(X,\theta)$
(\Cref{prop:ch21-constants}\,\cref{itm:ch21-const:transl}) and does not change
$\theta+\ddc u$, because $\ddc$ annihilates constants by
\Cref{def:ch20-ddc}.  Hence $u\mapsto u+a$ is a bijection between the two
families with $\MA_\theta(u+a)=\MA_\theta(u)$.

\cref{itm:ch25-adm:scale}  Write $u/M=(1/M)u+(1-1/M)\cdot0$, a convex
combination of two $\theta$-psh functions with $1/M\in(0,1]$; apply
\cref{itm:ch25-adm:nonempty}.  The bounds $-1\leq u/M\leq0$ are immediate.

\cref{itm:ch25-adm:model}  A $\theta$-psh model function with values in
$[-1,0]$ lies in $\PSH(X,\theta)$ by
\Cref{lem:ch21-model-skeletal-properties}.  For the second assertion apply
\Cref{thm:bfj-decreasing-regularization} to $u$, obtaining a decreasing net
$u_\alpha\searrow u$ of $\theta$-psh model functions; then
$u_\alpha\geq u\geq-1$.  Put $d_\alpha:=\sup_{X^{\An}}u_\alpha$ and
$c_\alpha:=\max\{0,d_\alpha\}$, nonincreasing nets of reals.  Fix an SNC model determining $\theta$ and
write $\Delta$ for its dual complex.  By
\Cref{prop:ch21-sup-on-skeleton} the supremum of any $\theta$-psh function
over $X^{\An}$ equals its supremum over $\Delta$, so
$d_\alpha=\sup_\Delta u_\alpha$ and $\sup u=\sup_\Delta u$.  The
restrictions $u_\alpha|_\Delta$ are continuous and decrease pointwise to
$u|_\Delta$, which is continuous by
\Cref{thm:bfj-skeletal-equicontinuity}; since $\Delta$ is compact,
\Cref{prop:dini-nets} gives uniform convergence on $\Delta$ and hence
$d_\alpha\downarrow\sup u\leq0$ and $c_\alpha\downarrow0$.
\end{proof}

\begin{lemma}
\label{lem:ch25-inner-outer}
\begin{enumerate}
\item\label{itm:ch25-io:open}
 If $G\subseteq X^{\An}$ is open, then
 $\Capa_\theta^*(G)=\Capa_\theta(G)$.
\item\label{itm:ch25-io:mono}
 $\Capa_\theta^*$ is monotone, $\Capa_\theta\leq\Capa_\theta^*$ on Borel
 sets, and $\Capa^*_\theta(E)=0$ implies $\Capa_\theta(E)=0$ for Borel
 $E$.
\item\label{itm:ch25-io:radon}
 For an arbitrary set $E$ and any $u\in\mathcal U_\theta$, the outer
 measure $\MA_\theta(u)^*(E)$ satisfies
 $\MA_\theta(u)^*(E)\leq\Capa^*_\theta(E)$.
\end{enumerate}
\end{lemma}

\begin{proof}
\cref{itm:ch25-io:open}  The infimum defining $\Capa_\theta^*(G)$ is over
open sets containing $G$, of which $G$ itself is one; monotonicity of
$\Capa_\theta$ on Borel sets, proved in
\Cref{prop:capacity-elementary}, gives the reverse inequality.

\cref{itm:ch25-io:mono}  Monotonicity of the infimum is clear.  If
$E\subseteq G$ with $G$ open, then
$\Capa_\theta(E)\leq\Capa_\theta(G)$ by monotonicity on Borel sets, and
taking the infimum over $G$ gives
$\Capa_\theta(E)\leq\Capa_\theta^*(E)$; the last assertion follows.

\cref{itm:ch25-io:radon}  Each $\MA_\theta(u)$ is a Radon measure, hence
outer regular by \eqref{eq:radon-regularity}: $\MA_\theta(u)^*(E)
=\inf\{\MA_\theta(u)(G):G\supseteq E\ \text{open}\}$.  For each such $G$,
$\MA_\theta(u)(G)\leq\Capa_\theta(G)$, and taking the infimum over $G$
gives the claim.
\end{proof}

\subsection{Monotonicity, subadditivity, and the Choquet axioms}
\label{sec:ch25-elementary}

\begin{proposition}
\label{prop:capacity-elementary}
For Borel sets $E,E_j$,
\[
 0\leq\Capa_\theta(E)\leq V,\qquad
 E\subset F\Longrightarrow\Capa_\theta(E)\leq\Capa_\theta(F),
\]
and
\begin{equation}
\label{eq:capacity-subadditivity}
                  \Capa_\theta\Bigl(\bigcup_jE_j\Bigr)
                  \leq\sum_j\Capa_\theta(E_j).
\end{equation}
Every Borel pluripolar set has capacity zero.
\end{proposition}

\begin{proof}
Fix $u\in\mathcal U_\theta$.  By \Cref{thm:bounded-ma-extension} the
measure $\MA_\theta(u)$ is positive with total mass $V$, so
$0\leq\MA_\theta(u)(E)\leq V$ for every Borel $E$, and
$\MA_\theta(u)(E)\leq\MA_\theta(u)(F)$ when $E\subseteq F$.  Taking the
supremum over $u\in\mathcal U_\theta$, which is nonempty by
\Cref{lem:ch25-admissible-calculus}\,\cref{itm:ch25-adm:nonempty}, gives the
first two assertions.

For \eqref{eq:capacity-subadditivity} let $(E_j)_{j\geq1}$ be Borel and let
$u\in\mathcal U_\theta$.  Countable subadditivity of the measure
$\MA_\theta(u)$ gives
\[
 \MA_\theta(u)\Bigl(\bigcup_jE_j\Bigr)
 \leq\sum_j\MA_\theta(u)(E_j)
 \leq\sum_j\Capa_\theta(E_j),
\]
the second step because $u$ is admissible for each individual $E_j$.  The
right-hand side does not depend on $u$; take the supremum on the left.
Observe that the two suprema occur in the correct order here, and that the
same argument fails verbatim for the reverse inequality: there is no reason
for one potential to be nearly optimal for all the $E_j$ simultaneously.
\Cref{ex:ch25-subadditivity-strict} exhibits a strict inequality in
\eqref{eq:capacity-subadditivity} with explicit numbers.

Finally, suppose $E$ is Borel and pluripolar, that is
$E\subseteq\{v=-\infty\}$ for some $v\in\PSH(X,\theta)$ with
$v\not\equiv-\infty$.  For every $u\in\mathcal U_\theta$ the measure
$\MA_\theta(u)=\MA_\theta(u,\ldots,u)$ is a mixed Monge--Amp\`ere measure
of bounded $\theta$-psh potentials, so
\Cref{prop:bounded-ma-no-polar-mass} applies and gives
$\MA_\theta(u)(\{v=-\infty\})=0$; monotonicity of the measure gives
$\MA_\theta(u)(E)=0$.  Taking the supremum over $u$ gives
$\Capa_\theta(E)=0$.  No outer-regularity statement for an arbitrary
non-Borel subset is used here; the sharper assertion that \emph{every}
pluripolar set, Borel or not, has outer capacity zero is
\Cref{cor:ch25-polar-negligible}, and it needs the quantitative estimate of
\Cref{prop:ch25-sublevel-capacity} rather than
\Cref{prop:bounded-ma-no-polar-mass}.
\end{proof}

\begin{corollary}
\label{cor:ch25-outer-subadditivity}
For arbitrary subsets $E_j\subseteq X^{\An}$,
\[
 \Capa_\theta^*\Bigl(\bigcup_{j\geq1}E_j\Bigr)
 \leq\sum_{j\geq1}\Capa_\theta^*(E_j).
\]
\end{corollary}

\begin{proof}
We may assume the right-hand side finite.  Let $\varepsilon>0$ and choose
open $G_j\supseteq E_j$ with
$\Capa_\theta(G_j)\leq\Capa_\theta^*(E_j)+2^{-j}\varepsilon$.  Then
$G=\bigcup_jG_j$ is open and contains $\bigcup_jE_j$, so by
\Cref{lem:ch25-inner-outer}\,\cref{itm:ch25-io:open} and
\eqref{eq:capacity-subadditivity},
\[
 \Capa^*_\theta\Bigl(\bigcup_jE_j\Bigr)\leq\Capa_\theta(G)
 \leq\sum_j\Capa_\theta(G_j)
 \leq\sum_j\Capa^*_\theta(E_j)+\varepsilon .
\]
Let $\varepsilon\downarrow0$.
\end{proof}

The capacity of \Cref{def:bfj-capacity} is therefore an admissible input to
the abstract calculus of \Cref{sec:capacity-preview}, once the two remaining
Choquet axioms have been checked.  The first is a formal interchange of two
suprema; the second is a compactness argument and is available only for the
outer capacity.

\begin{proposition}
\label{prop:ch25-choquet-axioms}
\begin{enumerate}
\item\label{itm:ch25-chq:incr}
 If $E_1\subseteq E_2\subseteq\cdots$ are Borel and
 $E=\bigcup_jE_j$, then
 $\Capa_\theta(E_j)\nearrow\Capa_\theta(E)$.
\item\label{itm:ch25-chq:decr}
 If $K_1\supseteq K_2\supseteq\cdots$ are compact and
 $K=\bigcap_jK_j$, then
 $\Capa^*_\theta(K_j)\searrow\Capa^*_\theta(K)$.
\item\label{itm:ch25-chq:choquet}
 Consequently $\Capa_\theta$ is monotone and countably subadditive on the
 Borel sets and continuous along increasing sequences of Borel sets, and
 $\Capa^*_\theta$ is monotone, countably subadditive on arbitrary sets,
 and continuous along decreasing sequences of compact sets.  These are the
 properties quoted from \Cref{sec:capacity-preview} in the sequel.
\end{enumerate}
\end{proposition}

\begin{proof}
\cref{itm:ch25-chq:incr}  For each fixed $u\in\mathcal U_\theta$ the measure
$\MA_\theta(u)$ is countably additive, so
$\MA_\theta(u)(E_j)\nearrow\MA_\theta(u)(E)$; an increasing limit of real
numbers is their supremum.  Hence
\[
 \Capa_\theta(E)
 =\sup_u\sup_j\MA_\theta(u)(E_j)
 =\sup_j\sup_u\MA_\theta(u)(E_j)
 =\sup_j\Capa_\theta(E_j),
\]
the middle step being the interchange of two suprema, which is
unconditionally valid.  The sequence $\Capa_\theta(E_j)$ is nondecreasing
by monotonicity, so its supremum is its limit.

\cref{itm:ch25-chq:decr}  Monotonicity gives
$\Capa_\theta^*(K_j)\geq\Capa^*_\theta(K)$ for every $j$, and the
sequence is nonincreasing; only the reverse inequality in the limit requires
proof.  Let $G\supseteq K$ be open.  The sets $K_j\setminus G$ are
compact, decreasing, and have empty intersection, since
$\bigcap_j(K_j\setminus G)=K\setminus G=\varnothing$.  In a Hausdorff
space a decreasing sequence of compact sets with empty intersection has an
empty member: otherwise the closed sets $K_j\setminus G$ would form a
family with the finite intersection property inside the compact set
$K_1$, forcing a common point.  Hence $K_j\subseteq G$ for all large
$j$, and then
$\Capa^*_\theta(K_j)\leq\Capa_\theta(G)$ by
\Cref{lem:ch25-inner-outer}.  Taking the limit in $j$ and then the infimum
over open $G\supseteq K$ gives
$\lim_j\Capa^*_\theta(K_j)\leq\Capa^*_\theta(K)$.

\cref{itm:ch25-chq:choquet}  Monotonicity and
$\Capa_\theta(\varnothing)=\Capa^*_\theta(\varnothing)=0$ are
\Cref{prop:capacity-elementary} and
\Cref{lem:ch25-inner-outer}\,\cref{itm:ch25-io:mono}; countable subadditivity
is \eqref{eq:capacity-subadditivity} and
\Cref{cor:ch25-outer-subadditivity}; the two monotone-continuity assertions
are \cref{itm:ch25-chq:incr} and \cref{itm:ch25-chq:decr}.
\end{proof}

\begin{warning}{What is not asserted about the Choquet axioms}
\Cref{prop:ch25-choquet-axioms}\,\cref{itm:ch25-chq:incr} is proved for
increasing sequences of \emph{Borel} sets.  For an increasing sequence of
arbitrary subsets, continuity of the outer capacity
$\Capa^*_\theta(A_j)\nearrow\Capa^*_\theta(\bigcup_jA_j)$ is not proved
here.  The classical route to that statement passes through \emph{strong}
subadditivity,
\[
 \Capa(A\cup B)+\Capa(A\cap B)\leq\Capa(A)+\Capa(B),
\]
which for a supremum of measures is not a formal consequence of the
corresponding identity for each measure separately: the two terms on the left
are computed with two different, and possibly unrelated, competitors.  We
have not verified strong subadditivity for $\Capa_\theta$, and no statement
below uses it.  Consequently Choquet's capacitability theorem is not invoked
anywhere in this book for $\Capa_\theta$, and no analytic set is claimed to
be capacitable.  \Cref{ex:capacity-not-subadditive} is the reminder that such
axioms have to be checked and not assumed.
\end{warning}

\begin{remark}[In dimension one the missing axioms are theorems]
\label{rem:ch25-curve-choquet-unconditional}
For Thuillier's relative capacity, strong subadditivity, capacitability,
and the equivalence of capacity zero with polarity are theorems in
dimension one.  For this \emph{relative} capacity of a smooth
Berkovich curve, defined not as a supremum of Monge--Amp\`ere masses but as
the reciprocal of a minimal energy, Thuillier proves all three
unconditionally \cite[\S3.6]{Thu05}; the statements are collected in
\Cref{thm:ch25-thuillier-capacity} and
\Cref{thm:ch25-thuillier-polar} below, and
\Cref{rem:ch25-thuillier-vs-bfj} explains exactly why they do not transfer
to $\Capa_\theta$ as it stands.  The reader should read the present
warning as concerning the global capacity $\Capa_\theta$; the relative
results do not settle its omitted properties.
\end{remark}

\subsection{Positivity: the capacity sees every open set}
\label{sec:ch25-positivity}

A capacity that vanished on a large set would be useless: the whole point of
\eqref{eq:capacity-definition} is that ``negligible'' should be a strong
restriction.  The following lemma is what prevents the theory from being
vacuous: a capacity-zero set contains no divisorial point, and an open
capacity-zero set is empty.  These conclusions are used in domination
(\Cref{prop:bounded-capacity-domination}) and to
extract a sequential regularization
(\Cref{cor:psh-sequential-regularization}).

\begin{lemma}
\label{lem:divisorial-positive-capacity}
If $x$ is divisorial, then $\Capa_\theta(\{x\})>0$. Consequently every
nonempty open subset of $X^{\An}$ has positive capacity.
\end{lemma}

\begin{proof}
Let $x$ be divisorial.  By \Cref{rem:quasimonomial-vs-divisorial} there is
a model $\mathcal X_0$ with a component $E_0$ of its special fibre such
that $x=x_{E_0}$, and the same is true, with the strict transform of
$E_0$ and the same multiplicity, on every model dominating
$\mathcal X_0$.  Choose a model dominating both $\mathcal X_0$ and a
determination of $\theta$.  By
\Cref{prop:ch20-positive-representative}, with $\varepsilon=1$ in
\eqref{eq:ch20-small-positive-representative}, a further dominating
model $\mathcal X$ carries a model function $\psi$ with
\[
                       -1\leq\psi\leq0
\]
such that $\theta+\ddc\psi$ is $\mathcal X$-positive, that is, is
represented on $\mathcal X$ by a relatively ample class.  In particular
$\psi\in\mathcal U_\theta$: it is $\theta$-psh because relatively ample
implies relatively nef
(\Cref{lem:ch20-special-fibre-test}\,\cref{itm:ch20-sft:cone}), and it obeys
the normalization.

Let $E\subseteq\mathcal X_s$ be the component with $x=x_E$ and let
$b_E$ be its multiplicity in $\mathcal X_s$.  By
\Cref{lem:ch20-special-fibre-test}\,\cref{itm:ch20-sft:ample} the restriction
$(\theta+\ddc\psi)|_E$ is an ample class on the projective $k$-scheme
$E$, which is of pure dimension $n$ because
$\mathcal X\to\Spec R$ is flat and projective of relative dimension $n$.
The Nakai--Moishezon criterion therefore gives
$\bigl((\theta+\ddc\psi)|_E\bigr)^n>0$, and
\eqref{eq:model-mixed-measure} evaluates the model measure at the atom
$x_E$:
\begin{equation}
\label{eq:ch25-divisorial-atom}
 \Capa_\theta(\{x\})\geq\MA_\theta(\psi)(\{x_E\})
 =b_E\bigl((\theta+\ddc\psi)|_E\bigr)^n>0 .
\end{equation}

For the second assertion let $G\neq\varnothing$ be open.  Divisorial points
are dense in $X^{\An}$
(\Cref{rem:quasimonomial-vs-divisorial}, \Cref{thm:skeleton-inverse-limit}),
so $G$ contains one, say $x$, and monotonicity gives
$\Capa_\theta(G)\geq\Capa_\theta(\{x\})>0$.
\end{proof}

\begin{corollary}
\label{cor:ch25-negligible-empty-interior}
A negligible set has empty interior, and its complement is dense.  In
particular $X^{\An}$ itself is not negligible, and
$\Capa_\theta(X^{\An})=V$.
\end{corollary}

\begin{proof}
If $E$ is negligible and $G\subseteq E$ is open and nonempty, then
$\Capa_\theta(G)\leq\Capa_\theta^*(E)=0$ by monotonicity, contradicting
\Cref{lem:divisorial-positive-capacity}.  A set with empty interior has dense
complement.  Finally $\Capa_\theta(X^{\An})=V$ because every admissible
$\MA_\theta(u)$ has total mass $V$ by
\eqref{eq:mixed-total-mass}, and the constant $0$ is admissible.
\end{proof}

\begin{remark}[The lower bound \eqref{eq:ch25-divisorial-atom} is explicit]
\label{rem:ch25-divisorial-explicit}
Nothing in the proof is asymptotic: given a model on which $\theta$ becomes
relatively ample after a potential bounded by $1$, the number
$b_E((\theta+\ddc\psi)|_E)^n$ is an intersection number on the special
fibre and can be computed.  \Cref{ex:ch25-segment-capacities} carries this
out on a semistable curve, where the answer is a rational number written in
terms of the degrees of $\theta$ on the components and the lengths of the
edges of the dual graph.  There is no bound depending only on $n$ and $V$ across all polarized
curves: on a segment of length $\ell$ with $m_0=m_1=1$, the midpoint has
capacity $4/\ell$ once $\ell\geq2$
(\Cref{prop:ch25-segment-capacities}), tending to $0$ as the geometry
varies.  This example does not settle a uniform bound on a fixed curve.  Thus \Cref{lem:divisorial-positive-capacity} is a
qualitative statement, and every quantitative use of it below is made at one
fixed point.
\end{remark}

\begin{warning}{Capacity zero versus polar}
The implication proved so far is
\[
       \text{Borel pluripolar}\Longrightarrow\text{capacity zero}
\]
(\Cref{prop:capacity-elementary}), strengthened in
\Cref{cor:ch25-polar-negligible} to arbitrary pluripolar sets and to the
outer capacity.  Its converse is not a formal consequence of the definition,
is not proved in this book, and is not needed for the BFJ variational
argument.  In particular the words \emph{negligible} and \emph{pluripolar}
are not used as synonyms here; \Cref{sec:ch25-negligible} fixes the
terminology, and the warning after
\Cref{cor:ch25-polar-countable} there says exactly what would be
required to close the gap.
\end{warning}

\section{Two estimates}
\label{sec:ch25-estimates}

Two inequalities carry the entire chapter.  The first is a
Chern--Levine--Nirenberg bound: it says that a uniformly bounded psh family
is uniformly absolutely continuous with respect to capacity, so that a set of
small capacity is small for \emph{every} such potential at once.  The second
is a Cauchy--Schwarz bootstrap which converts a smallness statement for one
distinguished measure into a smallness statement uniform over the admissible
family, and therefore into a statement about the capacity.  The first is
half a line; the second is the technical heart of the theory.

\subsection{The Chern--Levine--Nirenberg bound}
\label{sec:ch25-cln}

\begin{lemma}
\label{lem:ma-capacity-bound}
If $u\in\PSH(X,\theta)$ and $-M\leq u\leq0$ with $M\geq1$, then for
every Borel set $E$,
\begin{equation}
\label{eq:ma-capacity-bound}
                       \MA_\theta(u)(E)
                       \leq M^n\Capa_\theta(E).
\end{equation}
\end{lemma}

\begin{proof}
Since $M\geq1$ and $\theta$ is semipositive, the class
$(M-1)\theta$ is semipositive, and expanding the $n$-fold mixed product
by multilinearity (\Cref{cor:bounded-ma-multilinearity}) gives
\[
 \begin{aligned}
 (M\theta+\ddc u)^n
 &=\bigl((\theta+\ddc u)+(M-1)\theta\bigr)^n\\
 &=(\theta+\ddc u)^n+\sum_{k=1}^n\binom nk(M-1)^k
   \theta^{\wedge k}\wedge(\theta+\ddc u)^{\wedge(n-k)} .
 \end{aligned}
\]
Every summand on the right is a positive measure, by
\Cref{prop:model-ma-basic} at model level and by
\Cref{thm:bounded-ma-extension} in the limit, so
\begin{equation}
\label{eq:ch25-cln-expansion}
 (\theta+\ddc u)^n\leq(M\theta+\ddc u)^n
 =M^n\Bigl(\theta+\ddc\tfrac uM\Bigr)^{\! n}
\end{equation}
as measures.  By
\Cref{lem:ch25-admissible-calculus}\,\cref{itm:ch25-adm:scale} the function
$u/M$ lies in $\mathcal U_\theta$, so integrating
\eqref{eq:ch25-cln-expansion} over $E$ and using
\eqref{eq:capacity-definition} gives
$\MA_\theta(u)(E)\leq M^n\MA_\theta(u/M)(E)\leq M^n\Capa_\theta(E)$.
\end{proof}

\begin{corollary}
\label{cor:ch25-outer-ma-bound}
Let $u\in\PSH(X,\theta)$ with $-M\leq u\leq0$ and $M\geq1$.
\begin{enumerate}
\item\label{itm:ch25-om:outer}
 For every subset $E\subseteq X^{\An}$, the outer measure satisfies
 $\MA_\theta(u)^*(E)\leq M^n\Capa_\theta^*(E)$.
\item\label{itm:ch25-om:null}
 Every negligible set is $\MA_\theta(u)$-null, in the sense that it is
 contained in a Borel set of $\MA_\theta(u)$-measure zero.  The same holds
 for every mixed measure $\MA_\theta(u_1,\ldots,u_n)$ with all
 $-M\leq u_i\leq0$.
\item\label{itm:ch25-om:mixed}
 That mixed statement holds in the quantitative form: for all
 $u_1,\ldots,u_n\in\PSH(X,\theta)$ with $-M\leq u_i\leq0$ and every
 subset $E\subseteq X^{\An}$,
 \[
  \MA_\theta(u_1,\ldots,u_n)^*(E)
  \leq\frac{M^nn^n}{n!}\,\Capa^*_\theta(E).
 \]
\end{enumerate}
\end{corollary}

\begin{proof}
\cref{itm:ch25-om:outer}  By outer regularity of the Radon measure
$\MA_\theta(u)$, \eqref{eq:radon-regularity},
$\MA_\theta(u)^*(E)=\inf\{\MA_\theta(u)(G):G\supseteq E\ \text{open}\}$,
and for each such $G$ we have
$\MA_\theta(u)(G)\leq M^n\Capa_\theta(G)$ by
\eqref{eq:ma-capacity-bound}.  Take the infimum.

\cref{itm:ch25-om:null}  The first assertion is immediate from
\cref{itm:ch25-om:outer}.  For the mixed statement, set
$v_i:=u_i/M\in\mathcal U_\theta$
(\Cref{lem:ch25-admissible-calculus}\,\cref{itm:ch25-adm:scale}) and
$v:=(v_1+\cdots+v_n)/n$, which again lies in $\mathcal U_\theta$ by
convexity.  Multilinearity (\Cref{cor:bounded-ma-multilinearity}) expands
\[
 \MA_\theta(v)=\Bigl(\tfrac1n\sum_{i=1}^n(\theta+\ddc v_i)\Bigr)^{\!n}
 =\frac1{n^n}\sum_{i_1,\ldots,i_n=1}^n
    \MA_\theta(v_{i_1},\ldots,v_{i_n}),
\]
and every summand is a positive measure.  Discarding all summands except the
$n!$ permutations of $(v_1,\ldots,v_n)$, which are equal by symmetry,
gives $\MA_\theta(v)\geq n^{-n}n!\,\MA_\theta(v_1,\ldots,v_n)$.  On the
other hand $\theta+\ddc u_i\leq M\theta+\ddc u_i=M(\theta+\ddc v_i)$, and
replacing each factor of a mixed product by a larger semipositive one
increases the product, because the difference expands as in
\eqref{eq:ch25-cln-expansion} into positive mixed terms; hence
$\MA_\theta(u_1,\ldots,u_n)\leq M^n\MA_\theta(v_1,\ldots,v_n)$ and we
obtain
\[
 \MA_\theta(u_1,\ldots,u_n)
 \leq\frac{M^nn^n}{n!}\,\MA_\theta(v),
 \qquad v\in\mathcal U_\theta .
\]
Since $v\in\mathcal U_\theta$ gives $\MA_\theta(v)(G)\leq\Capa_\theta(G)$
for every open $G$, taking the infimum over open $G\supseteq E$ as in
\cref{itm:ch25-om:outer} turns this into
\cref{itm:ch25-om:mixed}:
$\MA_\theta(u_1,\ldots,u_n)^*(E)\leq(M^nn^n/n!)\Capa^*_\theta(E)$.  In
particular a set of capacity zero is null for the mixed measure as well,
with the constant $M^nn^n/n!$ in place of $M^n$.
\end{proof}

\Cref{cor:ch25-outer-ma-bound} is exactly the hypothesis of
\Cref{prop:polar-sigma-ideal}\,(3), with $C=M^n$; this is the interface
announced in \Cref{rem:capacity-interface}.  In particular the negligible
sets form a $\sigma$-ideal, and every Borel negligible set is null for
every Monge--Amp\`ere measure of bounded potentials.  We restate this in
\Cref{prop:ch25-negligible-sigma-ideal} once the terminology of
\Cref{sec:ch25-negligible} is in place.

\subsection{A uniform integrability estimate}
\label{sec:ch25-uniform-integrability}

The Chern--Levine--Nirenberg bound compares one Monge--Amp\`ere measure with
the capacity.  The next estimate compares the integral of an
\emph{unbounded} psh function against a bounded Monge--Amp\`ere measure with
its integral against the fixed model measure $\theta^n$, and does so with a
constant independent of everything but $n$, $M$ and $V$.  It is the
only place in the chapter where an unbounded potential is integrated, and it
is what makes polar sets negligible with a rate.

\begin{lemma}
\label{lem:ch25-uniform-integrability}
Let $u\in\PSH(X,\theta)$ with $\sup_{X^{\An}}u=0$, and let
$\varphi_1,\ldots,\varphi_n\in\PSH(X,\theta)$ satisfy
$-M\leq\varphi_i\leq0$ with $M\geq0$.  Then
\begin{equation}
\label{eq:ch25-uniform-integrability}
 0\leq\int_{X^{\An}}(-u)\,\MA_\theta(\varphi_1,\ldots,\varphi_n)
 \leq\int_{X^{\An}}(-u)\,\theta^n+nMV .
\end{equation}
\end{lemma}

\begin{proof}
Both sides are monotone under replacing $u$ by the truncation
$u^{(m)}:=\max\{u,-m\}$, which is $\theta$-psh by
\Cref{thm:ch21-stability}\,\cref{itm:ch21-stab:max} and again satisfies
$\sup u^{(m)}=0$; and $-u^{(m)}\nearrow-u$ pointwise.  By monotone
convergence for the two fixed positive measures involved
(\Cref{lem:integration-ledger}) it suffices to prove
\eqref{eq:ch25-uniform-integrability} for bounded $u$, which we now assume.

Write $T_k:=(\theta+\ddc\varphi_{k+1})\wedge\cdots\wedge
(\theta+\ddc\varphi_n)$ as a formal placeholder for $0\leq k\leq n$;
only the completed $n$-fold products below denote measures.  Thus $T_n$ is the empty
product and $\theta^{\wedge0}\wedge T_0=\MA_\theta(\varphi_1,\ldots,
\varphi_n)$.  We prove by induction on $k$ that
\begin{equation}
\label{eq:ch25-ui-induction}
 \int(-u)\,\theta^{\wedge k}\wedge T_k
 \leq\int(-u)\,\theta^{\wedge(k+1)}\wedge T_{k+1}+MV ,
 \qquad 0\leq k\leq n-1 ,
\end{equation}
where $\theta^{\wedge k}\wedge T_k$ denotes the mixed measure with $k$
slots filled by $\theta$, that is by the potential $0$.  Granting
\eqref{eq:ch25-ui-induction}, summing over $k=0,\ldots,n-1$ telescopes to
\eqref{eq:ch25-uniform-integrability}, the case $k=0$ being the left-hand
side and the case $k=n-1$ producing $\theta^{\wedge n}=\theta^n$.

To prove \eqref{eq:ch25-ui-induction}, split the $(k+1)$-st slot:
\[
 \theta^{\wedge k}\wedge T_k
 =\theta^{\wedge k}\wedge(\theta+\ddc\varphi_{k+1})\wedge T_{k+1}
 =\theta^{\wedge(k+1)}\wedge T_{k+1}
  +\theta^{\wedge k}\wedge\ddc\varphi_{k+1}\wedge T_{k+1},
\]
an identity of signed measures in which each of the first two terms is a
positive measure of mass $V$ by \eqref{eq:mixed-total-mass}.  Integration
by parts, valid for differences of bounded $\theta$-psh functions by
\Cref{cor:bounded-ma-multilinearity} and \eqref{eq:model-ibp}, gives
\[
 \int(-u)\,\theta^{\wedge k}\wedge\ddc\varphi_{k+1}\wedge T_{k+1}
 =\int(-\varphi_{k+1})\,\theta^{\wedge k}\wedge\ddc u\wedge T_{k+1}.
\]
Now write $\ddc u=(\theta+\ddc u)-\theta$ and use
$0\leq-\varphi_{k+1}\leq M$:
\[
\begin{split}
 \int(-\varphi_{k+1})\,\theta^{\wedge k}\wedge\ddc u\wedge T_{k+1}
 ={}&\int(-\varphi_{k+1})\,
     \theta^{\wedge k}\wedge(\theta+\ddc u)\wedge T_{k+1}\\
    &-\int(-\varphi_{k+1})\,\theta^{\wedge(k+1)}\wedge T_{k+1}\\
 \leq{}&\int(-\varphi_{k+1})\,
     \theta^{\wedge k}\wedge(\theta+\ddc u)\wedge T_{k+1}
 \ \leq\ MV,
\end{split}
\]
because the subtracted term is an integral of a nonnegative function against
a positive measure, and the remaining measure is positive of mass $V$.
This is \eqref{eq:ch25-ui-induction}.  The left-hand inequality in
\eqref{eq:ch25-uniform-integrability} holds because $u\leq\sup u=0$.
\end{proof}

\begin{lemma}
\label{lem:ch25-normalized-bound}
There is a finite constant $C_\theta$, depending only on $X$, $\theta$
and a choice of SNC model determining $\theta$, such that
\begin{equation}
\label{eq:ch25-normalized-bound}
 \int_{X^{\An}}(-u)\,\theta^n\leq C_\theta
 \qquad\text{for every }u\in\PSH(X,\theta)\text{ with }\sup u=0 .
\end{equation}
One may take $C_\theta=L\cdot\operatorname{diam}(\Delta_{\mathcal X})\cdot
V$, where $\mathcal X$ is an SNC model determining $\theta$,
$\Delta_{\mathcal X}$ its dual complex with the path metric of
\Cref{lem:ch22-path-metric}, and $L=L(\mathcal X,\theta)$ the equicontinuity
constant of \Cref{thm:bfj-skeletal-equicontinuity}.
\end{lemma}

\begin{proof}
Let $\mathcal X$ be an SNC model determining $\theta$, and write
$\Delta=\Delta_{\mathcal X}$.  By
\eqref{eq:model-mixed-measure} the measure $\theta^n$ is a finite sum
$\sum_Eb_E(\theta_{\mathcal X}|_E)^n\delta_{x_E}$ of point masses at
divisorial points $x_E\in\Delta$, with nonnegative coefficients summing to
$V$ by \eqref{eq:mixed-total-mass}.  Therefore
\[
 \int(-u)\,\theta^n=\sum_Eb_E(\theta_{\mathcal X}|_E)^n\,(-u(x_E))
 \leq V\cdot\max_{x\in\Delta}(-u(x)) .
\]
It remains to bound $-u$ on $\Delta$.  By
\Cref{prop:ch21-sup-on-skeleton} the supremum of $u$ over $X^{\An}$ is
attained on $\Delta$, so there is $y_0\in\Delta$ with $u(y_0)=0$.  By
\Cref{thm:bfj-skeletal-equicontinuity} and
\Cref{cor:ch22-explicit-modulus}, the restriction $u|_\Delta$ is
$L$-Lipschitz for the path metric $d_\Delta$ of
\Cref{lem:ch22-path-metric}, with $L$ independent of $u$.  Hence for
every $x\in\Delta$,
\[
 -u(x)=u(y_0)-u(x)\leq L\,d_\Delta(x,y_0)
 \leq L\operatorname{diam}(\Delta),
\]
the diameter being finite because $\Delta$ is a compact connected
polyhedral complex.  Combining the two displays gives
\eqref{eq:ch25-normalized-bound} with the stated constant.
\end{proof}

\begin{proposition}
\label{prop:ch25-sublevel-capacity}
Let $u\in\PSH(X,\theta)$ with $\sup_{X^{\An}}u=0$.  Then for every
$t>0$,
\begin{equation}
\label{eq:ch25-sublevel-capacity}
 \Capa_\theta\bigl(\{u<-t\}\bigr)
 \leq\frac{C_\theta+nV}{t},
\end{equation}
with $C_\theta$ as in \eqref{eq:ch25-normalized-bound}.  The set
$\{u<-t\}$ is open, so the same bound holds for the outer capacity.
\end{proposition}

\begin{proof}
The set $\{u<-t\}$ is open because $u$ is upper semicontinuous, and for
an open set inner and outer capacity agree
(\Cref{lem:ch25-inner-outer}\,\cref{itm:ch25-io:open}).  Let
$v\in\mathcal U_\theta$.  On $\{u<-t\}$ we have $-u>t$, so Chebyshev's
inequality for the positive measure $\MA_\theta(v)$ gives
\[
 \MA_\theta(v)\bigl(\{u<-t\}\bigr)
 \leq\frac1t\int(-u)\,\MA_\theta(v)
 \leq\frac1t\Bigl(\int(-u)\theta^n+nV\Bigr)
 \leq\frac{C_\theta+nV}{t},
\]
the middle step being \eqref{eq:ch25-uniform-integrability} with
$\varphi_1=\cdots=\varphi_n=v$ and $M=1$, and the last
\eqref{eq:ch25-normalized-bound}.  The bound is independent of $v$; take
the supremum.
\end{proof}

\begin{corollary}
\label{cor:ch25-polar-negligible}
Let $v\in\PSH(X,\theta)$, $v\not\equiv-\infty$.  Then
$\Capa^*_\theta(\{v=-\infty\})=0$.  Consequently every subset of the polar
set of a $\theta$-psh function is negligible, with no measurability
hypothesis, and a countable union of such sets is negligible.
\end{corollary}

\begin{proof}
Normalize $\sup v=0$, which is legitimate because $v$ is bounded above
(\Cref{prop:ch21-sup-on-skeleton}) and translation preserves
$\PSH(X,\theta)$.  The polar set is contained in the open set
$\{v<-t\}$ for every $t>0$, so
\[
 \Capa^*_\theta(\{v=-\infty\})
 \leq\Capa_\theta(\{v<-t\})\leq\frac{C_\theta+nV}{t}
 \xrightarrow[\ t\to\infty\ ]{}0 .
\]
Monotonicity of $\Capa^*_\theta$ gives the statement for subsets, and
\Cref{cor:ch25-outer-subadditivity} the statement for countable unions.
\end{proof}

\begin{remark}[Two proofs, two strengths]
\label{rem:ch25-two-polar-proofs}
The Borel case was already proved in \Cref{prop:capacity-elementary}:
it gives $\Capa_\theta(E)=0$ for Borel pluripolar $E$, using
\Cref{prop:bounded-ma-no-polar-mass}.  The present corollary is stronger in
three ways: it applies to arbitrary, possibly non-measurable, subsets; it
controls the \emph{outer} capacity; and it comes with the explicit rate
\eqref{eq:ch25-sublevel-capacity}, which is what
\Cref{chap:bfj-energy} sharpens in \Cref{cor:e1-capacity-tail}.  It also
gives the tail input for the extended-real-valued quasicontinuity theorem
\Cref{thm:ch25-psh-quasicontinuity}.  The rate is in fact optimal:
\Cref{ex:ch25-pluripolar-computed} computes both sides on $\mathbf P^1$
with the standard model form and finds that
\eqref{eq:ch25-sublevel-capacity} holds there with equality, exponent and
constant included, because the dual complex of $\mathbf P^1_R$ is a single
point and $C_\theta=0$.
\end{remark}

\subsection{The fundamental capacity estimate}
\label{sec:ch25-fundamental}

We now come to the inequality that makes quasicontinuity work.  The
mechanism is the following.  Suppose one wants to make
$\int(\psi-\varphi)\MA_\theta(u)$ small \emph{uniformly} in
$u\in\mathcal U_\theta$; this is what a capacity statement requires, since
the capacity is a supremum over $u$.  Direct control is hopeless, because
the measures $\MA_\theta(u)$ vary wildly with $u$.  The estimate below
replaces this uniform requirement by a single condition, namely smallness of
$\int(\psi-\varphi)\MA_\theta(\varphi/2)$, in which the measure is
determined by the data.  For a decreasing net that single condition is the
fixed-entry case of the weighted decreasing-net continuity in
\Cref{thm:bounded-ma-extension}; and that is the whole proof of
\Cref{thm:bounded-quasicontinuity}.

The mechanism by which one factor is exchanged for another is a
Cauchy--Schwarz inequality for the Dirichlet form associated with a mixed
product.  We isolate it first.

\begin{lemma}
\label{lem:ch25-dirichlet-form}
Let $w_2,\ldots,w_n$ be bounded $\theta$-psh functions and put
$T:=(\theta+\ddc w_2)\wedge\cdots\wedge(\theta+\ddc w_n)$ as a formal
placeholder in the completed bounded pairings of
\Cref{cor:bounded-ma-multilinearity}.  For bounded
$\theta$-psh $a,b$ and their differences define
\[
 \langle a,b\rangle_T:=-\int a\,\ddc b\wedge T .
\]
Then $\langle\cdot,\cdot\rangle_T$ is a symmetric, positive semidefinite
bilinear form on the real vector space spanned by the bounded
$\theta$-psh functions, and
\begin{equation}
\label{eq:ch25-cauchy-schwarz}
 \bigl|\langle a,b\rangle_T\bigr|
 \leq\langle a,a\rangle_T^{1/2}\,\langle b,b\rangle_T^{1/2}.
\end{equation}
\end{lemma}

\begin{proof}
Bilinearity is \Cref{cor:bounded-ma-multilinearity}, which asserts that mixed
products are affine in each psh entry; the expression
$\ddc b\wedge T=(\theta+\ddc b)\wedge T-\theta\wedge T$ is a difference of
two positive measures, and both are affine in $b$, so
$b\mapsto\ddc b\wedge T$ is additive on differences.  Symmetry is
integration by parts \eqref{eq:model-ibp}, extended to bounded psh
differences by \Cref{cor:bounded-ma-multilinearity}.  Positive
semidefiniteness is \eqref{eq:model-dirichlet-negativity} in the same
extended form: $\int a\,\ddc a\wedge T\leq0$, that is
$\langle a,a\rangle_T\geq0$.  A symmetric positive semidefinite bilinear
form satisfies the Cauchy--Schwarz inequality
\eqref{eq:ch25-cauchy-schwarz} by the usual discriminant argument applied to
$\lambda\mapsto\langle a+\lambda b,a+\lambda b\rangle_T\geq0$.
\end{proof}

\begin{lemma}
\label{lem:fundamental-capacity-estimate}
Suppose
\[
 -M\leq\varphi\leq\psi\leq0,\qquad
 -M\leq u_i\leq0
\]
are bounded $\theta$-psh functions and $M\geq1$. Then
\begin{equation}
\label{eq:fundamental-capacity-estimate}
\begin{split}
 0\leq{}&\int(\psi-\varphi)
       \MA_\theta(u_1,\ldots,u_n)\\
 \leq{}&4M V^{1-2^{-n}}
 \left(\int(\psi-\varphi)
       \MA_\theta(\varphi/2)\right)^{2^{-n}}.
\end{split}
\end{equation}
\end{lemma}

\begin{proof}
All the integrands are bounded Borel functions and all the measures are
positive of mass $V$, so every integral written below is a finite real
number; the manipulations are the multilinear ones of
\Cref{cor:bounded-ma-multilinearity}, valid for bounded $\theta$-psh
entries, and no regularization is needed once those are granted.  Write
\[
 f:=\psi-\varphi,\qquad 0\leq f\leq M ,
\]
the bounds because $-M\leq\varphi\leq\psi\leq0$.  The left-hand inequality
in \eqref{eq:fundamental-capacity-estimate} is the positivity of the mixed
measure.

\emph{Step 1: the exchange inequality.}  We claim that for every
$(n-1)$-fold mixed product
$T=(\theta+\ddc w_2)\wedge\cdots\wedge(\theta+\ddc w_n)$ of bounded
$\theta$-psh functions, interpreted as a formal placeholder in completed
$n$-fold products, and every bounded $\theta$-psh $w$ with $-M\leq w\leq0$,
\begin{equation}
\label{eq:ch25-exchange}
 \int f\,(\theta+\ddc w)\wedge T
 \ \leq\ 2(MV)^{1/2}
   \Bigl(\int f\,\bigl(\theta+\ddc\tfrac\varphi2\bigr)\wedge T\Bigr)^{1/2}.
\end{equation}
To see this, split the free factor:
\[
 \int f\,(\theta+\ddc w)\wedge T
 =\underbrace{\int f\,\theta\wedge T}_{=:A}
  +\int f\,\ddc w\wedge T .
\]
The measure $\theta\wedge T$ is positive of mass $V$ by
\eqref{eq:mixed-total-mass}, so $0\leq A\leq MV$ and therefore
\begin{equation}
\label{eq:ch25-A-term}
 A=A^{1/2}A^{1/2}\leq(MV)^{1/2}A^{1/2}.
\end{equation}
For the second term use the Dirichlet form of
\Cref{lem:ch25-dirichlet-form}:
\[
 \int f\,\ddc w\wedge T=-\langle f,w\rangle_T
 \leq\langle f,f\rangle_T^{1/2}\langle w,w\rangle_T^{1/2}.
\]
Now
\[
 \begin{aligned}
 \langle w,w\rangle_T
 &=-\int w\,\ddc w\wedge T\\
 &=\int(-w)(\theta+\ddc w)\wedge T-\int(-w)\theta\wedge T\\
 &\leq\int(-w)(\theta+\ddc w)\wedge T\leq MV,
 \end{aligned}
\]
because $0\leq-w\leq M$, the subtracted term is nonnegative, and
$(\theta+\ddc w)\wedge T$ is positive of mass $V$.  Likewise, since
$f=\psi-\varphi$,
\[
 \begin{aligned}
 \langle f,f\rangle_T
 &=\int f\,\ddc(\varphi-\psi)\wedge T\\
 &=\int f\,(\theta+\ddc\varphi)\wedge T
   -\int f\,(\theta+\ddc\psi)\wedge T\\
 &\leq\underbrace{\int f\,(\theta+\ddc\varphi)\wedge T}_{=:C},
 \end{aligned}
\]
because $f\geq0$ and $(\theta+\ddc\psi)\wedge T$ is a positive measure.
Combining,
\[
 \int f\,(\theta+\ddc w)\wedge T\leq(MV)^{1/2}\bigl(A^{1/2}+C^{1/2}\bigr).
\]
Finally $A^{1/2}+C^{1/2}\leq\sqrt2\,(A+C)^{1/2}$ by the
Cauchy--Schwarz inequality in $\mathbf R^2$, and
\[
 A+C=\int f\,\bigl(2\theta+\ddc\varphi\bigr)\wedge T
 =2\int f\,\bigl(\theta+\ddc\tfrac\varphi2\bigr)\wedge T ,
\]
the last equality because $\ddc$ is linear
(\Cref{def:ch20-ddc}).  Hence
$A^{1/2}+C^{1/2}\leq2\bigl(\int f(\theta+\ddc\frac\varphi2)\wedge
T\bigr)^{1/2}$, which is \eqref{eq:ch25-exchange}.  Note that
$\varphi/2=\frac12\varphi+\frac12\cdot0$ is $\theta$-psh by
\Cref{lem:ch25-admissible-calculus}\,\cref{itm:ch25-adm:nonempty}, so the
right-hand side of \eqref{eq:ch25-exchange} is again an integral against a
positive measure of mass $V$.

\emph{Step 2: iteration.}  Put, for $0\leq k\leq n$,
\[
 J_k:=\int f\,
  \bigl(\theta+\ddc\tfrac\varphi2\bigr)^{\wedge k}
  \wedge(\theta+\ddc u_{k+1})\wedge\cdots\wedge(\theta+\ddc u_n),
\]
so that $J_0=\int f\,\MA_\theta(u_1,\ldots,u_n)$ is the left-hand side of
\eqref{eq:fundamental-capacity-estimate} and
$J_n=\int f\,\MA_\theta(\varphi/2)$ is the quantity on the right.  Applying
\eqref{eq:ch25-exchange} with $w=u_{k+1}$ and
$T=(\theta+\ddc\frac\varphi2)^{\wedge k}\wedge(\theta+\ddc u_{k+2})
\wedge\cdots\wedge(\theta+\ddc u_n)$ gives
\begin{equation}
\label{eq:ch25-iteration-step}
 J_k\leq2(MV)^{1/2}J_{k+1}^{1/2},\qquad 0\leq k\leq n-1 .
\end{equation}
The hypotheses of Step~1 are met at every stage: the entries
$u_{k+2},\ldots,u_n$ and $\varphi/2$ are bounded $\theta$-psh, and the
free entry $u_{k+1}$ satisfies $-M\leq u_{k+1}\leq0$.

Iterating \eqref{eq:ch25-iteration-step} from $k=0$ to $k=n-1$,
\[
 \begin{aligned}
 J_0&\leq\prod_{k=0}^{n-1}\bigl(2(MV)^{1/2}\bigr)^{2^{-k}}\cdot
 J_n^{2^{-n}}\\
 &=\bigl(2(MV)^{1/2}\bigr)^{2(1-2^{-n})}J_n^{2^{-n}}\\
 &=2^{2(1-2^{-n})}(MV)^{1-2^{-n}}J_n^{2^{-n}},
 \end{aligned}
\]
using $\sum_{k=0}^{n-1}2^{-k}=2(1-2^{-n})$.  Since $n\geq1$ we have
$2^{2(1-2^{-n})}\leq4$, and since $M\geq1$ we have
$M^{1-2^{-n}}\leq M$.  Therefore
$J_0\leq4MV^{1-2^{-n}}J_n^{2^{-n}}$, which is
\eqref{eq:fundamental-capacity-estimate}.
\end{proof}

\begin{remark}[Reading the estimate]
\label{rem:ch25-reading-fundamental}
Three features of \eqref{eq:fundamental-capacity-estimate} are used later and
are worth isolating.

\emph{The left-hand side is uniform in $u_1,\ldots,u_n$.}  The right-hand
side does not mention them at all.  This is precisely what a capacity
statement needs, since $\Capa_\theta$ is a supremum over the admissible
family; taking the supremum over $u_1=\cdots=u_n=u\in\mathcal U_\theta$
with $M=1$ yields
\begin{equation}
\label{eq:ch25-uniform-form}
 \sup_{u\in\mathcal U_\theta}\int(\psi-\varphi)\MA_\theta(u)
 \leq4V^{1-2^{-n}}
 \Bigl(\int(\psi-\varphi)\MA_\theta(\varphi/2)\Bigr)^{2^{-n}}
\end{equation}
whenever $-1\leq\varphi\leq\psi\leq0$.

\emph{The measure on the right is fixed by $\varphi$ alone.}  In every
application $\varphi$ is the limit of a decreasing net $\psi_\alpha$, and
one applies the estimate with $\psi=\psi_\alpha$.  The required limit
\[
 \int(\psi_\alpha-\varphi)\MA_\theta(\varphi/2)\longrightarrow0
\]
is the fixed-entry case of the weighted decreasing-net continuity in
\Cref{thm:bounded-ma-extension}: apply it to the weight
$\psi_\alpha\searrow\varphi$ while keeping all Monge--Amp\`ere entries equal
to $\varphi/2$.  This point cannot be replaced by ordinary monotone
convergence for nets, since $\psi_\alpha-\varphi$ is lower, not upper,
semicontinuous in general.

\emph{The exponent degrades with the dimension.}  The power $2^{-n}$ is
what the iteration produces and we make no claim of sharpness; for
$n=1$ it is $1/2$ and \eqref{eq:fundamental-capacity-estimate} reads
$\int(\psi-\varphi)\MA_\theta(u)\leq4MV^{1/2}
(\int(\psi-\varphi)\MA_\theta(\varphi/2))^{1/2}$.
\Cref{ex:ch25-fundamental-numbers} tests this case against an explicit
computation on a metric graph and finds a loss of a factor of about
$4V^{1/2}$, so the shape of the estimate is right and the constant is not
optimal.
\end{remark}

\begin{warning}{The estimate is a bounded-regime statement}
Every entry in \eqref{eq:fundamental-capacity-estimate} is assumed bounded,
and the bound $M$ appears in the conclusion.  This lemma does not supply
an estimate for unbounded potentials in which $M$ is replaced by an energy
quantity: the reader should not read \eqref{eq:ch25-uniform-form} as a
statement about $\cE^1$.  The finite-energy substitute is a different
argument, carried out in
\Cref{chap:bfj-energy,chap:bfj-finite-energy-measures} by truncation, and it
uses this lemma only through the truncated potentials
$\max\{\varphi,-t\}$, for which $M=t$.
\end{warning}

\section{Negligible, polar, and pluripolar sets}
\label{sec:ch25-negligible}

Three notions of smallness are in play, and the literature uses the words
for them inconsistently.  This section fixes the terminology once, proves the
implications that hold, and records what is not known.

\begin{definition}[Three kinds of small set]
\label{def:ch25-set-classes}
\index{negligible set}\index{polar set!BFJ}\index{pluripolar set!BFJ}
Let $E\subseteq X^{\An}$.
\begin{enumerate}
\item\label{itm:ch25-cls:neg}
 $E$ is \emph{negligible} if $\Capa^*_\theta(E)=0$, as in
 \Cref{def:bfj-capacity}.
\item\label{itm:ch25-cls:plp}
 $E$ is \emph{pluripolar} if $E\subseteq\{v=-\infty\}$ for some
 $v\in\PSH(X,\theta)$ with $v\not\equiv-\infty$; that is, if $E$ is
 contained in the polar set of a $\theta$-psh function in the sense of
 \Cref{def:ch21-polar-singularity}.
\item\label{itm:ch25-cls:sig}
 $E$ is \emph{$\sigma$-pluripolar} if it is a countable union of
 pluripolar sets.
\end{enumerate}
\end{definition}

\begin{remark}[Dependence on the reference form]
\label{rem:ch25-independence-of-theta}
All three notions are defined relative to the fixed form
\eqref{eq:ch25-reference-form}, and we make no claim that they are
independent of it.  One direction is easy and occasionally useful.  Suppose
$\theta'$ is another closed model form and $\varepsilon>0$ is such that
$\theta'-\varepsilon\theta$ is semipositive.  If $\psi$ is a
$\theta$-psh model function, then
$\theta'+\ddc(\varepsilon\psi)
 =(\theta'-\varepsilon\theta)+\varepsilon(\theta+\ddc\psi)$
is a sum of two semipositive classes on a common model, hence semipositive,
so $\varepsilon\psi$ is $\theta'$-psh; passing to decreasing limits
(\Cref{thm:ch21-decreasing-stability},
\Cref{thm:bfj-decreasing-regularization}) shows that $\varepsilon v$ is
$\theta'$-psh for every $v\in\PSH(X,\theta)$, with the same polar set.
Hence every $\theta$-pluripolar set is $\theta'$-pluripolar.  Nothing
below uses this, and no comparison of $\Capa_\theta$ with
$\Capa_{\theta'}$ is asserted anywhere in the book.
\end{remark}

\begin{proposition}
\label{prop:ch25-negligible-sigma-ideal}
\begin{enumerate}
\item\label{itm:ch25-nsi:ideal}
 A subset of a negligible set is negligible, and a countable union of
 negligible sets is negligible.
\item\label{itm:ch25-nsi:null}
 Every negligible set is contained in a Borel set of measure zero for every
 mixed measure
 \[
  \MA_\theta(\varphi_1,\ldots,\varphi_n),
 \]
 where the $\varphi_i$ are bounded $\theta$-psh functions.
\item\label{itm:ch25-nsi:interior}
 A negligible set has empty interior and contains no divisorial point.
\end{enumerate}
\end{proposition}

\begin{proof}
\cref{itm:ch25-nsi:ideal} is
\Cref{lem:ch25-inner-outer}\,\cref{itm:ch25-io:mono} and
\Cref{cor:ch25-outer-subadditivity}; equivalently, it is
\Cref{prop:polar-sigma-ideal}\,(1)--(2) applied to the countably subadditive
capacity $\Capa^*_\theta$.  \cref{itm:ch25-nsi:null} is
\Cref{cor:ch25-outer-ma-bound}\,\cref{itm:ch25-om:null}, which supplies the
hypothesis $\mu\leq C\Capa$ of \Cref{prop:polar-sigma-ideal}\,(3);
inner regularity of a Radon measure is automatic by
\eqref{eq:radon-regularity}.  \cref{itm:ch25-nsi:interior} is
\Cref{cor:ch25-negligible-empty-interior} and
\Cref{lem:divisorial-positive-capacity}.
\end{proof}

\begin{theorem}
\label{thm:ch25-pluripolar-negligible}
Every $\sigma$-pluripolar subset of $X^{\An}$ is negligible.  Moreover
the vanishing is quantitative for a finite family: if
$E\subseteq\bigcup_{j=1}^N\{v_j=-\infty\}$ with
$v_j\in\PSH(X,\theta)$ normalized by $\sup v_j=0$, then for every
$t>0$,
\begin{equation}
\label{eq:ch25-sigma-pluripolar-rate}
 \Capa^*_\theta(E)\leq\sum_{j=1}^N
 \Capa_\theta\bigl(\{v_j<-t\}\bigr)
 \leq\frac{N(C_\theta+nV)}{t}.
\end{equation}
\end{theorem}

\begin{proof}
By \Cref{cor:ch25-polar-negligible} each $\{v_j=-\infty\}$ is negligible.
By \Cref{prop:ch25-negligible-sigma-ideal}\,\cref{itm:ch25-nsi:ideal} a
countable union of such sets is negligible, and so is any subset of it.  For
\eqref{eq:ch25-sigma-pluripolar-rate}, note that
$E\subseteq\bigcup_{j=1}^N\{v_j<-t\}$ for every $t>0$, and apply
\Cref{cor:ch25-outer-subadditivity} followed by
\eqref{eq:ch25-sublevel-capacity}.  Letting $t\to\infty$ recovers
$\Capa^*_\theta(E)=0$ in the finite case.
\end{proof}

\begin{corollary}
\label{cor:ch25-polar-countable}
A countable union of pluripolar sets is negligible.  In particular the
family of $\sigma$-pluripolar sets is contained in the
$\sigma$-ideal of negligible sets, and every countable union of polar sets
of $\theta$-psh functions is a set of outer capacity zero.
\end{corollary}

\begin{proof}
Immediate from \Cref{thm:ch25-pluripolar-negligible}.
\end{proof}

\begin{warning}{The converse is not proved here}
Whether a negligible set is pluripolar, or even
$\sigma$-pluripolar, is not established in this book.  In complex
pluripotential theory the corresponding statement is true and classical: a
set of zero Monge--Amp\`ere capacity on a compact K\"ahler manifold is
pluripolar, by the theory of extremal functions of \cite{BT82} in the local
case and \cite{GZ05} in the compact case, both of which construct a
\emph{global} psh function with prescribed $-\infty$ locus out of local
data.  The construction rests on Josefson-type gluing, which in turn rests on
a local-to-global principle for psh functions.  In the BFJ theory no such
principle is available: $\PSH(X,\theta)$ is defined globally, by
positivity on all models (\Cref{def:ch21-bfj-psh}), and
\Cref{prop:ch21-bfj-vs-cld} records that a global BFJ-psh function is not
characterized by a local condition.  We therefore state only what is proved:
pluripolar $\Rightarrow$ negligible, and no converse.  This is the
discipline announced in the introduction to \Cref{part:bfj}, and it costs
nothing, because every use of the words below is in the proved direction.
\end{warning}

\begin{remark}[The converse holds on curves, for a different capacity]
\label{rem:ch25-curve-converse}
On a smooth Berkovich curve the converse is available, and its proof does
not go through Josefson gluing at all.  Thuillier shows that a subset of a
connected open $\Omega$ with $\partial\Omega$ a nonempty set of
quasimonomial points has vanishing outer relative capacity if and only if it
is contained in the $-\infty$ locus of a single subharmonic function on
$\Omega$ \cite[Th\'eor\`eme~3.6.11]{Thu05}; this is
\Cref{thm:ch25-thuillier-polar}.  The mechanism is the relative extremal
function $u^*_E$, which on a curve is available with no regularization
theorem and no local-to-global principle, because subharmonicity there is
already a local condition (\Cref{thm:wanner-curve-comparison}).  What the
present chapter lacks in higher dimension is precisely that: a local
characterization of membership in $\PSH(X,\theta)$.
\Cref{rem:ch25-thuillier-vs-bfj} records that
\Cref{thm:ch25-thuillier-polar} is a statement about Thuillier's capacity
and is not being asserted for $\Capa_\theta$, not even when $n=1$.
\end{remark}

\begin{remark}[Terminology: a translation table]
\label{rem:ch25-terminology}
The words are used as follows in this book.
\begin{center}
\small
\begin{tabular}{p{.24\textwidth}|p{.34\textwidth}|p{.34\textwidth}}
\toprule
\textbf{Word}&\textbf{Meaning here}&\textbf{Fixed at}\\
\midrule
polar set of $\varphi$
 &$\{\varphi=-\infty\}$ for one $\varphi\in\PSH(X,\theta)$
 &\Cref{def:ch21-polar-singularity}\\
pluripolar
 &contained in the polar set of some $\varphi\in\PSH(X,\theta)$
 &\Cref{def:ch25-set-classes}\,\cref{itm:ch25-cls:plp}\\
$\sigma$-pluripolar
 &countable union of pluripolar sets
 &\Cref{def:ch25-set-classes}\,\cref{itm:ch25-cls:sig}\\
negligible
 &outer $\theta$-capacity zero
 &\Cref{def:bfj-capacity}\\
polar for the capacity
 &synonym of negligible, in the sense of
  \Cref{sec:capacity-preview}
 &\Cref{prop:ch25-negligible-sigma-ideal}\\
\bottomrule
\end{tabular}
\end{center}
The last line uses the abstract capacity terminology of
\Cref{sec:capacity-preview}.  The distinction announced in
\Cref{sec:ch21-polar} is retained: subsets of psh pole loci are
pluripolar, and their countable unions are negligible by
\Cref{thm:ch25-pluripolar-negligible}; they are not proved to exhaust the
negligible sets.  The
classical word \emph{negligible} for a set of the form
$\{\sup_i u_i<(\sup_iu_i)^*\}$ is not used; that set is called the
\emph{regularization defect} here, and it is the subject of
\Cref{prop:ch25-regularization-defect}.
\end{remark}

\begin{proposition}
\label{prop:ch25-regularization-defect}
Let $(\varphi_i)_{i\in I}$ be a nonempty family in $\PSH(X,\theta)$,
uniformly bounded above, let $\varphi=\sup_i\varphi_i$ and let
$\varphi^*$ be its upper semicontinuous regularization.  Then the
\emph{regularization defect}
\[
                 D:=\{\varphi<\varphi^*\}
\]
contains no quasimonomial point; in particular it contains no divisorial
point, it has empty interior, and its complement is dense.  If $\varphi$
is upper semicontinuous---for instance if the family is finite, or is a
family with a largest member---then $D=\varnothing$.
\end{proposition}

\begin{proof}
By \Cref{thm:ch23-upper-envelope}, $\varphi^*\in\PSH(X,\theta)$ and
$\varphi^*=\varphi$ at every quasimonomial point, so $D$ is disjoint
from $\bigcup_{\mathcal X}\Delta_{\mathcal X}$.  That union is dense in
$X^{\An}$ by \Cref{thm:skeleton-inverse-limit}, so $D$ contains no
nonempty open set.  Divisorial points are quasimonomial
(\Cref{rem:quasimonomial-vs-divisorial}).  If $\varphi$ is upper
semicontinuous, then $\varphi=\varphi^*$ by the last clause of
\Cref{prop:usc-majorant} and $D=\varnothing$; a finite family has a
maximum, which is upper semicontinuous, and a family with a largest member
has that member as its supremum.
\end{proof}

\begin{warning}{The regularization defect is not proved to be negligible}
The classical theorem of \cite{BT82}, that negligible sets in the sense of
regularization defects are pluripolar, is not transported here, and neither
is the weaker statement that $D$ has capacity zero.  The obstruction is
visible in the proof one would attempt.  Applying
\eqref{eq:ch25-uniform-form} to the pair $(\varphi_j,\varphi^*)$ for an
increasing sequence $\varphi_j\nearrow\varphi$ requires
$\int(\varphi^*-\varphi_j)\MA_\theta(\varphi_j/2)\to0$, in which the
measure \emph{moves with $j$}.  The integrands have no useful
semicontinuity in general, and the weighted continuity of
\Cref{thm:bounded-ma-extension} applies to decreasing psh weights and
decreasing psh entries, whereas $\varphi_j$ increases.  Thus neither the
fixed-measure monotone theorem nor the weighted Monge--Amp\`ere theorem gives
this limit.  What \emph{is} proved above is that
$D$ misses every skeleton, which is a strong structural statement with no
analogue in the complex theory, and which suffices for every use made of
upper envelopes in \Cref{chap:bfj-envelopes}.  We do not assert
$\Capa^*_\theta(D)=0$.
\end{warning}

\begin{remark}[Comparison with the CLD negligibility theorem]
\label{rem:ch25-cld-comparison}
\Cref{thm:cld-polar-negligibility} asserts that a Chambert--Loir--Ducros
mixed measure of continuous, locally approachable psh functions does not
charge the polar locus of a locally approachable psh function.  It is a
\emph{local} statement about a \emph{different} product, proved by a weighted
Chern--Levine--Nirenberg estimate on affinoid charts, and its hypotheses
(goodness, local approachability) have no counterpart here.  The BFJ
statement \Cref{cor:ch25-polar-negligible} is global, applies to all bounded
BFJ-psh potentials without an approachability hypothesis, and controls a
supremum over potentials rather than one measure.  Neither implies the other,
and the two are compared only on the common locus identified in
\Cref{thm:bfj-cld-ma-common-locus}: for semipositive model metrics, both
measures are the same atomic measure, and both statements reduce to the
observation that an atomic measure supported at divisorial points cannot
charge a polar set, since polar sets contain no divisorial point
(\Cref{prop:ch21-polar-structure}\,\cref{itm:ch21-pol:skeleta}).
\end{remark}

\section{Quasicontinuity and a sequential regularization}
\label{sec:ch25-quasicontinuity}

\begin{definition}
\label{def:ch25-quasicontinuous}
A function $h:X^{\An}\to\overline{\mathbf R}$ is
\emph{quasicontinuous} if, for every $\varepsilon>0$, there is an open set
$G$ with $\Capa_\theta(G)<\varepsilon$ such that
$h|_{X^{\An}\setminus G}$ is continuous for the order topology on
$\overline{\mathbf R}$.  Equivalently, one may use the metric $\rho$ fixed
before \Cref{def:capacity-convergence}.  For real-valued functions this is
the definition used below in the bounded theory.
\end{definition}

This is \Cref{sec:capacity-preview} specialized to $\Capa_\theta$, and the
formal calculus proved there for real-valued functions---finite sums,
maxima, products, uniform limits, and simultaneous continuity outside a
small-capacity open set for a countable family---is available at once.
For extended-real functions one must avoid undefined expressions such as
$+\infty-\infty$.  The real-valued calculus applies here, because $\Capa_\theta$ is countably
subadditive by \eqref{eq:capacity-subadditivity}.  What is \emph{not}
available from the abstract theory is the existence of a single
quasicontinuous function; that is the theorem below.

\begin{theorem}
\label{thm:bounded-quasicontinuity}
Every bounded $\theta$-psh function is quasicontinuous. More precisely,
for bounded $\varphi$ and $\varepsilon>0$, there are an open set $G$
with $\Capa_\theta(G)<\varepsilon$ and a decreasing sequence of psh model
functions $\varphi_m$ converging uniformly to $\varphi$ on $G^c$.
\end{theorem}

\begin{proof}
The proof has four steps.  Steps~1--2 are the analytic content; Steps~3--4
are the extraction of a sequence and the verification that the exceptional
set can be taken open.

\emph{Step 1: normalization and a decreasing model net.}  Let $\varphi$ be
bounded $\theta$-psh.  Adding a constant, which changes neither
quasicontinuity nor any capacity, we may assume $\varphi\leq0$; let
$M\geq1$ be such that $-M\leq\varphi\leq0$.  By
\Cref{thm:bfj-decreasing-regularization} there is a decreasing net
$(\varphi_\alpha)_{\alpha\in A}$ of $\theta$-psh model functions with
$\varphi_\alpha\searrow\varphi$ pointwise.  One cannot simply truncate the
approximants at $0$, since a minimum of $\theta$-psh functions is not
$\theta$-psh (\Cref{ex:ch20-min-fails}); instead we use the constants
$c_\alpha:=\sup_{X^{\An}}\varphi_\alpha$, which decrease to
$\sup\varphi\leq0$ by
\Cref{lem:ch25-admissible-calculus}\,\cref{itm:ch25-adm:model}.  Discarding an
initial segment of the net we may therefore assume $c_\alpha\leq1$, so that
\begin{equation}
\label{eq:ch25-qc-normalization}
 -M\leq\varphi\leq\varphi_\alpha\leq1
 \qquad\text{for all }\alpha .
\end{equation}
Replacing $M$ by $M+1$ and each $\varphi_\alpha$ by
$\varphi_\alpha-1$ and $\varphi$ by $\varphi-1$---a translation of all
functions by the same constant, which changes nothing---we may and do assume
\begin{equation}
\label{eq:ch25-qc-normalization-two}
 -M\leq\varphi\leq\varphi_\alpha\leq0
 \qquad\text{for all }\alpha ,\qquad M\geq1 .
\end{equation}

\emph{Step 2: the exchange, uniformly over the admissible family.}  Fix
$\alpha$.  The pair $(\varphi,\varphi_\alpha)$ satisfies the hypotheses
of \Cref{lem:fundamental-capacity-estimate}, so for every
$u\in\mathcal U_\theta$, taking $u_1=\cdots=u_n=u$ there,
\begin{equation}
\label{eq:ch25-qc-uniform}
 \int(\varphi_\alpha-\varphi)\,\MA_\theta(u)
 \ \leq\ 4MV^{1-2^{-n}}\,\Lambda_\alpha^{2^{-n}},
 \qquad
 \Lambda_\alpha:=\int(\varphi_\alpha-\varphi)\,
   \MA_\theta(\varphi/2).
\end{equation}
The weighted continuity in \Cref{thm:bounded-ma-extension}, applied with
weight $\varphi_\alpha\searrow\varphi$ and with all measure entries fixed at
$\varphi/2$, gives
\begin{equation}
\label{eq:ch25-qc-lambda}
 \Lambda_\alpha\longrightarrow0 .
\end{equation}
Combining \eqref{eq:ch25-qc-uniform} and \eqref{eq:ch25-qc-lambda},
\begin{equation}
\label{eq:ch25-qc-sup}
 \Theta_\alpha:=\sup_{u\in\mathcal U_\theta}
 \int(\varphi_\alpha-\varphi)\,\MA_\theta(u)\longrightarrow0 .
\end{equation}
This uniformity is the reason the chapter is organized as it is:
\eqref{eq:ch25-qc-sup} is a statement uniform over the admissible family, and
it has been obtained from the weighted Monge--Amp\`ere continuity theorem.

\emph{Step 3: extraction of a sequence of indices.}  Let $\varepsilon>0$.
For each $m\geq1$ put
\[
 A_m:=\Bigl\{\alpha\in A\ :\
   \Theta_\alpha<\frac{2^{-m}\varepsilon}{m}\Bigr\}.
\]
By \eqref{eq:ch25-qc-sup} and monotonicity of $\alpha\mapsto\Theta_\alpha$
along the net, each $A_m$ contains a tail of $A$, hence is cofinal.
Choose $\alpha_1\in A_1$; given $\alpha_{m}$, use directedness of $A$
and cofinality of $A_{m+1}$ to choose
$\alpha_{m+1}\geq\alpha_m$ with $\alpha_{m+1}\in A_{m+1}$.  The resulting
sequence $(\alpha_m)_{m\geq1}$ is increasing, so
$\varphi_m:=\varphi_{\alpha_m}$ is a \emph{decreasing sequence} of
$\theta$-psh model functions with
\begin{equation}
\label{eq:ch25-qc-theta-bound}
 \Theta_{\alpha_m}<\frac{2^{-m}\varepsilon}{m},\qquad m\geq1 .
\end{equation}
No claim is made that $(\alpha_m)$ is cofinal in $A$; indeed
\Cref{ex:ch22-no-countable-cofinal} shows that no countable subset of the
model system need be cofinal.  This observation does not rule out a
cofinal sequence for a particular regularizing net.  What is used is only
\eqref{eq:ch25-qc-theta-bound}.

\emph{Step 4: the exceptional set.}  Put
\[
 G_m:=\bigl\{\varphi_m-\varphi>m^{-1}\bigr\},
 \qquad
 G:=\bigcup_{m\geq1}G_m .
\]
Each $G_m$ is open: $\varphi_m$ is continuous, being a model function,
and $-\varphi$ is lower semicontinuous, so $\varphi_m-\varphi$ is lower
semicontinuous and its strict superlevel sets are open.  For
$u\in\mathcal U_\theta$, Chebyshev's inequality for the positive measure
$\MA_\theta(u)$ gives
\[
 \MA_\theta(u)(G_m)
 \leq m\int(\varphi_m-\varphi)\,\MA_\theta(u)
 \leq m\,\Theta_{\alpha_m}<2^{-m}\varepsilon ,
\]
and the bound is independent of $u$, so
$\Capa_\theta(G_m)\leq2^{-m}\varepsilon$.  By countable subadditivity
\eqref{eq:capacity-subadditivity},
\[
 \Capa_\theta(G)\leq\sum_{m\geq1}\Capa_\theta(G_m)
 \leq\sum_{m\geq1}2^{-m}\varepsilon=\varepsilon .
\]
Strict inequality is obtained by running the argument with
$\varepsilon/2$.  On $X^{\An}\setminus G$ we have, for every $m$,
\[
        0\leq\varphi_m-\varphi\leq m^{-1},
\]
so $\varphi_m\to\varphi$ uniformly there.  A uniform limit of continuous
functions on the subspace $X^{\An}\setminus G$ is continuous, so
$\varphi|_{X^{\An}\setminus G}$ is continuous and $\varphi$ is
quasicontinuous in the sense of
\Cref{def:ch25-quasicontinuous}.  This proves both assertions.
\end{proof}

\begin{remark}[Exactly which hypotheses are used]
\label{rem:ch25-qc-hypotheses}
Boundedness of $\varphi$ enters twice, and in different ways.  It is used
in Step~1 to obtain a finite $M$, and $M$ appears linearly in
\eqref{eq:ch25-qc-uniform}; and it is used implicitly in Step~2, since
$\MA_\theta(\varphi/2)$ is defined only for bounded $\varphi$ at this
stage (\Cref{thm:bounded-ma-extension}).  Semipositivity of $\theta$ is
used to know that $\mathcal U_\theta\neq\varnothing$ and that
$\varphi/2$ is $\theta$-psh.  Residue characteristic zero is used exactly
once, through \Cref{thm:bfj-decreasing-regularization}, to produce the model
net at all.  Nothing else is needed; in particular no continuity of
envelopes and no compactness of $\PSH(X,\theta)$ is used here.

The theorem is a bounded, real-valued statement.  Its extended-real-valued
counterpart for arbitrary $\theta$-psh functions is
\Cref{thm:ch25-psh-quasicontinuity}, proved by truncation in
\Cref{sec:ch25-e1}; none of the bounded estimates is extended there.
\end{remark}

\begin{corollary}
\label{cor:ch25-quasicontinuity-calculus}
Let $\varphi_1,\varphi_2,\ldots$ be countably many bounded $\theta$-psh
functions and let $\varepsilon>0$.  Then there is a single open set $G$
with $\Capa_\theta(G)<\varepsilon$ such that every $\varphi_j$ restricts
to a continuous function on $X^{\An}\setminus G$.  Consequently every
finite sum, difference, maximum, minimum or product of bounded
$\theta$-psh functions is quasicontinuous, as is $\chi\circ\varphi$ for
$\chi:\mathbf R\to\mathbf R$ continuous, and so is every uniform limit of
such functions.
\end{corollary}

\begin{proof}
\Cref{thm:bounded-quasicontinuity} makes each $\varphi_j$ quasicontinuous;
$\Capa_\theta$ is countably subadditive by
\eqref{eq:capacity-subadditivity}; now apply
\Cref{prop:quasicontinuity-calculus}\,(1)--(3) verbatim.
\end{proof}

\begin{remark}[What quasicontinuity does not say]
\label{rem:ch25-qc-not}
The exceptional set $G$ depends on $\varepsilon$ and, for a family, on
the family; a single negligible set need not suffice.  For example, in the branch
sum of \Cref{ex:ch24-branch-sum}, every selected point $x_k$ and
the Gauss point $\zeta$ is divisorial, so
\Cref{lem:divisorial-positive-capacity} prevents any negligible set from
containing them.  In its complement $x_k\to\zeta$ still holds, while
$\varphi(x_k)=-1$ and $\varphi(\zeta)=0$, precluding continuity there.
This argument uses a sequence that remains in the complement, rather than
a claim that removing a set must remove the original discontinuity locus.
Moreover
$\varphi|_{X^{\An}\setminus G}$ is continuous \emph{as a function on the
subspace}; this is weaker than continuity of $\varphi$ at the points of
that subspace, as \Cref{sec:capacity-preview} already warned.  Finally,
\Cref{rem:quasicontinuity-limits} recorded two statements that do not follow
formally, one of which---that a decreasing limit of quasicontinuous
functions is quasicontinuous---is proved here only in the special case of
\Cref{thm:bounded-quasicontinuity}, where the bounded limit is again
$\theta$-psh.
\end{remark}

\begin{corollary}
\label{cor:psh-sequential-regularization}
Every $\theta$-psh function, bounded or not, is the pointwise decreasing
limit of a sequence of psh model functions.
\end{corollary}

\begin{proof}
Let $\varphi\in\PSH(X,\theta)$.  Adding a constant we may assume
$\sup\varphi\leq0$.  By \Cref{thm:bfj-decreasing-regularization} choose a
decreasing net $(\varphi_\alpha)_{\alpha\in A}$ of $\theta$-psh model
functions with $\varphi_\alpha\searrow\varphi$; as in Step~1 of the proof
of \Cref{thm:bounded-quasicontinuity} we may assume $\varphi_\alpha\leq0$
for all $\alpha$, after discarding an initial segment and translating all
functions by one and the same constant.  For $m\geq1$ write
\[
 \varphi^{(m)}:=\max\{\varphi,-m\},\qquad
 \varphi^{(m)}_\alpha:=\max\{\varphi_\alpha,-m\} ,
\]
both $\theta$-psh by
\Cref{thm:ch21-stability}\,\cref{itm:ch21-stab:max}, both bounded with values
in $[-m,0]$, and satisfying
$\varphi^{(m)}\leq\varphi^{(m)}_\alpha\leq0$.

\emph{Step 1: the diagonal choice.}  Fix $m$.  The weighted decreasing-net
continuity of \Cref{thm:bounded-ma-extension}, applied to the weight
$\varphi^{(m)}_\alpha\searrow\varphi^{(m)}$ and fixed measure entries
$\varphi^{(m)}/2$, gives
\[
 \int(\varphi^{(m)}_\alpha-\varphi^{(m)})
       \MA_\theta(\varphi^{(m)}/2)\longrightarrow0.
\]
Hence the set of indices $\alpha$ with
\begin{equation}
\label{eq:diagonal-capacity-choice}
\begin{split}
0\leq{}&\int\bigl(\max\{\varphi_{\alpha},-m\}-\varphi^{(m)}\bigr)
               \MA_\theta(\varphi^{(m)}/2)\\
&\leq \bigl(4m^2V^{1-2^{-n}}\bigr)^{-2^n}
\end{split}
\end{equation}
is cofinal in $A$.  Exactly as in Step~3 of the proof of
\Cref{thm:bounded-quasicontinuity}, directedness of $A$ allows us to choose
indices $\alpha_1\leq\alpha_2\leq\cdots$ with $\alpha_m$ satisfying
\eqref{eq:diagonal-capacity-choice} for that $m$.  Put
$\varphi_m:=\varphi_{\alpha_m}$; this is a decreasing sequence of
$\theta$-psh model functions.

\emph{Step 2: the uniform estimate at level $m$.}  Apply
\Cref{lem:fundamental-capacity-estimate} to the pair
$\varphi^{(m)}\leq\varphi^{(m)}_{\alpha_m}$, with $M=m\geq1$ and with
arbitrary $u_1,\ldots,u_n\in\mathcal U_\theta$.  Its hypotheses hold:
$-m\leq\varphi^{(m)}\leq\varphi^{(m)}_{\alpha_m}\leq0$ and
$-m\leq u_i\leq0$.  Using \eqref{eq:diagonal-capacity-choice},
\[
 \begin{aligned}
 \int\bigl(\varphi^{(m)}_{\alpha_m}-\varphi^{(m)}\bigr)
       \MA_\theta(u_1,\ldots,u_n)
 &\leq4mV^{1-2^{-n}}\cdot
   \Bigl(\bigl(4m^2V^{1-2^{-n}}\bigr)^{-2^n}\Bigr)^{2^{-n}}
 \\
 &=\frac{4mV^{1-2^{-n}}}{4m^2V^{1-2^{-n}}}
  =\frac1m .
 \end{aligned}
\]

\emph{Step 3: passing to divisorial points.}  Let $x$ be a divisorial
point, so that $\varphi(x)>-\infty$ by
\Cref{prop:ch21-first-consequences}\,\cref{itm:ch21-fc:quasimonomial}.  For
$m\geq-\varphi(x)$ both truncations are inactive at $x$: indeed
$\varphi(x)\geq-m$ gives $\varphi^{(m)}(x)=\varphi(x)$, and
$\varphi_{\alpha_m}(x)\geq\varphi(x)\geq-m$ gives
$\varphi^{(m)}_{\alpha_m}(x)=\varphi_{\alpha_m}(x)$.  Taking
$u_1=\cdots=u_n=u$ in Step~2 and restricting the integral to the single
point $x$,
\[
 \bigl(\varphi_{m}(x)-\varphi(x)\bigr)\,\MA_\theta(u)(\{x\})
 \leq\int\bigl(\varphi^{(m)}_{\alpha_m}-\varphi^{(m)}\bigr)\MA_\theta(u)
 \leq\frac1m ,
\]
the first inequality because the integrand is nonnegative.  Taking the
supremum over $u\in\mathcal U_\theta$,
\[
 0\leq\bigl(\varphi_{m}(x)-\varphi(x)\bigr)\Capa_\theta(\{x\})
 \leq\frac1m .
\]
By \Cref{lem:divisorial-positive-capacity}, $\Capa_\theta(\{x\})>0$, so
$\varphi_m(x)\to\varphi(x)$.

\emph{Step 4: conclusion.}  The sequence $(\varphi_m)$ is decreasing and
bounded below by $\varphi$, so $\chi:=\lim_m\varphi_m=\inf_m\varphi_m$
exists pointwise and satisfies $\chi\geq\varphi$.  Since
$\chi\geq\varphi\not\equiv-\infty$, the decreasing-limit stability theorem
\Cref{thm:ch21-decreasing-stability} gives $\chi\in\PSH(X,\theta)$.  By
Step~3, $\chi=\varphi$ at every divisorial point, and two $\theta$-psh
functions that agree at all divisorial points are equal by
\Cref{prop:ch21-divisorial-order}.  Hence $\varphi_m\searrow\varphi$
pointwise on $X^{\An}$.
\end{proof}

\begin{remark}
\label{rem:ch25-sequential-specific}
\Cref{cor:psh-sequential-regularization} is a theorem specific to
this psh class. It does not produce a countable cofinal system of models and
does not make $X^{\An}$ first countable.
\end{remark}

\begin{remark}[What the sequential regularization improves, and what it does
not]
\label{rem:ch25-sequential-vs-ch22}
\Cref{rem:ch22-decreasing-selection-improvement} recorded that
\Cref{chap:bfj-compactness} produces a decreasing \emph{sequence} of model
approximants only in the continuous case
(\Cref{cor:ch22-continuous-decreasing-sequence}), and that ``no argument in
this chapter produces a sequence'' in the singular case.
\Cref{cor:psh-sequential-regularization} supplies exactly that missing
argument, and it does so using an ingredient unavailable in
\Cref{chap:bfj-compactness}: the positivity of the capacity of a single
divisorial point.  Three caveats attach to the statement.

First, convergence is pointwise globally and uniform on each fixed dual
complex: the restrictions decrease to the continuous skeletal restriction
of $\varphi$, so \Cref{prop:dini-nets} applies.  No global uniform
convergence is claimed unless $\varphi$ is continuous.  Second, the
sequence is adapted to $\varphi$ and does not give a countably cofinal
model system; \Cref{ex:ch22-countable-blind-spot} explains the limitation
of a fixed countable determining family.  Third, the proof uses the
bounded Monge--Amp\`ere operator on the truncations
$\varphi^{(m)}$ and therefore does not survive the removal of
\Cref{thm:bounded-ma-extension}; it is not a statement of pure topology.
\end{remark}

Before the two limit lemmas we settle a measurability point once and for
all.

\begin{lemma}
\label{lem:ch25-quasicontinuous-measurable}
Let $h:X^{\An}\to\mathbf R$ be quasicontinuous.  Then there is a negligible
$G_\infty$ such that $h|_{X^{\An}\setminus G_\infty}$ is Borel
measurable.  Consequently $h$ is measurable, and integrable when bounded,
for the completion of every Monge--Amp\`ere measure of bounded
$\theta$-psh potentials.
\end{lemma}

\begin{proof}
For each $k\geq1$ choose an open $G_k$ with
$\Capa_\theta(G_k)<1/k$ such that $h|_{X^{\An}\setminus G_k}$ is
continuous, and put $G_\infty:=\bigcap_kG_k$, which is negligible by
monotonicity.  For an open $U\subseteq\mathbf R$, the set
$h^{-1}(U)\cap(X^{\An}\setminus G_k)$ is relatively open in the closed set
$X^{\An}\setminus G_k$, hence Borel in $X^{\An}$; taking the union over
$k$ shows that $h^{-1}(U)\setminus G_\infty$ is Borel.  The last
assertion follows from
\Cref{cor:ch25-outer-ma-bound}\,\cref{itm:ch25-om:null}, which makes
$G_\infty$ null for every such measure.
\end{proof}

\begin{lemma}
\label{lem:quasicontinuous-tests}
Let $h$ be bounded and quasicontinuous.  Let $u$ be a bounded
$\theta$-psh function, and suppose $u_\alpha$ are uniformly bounded
$\theta$-psh functions and
$\MA_\theta(u_\alpha)\to\MA_\theta(u)$ weakly. Then
\[
                    \int h\,\MA_\theta(u_\alpha)
                    \longrightarrow\int h\,\MA_\theta(u).
\]
\end{lemma}

\begin{proof}
Let $M\geq1$ be such that $-M\leq u_\alpha\leq0$ for all $\alpha$ and
$-M\leq u\leq0$; translating all potentials by a common constant, which
changes no Monge--Amp\`ere measure, we may assume this normalization.  Let
$C:=\sup|h|<\infty$ and let $\eta>0$.

By quasicontinuity choose an open $G$ with
$\Capa_\theta(G)<\eta$ such that $h|_{X^{\An}\setminus G}$ is continuous.
The space $X^{\An}$ is compact Hausdorff, hence normal, and
$X^{\An}\setminus G$ is closed; by the Tietze extension theorem there is
$H\in C^0(X^{\An})$ with $H=h$ on $X^{\An}\setminus G$ and
$\sup|H|\leq C$.  Then $|h-H|\leq2C$ everywhere and $h-H=0$ off
$G$, so for every bounded $\theta$-psh $w$ with $-M\leq w\leq0$,
\begin{equation}
\label{eq:ch25-tietze-error}
 \begin{aligned}
 \Bigl|\int h\,\MA_\theta(w)-\int H\,\MA_\theta(w)\Bigr|
 &\leq2C\,\MA_\theta(w)(G)\\
 &\leq2CM^n\Capa_\theta(G)<2CM^n\eta ,
 \end{aligned}
\end{equation}
by \eqref{eq:ma-capacity-bound}.  Applying
\eqref{eq:ch25-tietze-error} to $w=u_\alpha$ and to $w=u$, and using
weak convergence $\int H\MA_\theta(u_\alpha)\to\int H\MA_\theta(u)$ for the
continuous test function $H$, we get
\[
 \varlimsup_\alpha\Bigl|\int h\,\MA_\theta(u_\alpha)
   -\int h\,\MA_\theta(u)\Bigr|\leq4CM^n\eta .
\]
Let $\eta\downarrow0$.
\end{proof}

\begin{lemma}
\label{lem:ch25-joint-limit}
Let $M\geq1$ and let $(\mu_\alpha)$ be a net of Monge--Amp\`ere measures
$\mu_\alpha=\MA_\theta(u_{\alpha,1},\ldots,u_{\alpha,n})$ with all
$-M\leq u_{\alpha,i}\leq0$, converging weakly to a measure
$\mu=\MA_\theta(u_1,\ldots,u_n)$ of the same kind.  Let $(h_\alpha)$ be
quasicontinuous functions with $\sup_\alpha\|h_\alpha\|_\infty\leq C$,
converging in capacity---in the sense of
\Cref{def:ch25-convergence-in-capacity}, stated in the next section, of which
only the definition is used here---to a quasicontinuous $h$ with
$\|h\|_\infty\leq C$.  All the integrals below are defined by
\Cref{lem:ch25-quasicontinuous-measurable}.  Then
\[
 \int h_\alpha\,\rd\mu_\alpha\longrightarrow\int h\,\rd\mu .
\]
\end{lemma}

\begin{proof}
Write $M':=M^nn^n/n!$, so that $\nu(E)\leq M'\Capa^*_\theta(E)$ for every
set $E$ and every measure $\nu$ of the stated kind, by
\Cref{cor:ch25-outer-ma-bound}\,\cref{itm:ch25-om:mixed}.  Let
$\eta,\delta>0$.  Choose an open $G$ with $\Capa_\theta(G)<\eta$ off
which $h$ is continuous, and $H\in C^0(X^{\An})$ with $H=h$ off $G$
and $\|H\|_\infty\leq C$, by Tietze as above.  Then
\[
\begin{split}
 \Bigl|\int h_\alpha\rd\mu_\alpha-\int h\,\rd\mu\Bigr|
 \leq{}&\Bigl|\int(h_\alpha-H)\rd\mu_\alpha\Bigr|
       +\Bigl|\int H\rd\mu_\alpha-\int H\rd\mu\Bigr|\\
      &+\Bigl|\int(H-h)\rd\mu\Bigr| .
\end{split}
\]
The third term is at most $2CM'\eta$, since $H-h$ vanishes off $G$.
The second tends to $0$ by weak convergence.  For the first, let
$E_\alpha:=\{|h_\alpha-h|>\delta\}$; off $G\cup E_\alpha$ we have
$|h_\alpha-H|=|h_\alpha-h|\leq\delta$, so
\[
 \Bigl|\int(h_\alpha-H)\rd\mu_\alpha\Bigr|
 \leq\delta V+2C\mu_\alpha(G\cup E_\alpha)
 \leq\delta V+2CM'\bigl(\eta+\Capa^*_\theta(E_\alpha)\bigr).
\]
Since $h_\alpha\to h$ in capacity, $\Capa^*_\theta(E_\alpha)\to0$.
Collecting,
\[
 \varlimsup_\alpha\Bigl|\int h_\alpha\rd\mu_\alpha-\int h\,\rd\mu\Bigr|
 \leq\delta V+4CM'\eta .
\]
Let $\delta,\eta\downarrow0$.
\end{proof}

\begin{remark}[Why the joint lemma is needed]
\label{rem:ch25-why-joint}
\Cref{lem:quasicontinuous-tests} moves a \emph{fixed} quasicontinuous test
function past a weakly convergent net of Monge--Amp\`ere measures.  In the
proof of the maximum principle the test function itself varies with the
index, being built from the approximants, and no reformulation avoids this:
the model identity \eqref{eq:model-maximum-principle} holds on the set where
the \emph{approximants} are ordered, not on the set where their limits are.
\Cref{lem:ch25-joint-limit} is exactly the tool that lets both the measure
and the test function move, at the cost of requiring convergence in capacity
for the latter---which is precisely what
\Cref{prop:decreasing-capacity-convergence} supplies for decreasing nets.
\end{remark}

\section{Convergence in capacity}
\label{sec:ch25-convergence}

\begin{definition}
\label{def:ch25-convergence-in-capacity}
Let $(Y,d)$ be a metric space.  A net of maps
$f_\alpha:X^{\An}\to Y$ \emph{converges in $\theta$-capacity} to
$f:X^{\An}\to Y$ if for every $\delta>0$,
\[
             \Capa^*_\theta(\{d(f_\alpha,f)>\delta\})\longrightarrow0.
\]
For real-valued maps we use the ordinary distance, and for
$\overline{\mathbf R}$-valued maps the metric $\rho$ of
\Cref{def:capacity-convergence}.  The two notions agree on uniformly bounded
real-valued families.
\end{definition}

By \Cref{lem:ch25-inner-outer}\,\cref{itm:ch25-io:open}, the outer capacity
in this definition equals the Borel-set capacity whenever the exceptional
set is open.  This is the case that occurs most often below: if $f_\alpha$ is
continuous, $f$ is upper semicontinuous, and $f_\alpha\geq f$, then
$\{f_\alpha-f>\delta\}$ is open.

\begin{proposition}
\label{prop:decreasing-capacity-convergence}
If a uniformly bounded decreasing net of $\theta$-psh functions converges
pointwise to a bounded $\theta$-psh function, then it converges in
capacity. Uniform convergence implies convergence in capacity as well.
\end{proposition}

\begin{proof}
Let $\varphi_\alpha\searrow\varphi$ be a decreasing net in
$\PSH(X,\theta)$ with $-M\leq\varphi\leq\varphi_\alpha\leq0$ after a
common translation, $M\geq1$, and with $\varphi$ bounded
$\theta$-psh.  Fix $\delta>0$.  For $u\in\mathcal U_\theta$, Chebyshev
and \Cref{lem:fundamental-capacity-estimate} give
\[
 \begin{aligned}
 \MA_\theta(u)\bigl(\{\varphi_\alpha-\varphi>\delta\}\bigr)
 &\leq\frac1\delta\int(\varphi_\alpha-\varphi)\MA_\theta(u)\\
 &\leq\frac{4MV^{1-2^{-n}}}{\delta}\,\Lambda_\alpha^{2^{-n}},\\
 \Lambda_\alpha&=\int(\varphi_\alpha-\varphi)\MA_\theta(\varphi/2).
 \end{aligned}
\]
The weighted decreasing-net continuity in
\Cref{thm:bounded-ma-extension}, applied to the weight
$\varphi_\alpha\searrow\varphi$ and fixed entries $\varphi/2$, gives
$\Lambda_\alpha\to0$.  The bound is uniform in $u$, so
$\Capa_\theta(\{|\varphi_\alpha-\varphi|>\delta\})\to0$; note that
$|\varphi_\alpha-\varphi|=\varphi_\alpha-\varphi$ here.
This proves a Borel-capacity estimate, which must still be upgraded to
outer capacity as required by \Cref{def:capacity-convergence}.

Given $\eta>0$, the construction in the proof of
\Cref{thm:bounded-quasicontinuity} supplies an open $G$ with
$\Capa_\theta(G)<\eta$ and a continuous model majorant $\chi\geq\varphi$
such that $0\leq\chi-\varphi<\delta/3$ on $G^c$: choose a sufficiently
late member of the extracted sequence.  By \Cref{lem:ch24-sandwich},
eventually $\varphi_\alpha\leq\chi+\delta/3$ on all of $X^{\An}$.
Then $\varphi_\alpha-\varphi<2\delta/3$ on $G^c$, so
$\{\varphi_\alpha-\varphi>\delta\}\subseteq G$ and
$\Capa_\theta^*(\{\varphi_\alpha-\varphi>\delta\})<\eta$.
Letting $\eta\downarrow0$ proves convergence in outer capacity.

For the uniform assertion, if $\|f_\alpha-f\|_\infty\to0$ then for each
$\delta>0$ the set $\{|f_\alpha-f|>\delta\}$ is empty for large
$\alpha$, and the empty set has capacity $0$.
\end{proof}

\begin{proposition}
\label{prop:ch25-capacity-convergence-calculus}
Let $f_\alpha\to f$ and $g_\alpha\to g$ in capacity, all functions
bounded by $C$ in absolute value.
\begin{enumerate}
\item\label{itm:ch25-ccc:lin}
 $af_\alpha+bg_\alpha\to af+bg$ in capacity for $a,b\in\mathbf R$, and
 $\max\{f_\alpha,g_\alpha\}\to\max\{f,g\}$,
 $\min\{f_\alpha,g_\alpha\}\to\min\{f,g\}$,
 $f_\alpha g_\alpha\to fg$ in capacity.
\item\label{itm:ch25-ccc:lip}
 If $\chi:\mathbf R\to\mathbf R$ is $L$-Lipschitz, then
 $\chi\circ f_\alpha\to\chi\circ f$ in capacity.
\item\label{itm:ch25-ccc:unique}
 If $f_\alpha\to f$ and $f_\alpha\to f'$ in capacity, then
 $\{f\neq f'\}$ is negligible.
\end{enumerate}
\end{proposition}

\begin{proof}
\cref{itm:ch25-ccc:lin}  For sums, $\{|(af_\alpha+bg_\alpha)-(af+bg)|
>\delta\}\subseteq\{|f_\alpha-f|>\delta/(2|a|+1)\}\cup
\{|g_\alpha-g|>\delta/(2|b|+1)\}$, and $\Capa^*_\theta$ is subadditive by
\Cref{cor:ch25-outer-subadditivity}.  For maxima and minima use
$|\max\{s,t\}-\max\{s',t'\}|\leq\max\{|s-s'|,|t-t'|\}$.  For products,
$|f_\alpha g_\alpha-fg|\leq C(|f_\alpha-f|+|g_\alpha-g|)$.

\cref{itm:ch25-ccc:lip}  If $L=0$ the compositions are the same constant.
For $L>0$, $\{|\chi\circ f_\alpha-\chi\circ f|>\delta\}
\subseteq\{|f_\alpha-f|>\delta/L\}$.

\cref{itm:ch25-ccc:unique}  For each $\delta>0$,
$\{|f-f'|>2\delta\}\subseteq\{|f_\alpha-f|>\delta\}\cup
\{|f_\alpha-f'|>\delta\}$, whose outer capacity tends to $0$; hence
$\Capa^*_\theta(\{|f-f'|>2\delta\})=0$, and
$\{f\neq f'\}=\bigcup_k\{|f-f'|>2/k\}$ is negligible by
\Cref{cor:ch25-outer-subadditivity}.
\end{proof}

\begin{proposition}
\label{prop:ch25-convergence-comparison}
For nets of uniformly bounded functions on $X^{\An}$:
\begin{enumerate}
\item\label{itm:ch25-cc:unif}
 uniform convergence implies convergence in capacity;
\item\label{itm:ch25-cc:seq}
 if a \emph{sequence} converges in capacity, some subsequence converges
 pointwise outside a negligible set;
\item\label{itm:ch25-cc:none}
 convergence in capacity need not be uniform or pointwise for the whole
 sequence: \Cref{ex:ch25-capacity-not-uniform}
 gives a sequence of $\theta$-psh functions converging in capacity but not
 uniformly or pointwise at its pole.  Conversely,
 \Cref{ex:ch25-pointwise-not-capacity} gives a sequence of
 $\theta$-psh functions converging pointwise everywhere and not in
 capacity.
\end{enumerate}
\end{proposition}

\begin{proof}
\cref{itm:ch25-cc:unif} is the last assertion of
\Cref{prop:decreasing-capacity-convergence}.  \cref{itm:ch25-cc:seq} is
\Cref{prop:capacity-subsequence} applied to the countably subadditive
capacity $\Capa^*_\theta$, whose polar sets are by definition the
negligible sets.  \cref{itm:ch25-cc:none} is proved by the two examples
cited.
\end{proof}

\begin{remark}[Nets, sequences and convergence in capacity]
\label{rem:ch25-nets-in-capacity}
\Cref{prop:ch25-convergence-comparison}\,\cref{itm:ch25-cc:seq} is a
statement about sequences and is used only for sequences.  For a net
converging in capacity, no such subnet theorem is proved here.  The proof of
\Cref{prop:capacity-subsequence} builds the exceptional set as a countable
union $\bigcup_{m\geq N}$, and countable subadditivity has nothing to say
about the uncountable unions that a net would produce.  This is the
discipline announced in \Cref{rem:net-vs-sequence-checklist}, and the reason
\Cref{cor:psh-sequential-regularization} is stated as a theorem: it is a
legitimate reduction from nets to sequences, obtained by an estimate, not by
an appeal to countability.
\end{remark}

\begin{proposition}
\label{prop:ch25-truncation-capacity}
Let $\varphi\in\PSH(X,\theta)$ with $\sup\varphi=0$ and put
$\varphi^{(t)}=\max\{\varphi,-t\}$ for $t>0$.  Then
\[
 \Capa_\theta\bigl(\{\varphi^{(t)}\neq\varphi\}\bigr)
 =\Capa_\theta\bigl(\{\varphi<-t\}\bigr)
 \leq\frac{C_\theta+nV}{t}\xrightarrow[t\to\infty]{}0 .
\]
In particular the truncations converge to $\varphi$ in $\theta$-capacity as
$\overline{\mathbf R}$-valued functions.
\end{proposition}

\begin{proof}
The two sets $\{\varphi^{(t)}\neq\varphi\}$ and $\{\varphi<-t\}$
coincide, and the estimate is
\eqref{eq:ch25-sublevel-capacity}.  For every $\delta>0$, the set
$\{\rho(\varphi^{(t)},\varphi)>\delta\}$ is contained in
$\{\varphi<-t\}$, because the two functions agree off the latter.  Its outer
capacity therefore tends to zero.  The metric formulation avoids assigning
a value to an undefined difference at the polar points of $\varphi$.
\end{proof}

\section{Locality and the maximum principle}
\label{sec:ch25-maximum}

The set $\{\varphi>\psi\}$ need not be open: a difference of two upper
semicontinuous functions need not have either semicontinuity.  This is why the maximum principle is not a
formal consequence of weak convergence, and why it has to wait for
quasicontinuity.  The mechanism of the proof is to replace the indicator
$\one_{\{\varphi>\psi\}}$ by a continuous cutoff of the difference of the
approximants, to apply the model identity to that, and then to pass to the
limit in both the measure and the test function---which is exactly what
\Cref{lem:ch25-joint-limit} permits.

\begin{theorem}
\label{thm:bounded-maximum-principle}
If $\varphi,\psi$ are bounded $\theta$-psh functions, then
\begin{equation}
\label{eq:bounded-maximum-principle}
 \one_{\{\varphi>\psi\}}\MA_\theta(\max\{\varphi,\psi\})
 =\one_{\{\varphi>\psi\}}\MA_\theta(\varphi).
\end{equation}
Consequently, if $\varphi=\psi$ on an open set $U$, their
Monge--Amp\`ere measures agree on $U$.
\end{theorem}

\begin{proof}
\emph{Step 0: normalization and the cutoff family.}  Translating both
functions by a common constant, which changes neither side of
\eqref{eq:bounded-maximum-principle}, assume $-M\leq\varphi,\psi\leq0$ with
$M\geq1$.  For $\varepsilon>0$ let
\[
 \chi_\varepsilon:\mathbf R\to[0,1],\qquad
 \chi_\varepsilon(s)=\min\{1,\max\{0,s/\varepsilon\}\} .
\]
Each $\chi_\varepsilon$ is continuous and $\varepsilon^{-1}$-Lipschitz,
vanishes on $(-\infty,0]$, equals $1$ on $[\varepsilon,\infty)$, and
increases to $\one_{(0,\infty)}$ pointwise as
$\varepsilon\downarrow0$.

\emph{Step 1: the model identity, multiplied by a cutoff.}  By
\Cref{thm:bfj-decreasing-regularization} choose decreasing nets of
$\theta$-psh model functions $\varphi_\alpha\searrow\varphi$ and
$\psi_\alpha\searrow\psi$ indexed by a common directed set, obtained by
taking the product of the two index sets, and normalized as in Step~1 of the
proof of \Cref{thm:bounded-quasicontinuity} so that
$-M\leq\varphi\leq\varphi_\alpha\leq0$ and
$-M\leq\psi\leq\psi_\alpha\leq0$.  The functions
$\max\{\varphi_\alpha,\psi_\alpha\}$ are continuous $\theta$-psh
functions by \Cref{thm:ch21-stability} and decrease to
$\max\{\varphi,\psi\}$.  A maximum of real model functions need not
be a model function; \Cref{prop:model-maximum-principle} already includes
this continuous case.

Fix $\varepsilon>0$ and $f\in C^0(X^{\An})$.  The continuous function
\[
 h_\alpha:=f\cdot\chi_\varepsilon(\varphi_\alpha-\psi_\alpha)
\]
vanishes off $\{\varphi_\alpha>\psi_\alpha\}$.  On this strict set,
\Cref{prop:model-maximum-principle} gives
\eqref{eq:model-maximum-principle}.  Hence
\begin{equation}
\label{eq:ch25-max-model-tested}
 \int h_\alpha\,\MA_\theta(\max\{\varphi_\alpha,\psi_\alpha\})
 =\int h_\alpha\,\MA_\theta(\varphi_\alpha) .
\end{equation}

\emph{Step 2: passing to the limit in $\alpha$.}  By
\Cref{prop:decreasing-capacity-convergence} the nets
$(\varphi_\alpha)$ and $(\psi_\alpha)$ converge in capacity to
$\varphi$ and $\psi$; by
\Cref{prop:ch25-capacity-convergence-calculus}\,\cref{itm:ch25-ccc:lin} and
\cref{itm:ch25-ccc:lip} the functions
$h_\alpha$ converge in capacity to
\[
 h:=f\cdot\chi_\varepsilon(\varphi-\psi),
\]
uniformly bounded by $\|f\|_\infty$.  Each $h_\alpha$ is continuous,
hence quasicontinuous, and $h$ is quasicontinuous by
\Cref{cor:ch25-quasicontinuity-calculus}, being a product of a continuous
function with a Lipschitz function of a difference of two bounded
$\theta$-psh functions.  By \Cref{thm:bounded-ma-extension} the measures
$\MA_\theta(\max\{\varphi_\alpha,\psi_\alpha\})$ and
$\MA_\theta(\varphi_\alpha)$ converge weakly to
$\MA_\theta(\max\{\varphi,\psi\})$ and $\MA_\theta(\varphi)$; all
potentials involved lie in $[-M,0]$.  \Cref{lem:ch25-joint-limit} therefore
applies to both sides of \eqref{eq:ch25-max-model-tested} and yields
\begin{equation}
\label{eq:ch25-max-limit-tested}
 \int f\,\chi_\varepsilon(\varphi-\psi)\,
   \MA_\theta(\max\{\varphi,\psi\})
 =\int f\,\chi_\varepsilon(\varphi-\psi)\,\MA_\theta(\varphi)
\end{equation}
for every $f\in C^0(X^{\An})$.

\emph{Step 3: from test functions to measures.}  Both sides of
\eqref{eq:ch25-max-limit-tested} are integrals of $f$ against the positive
Radon measures $\chi_\varepsilon(\varphi-\psi)\MA_\theta(\max\{\varphi,
\psi\})$ and $\chi_\varepsilon(\varphi-\psi)\MA_\theta(\varphi)$, the
density being a bounded Borel function.  Since they agree against every
continuous $f$, the two measures are equal by
\Cref{cor:riesz-uniqueness}:
\begin{equation}
\label{eq:ch25-max-measure-eps}
 \chi_\varepsilon(\varphi-\psi)\,\MA_\theta(\max\{\varphi,\psi\})
 =\chi_\varepsilon(\varphi-\psi)\,\MA_\theta(\varphi) .
\end{equation}

\emph{Step 4: letting $\varepsilon\downarrow0$.}  The function
$\varphi-\psi$ is a real-valued Borel function, both potentials being
bounded and upper semicontinuous, so $\{\varphi>\psi\}$ is Borel.  As
$\varepsilon\downarrow0$ the densities
$\chi_\varepsilon(\varphi-\psi)$ increase pointwise to
$\one_{\{\varphi>\psi\}}$.  Applying the monotone convergence theorem
(\Cref{lem:integration-ledger}) to each of the two fixed positive measures in
\eqref{eq:ch25-max-measure-eps}, evaluated on an arbitrary Borel set, gives
\eqref{eq:bounded-maximum-principle}.

\emph{Step 5: locality.}  Suppose $\varphi=\psi$ on an open set $U$.  For
$c>0$ we have $U\subseteq\{\varphi+c>\psi\}$, so
\eqref{eq:bounded-maximum-principle}, applied to the pair
$(\varphi+c,\psi)$ and combined with
$\MA_\theta(\varphi+c)=\MA_\theta(\varphi)$, gives
\[
 \one_{U}\,\MA_\theta(\max\{\varphi+c,\psi\})
 =\one_{U}\,\MA_\theta(\varphi) ,
\]
and, exchanging the two functions,
$\one_U\MA_\theta(\max\{\psi+c,\varphi\})=\one_U\MA_\theta(\psi)$.  As
$c\downarrow0$ the functions $\max\{\varphi+c,\psi\}$ decrease to
$\max\{\varphi,\psi\}$, uniformly, since
$0\leq\max\{\varphi+c,\psi\}-\max\{\varphi,\psi\}\leq c$; the same holds
for $\max\{\psi+c,\varphi\}$.  By \Cref{thm:bounded-ma-extension} both
families of measures converge weakly to
$\MA_\theta(\max\{\varphi,\psi\})$.  Let $f\in C^0(X^{\An})$ have support
contained in $U$.  Then $\int f\MA_\theta(\varphi)
=\int f\MA_\theta(\max\{\varphi+c,\psi\})$ for every $c>0$, and letting
$c\downarrow0$ gives $\int f\MA_\theta(\varphi)
=\int f\MA_\theta(\max\{\varphi,\psi\})$.  The same computation with the
roles exchanged gives $\int f\MA_\theta(\psi)
=\int f\MA_\theta(\max\{\varphi,\psi\})$.  Since the two Radon measures
$\MA_\theta(\varphi)$ and $\MA_\theta(\psi)$ agree against every
continuous function supported in $U$, they agree on $U$ by
\Cref{cor:riesz-uniqueness}.
\end{proof}

\begin{corollary}
\label{cor:ch25-locality}
Let $U\subseteq X^{\An}$ be open.  If $\varphi,\psi$ are bounded
$\theta$-psh and $\varphi=\psi$ on $U$, then
$\MA_\theta(\varphi)|_U=\MA_\theta(\psi)|_U$.  More generally, if
$\varphi_1,\ldots,\varphi_n$ and $\psi_1,\ldots,\psi_n$ are bounded
$\theta$-psh with $\varphi_i=\psi_i$ on $U$ for every $i$, then
\[
 \MA_\theta(\varphi_1,\ldots,\varphi_n)|_U
 =\MA_\theta(\psi_1,\ldots,\psi_n)|_U .
\]
\end{corollary}

\begin{proof}
The first assertion is the last clause of
\Cref{thm:bounded-maximum-principle}.  For the mixed statement use the
bounded polarization identity in \Cref{cor:bounded-ma-multilinearity}.
For a nonempty $S\subseteq\{1,\ldots,n\}$ put
$\varphi_S:=(\#S)^{-1}\sum_{i\in S}\varphi_i$, a convex combination of
$\theta$-psh functions, hence $\theta$-psh by
\Cref{lem:ch25-admissible-calculus}\,\cref{itm:ch25-adm:nonempty} and bounded;
the bounded polarization identity gives
\[
 n!\,\MA_\theta(\varphi_1,\ldots,\varphi_n)
 =\sum_{\varnothing\neq S}(-1)^{\,n-\#S}(\#S)^n\,
    \MA_\theta(\varphi_S).
\]
Each $\varphi_S$ agrees on $U$ with the corresponding $\psi_S$, so the
diagonal case applies term by term and the two sides agree on $U$.
\end{proof}

\begin{remark}[Two sets that must not be confused]
\label{rem:ch25-max-sets}
Identity \eqref{eq:bounded-maximum-principle} concerns the \emph{strict}
inequality locus.  No general equality with either operand's measure is
asserted on the contact set $\{\varphi=\psi\}$: the maximum's measure may
differ from both $\MA_\theta(\varphi)$ and $\MA_\theta(\psi)$ there,
because the maximum acquires new curvature along the contact locus.
\Cref{ex:ch25-comparison-numbers} exhibits this on a metric graph with
explicit numbers: two model potentials whose maximum carries an extra atom
of mass $\tfrac12$ at the contact point, where both original measures
vanish.  This is why the proof of the comparison
principle below is arranged to avoid the contact set entirely, by shifting
one potential by $\varepsilon$ and letting $\varepsilon\downarrow0$ at
the end.
\end{remark}

\section{Comparison and domination}
\label{sec:ch25-comparison}

\begin{corollary}
\label{cor:bounded-comparison-principle}
For bounded $\theta$-psh functions $\varphi,\psi$,
\begin{equation}
\label{eq:bounded-comparison-principle}
       \int_{\{\varphi<\psi\}}\MA_\theta(\psi)
       \leq
       \int_{\{\varphi<\psi\}}\MA_\theta(\varphi).
\end{equation}
\end{corollary}

\begin{proof}
Fix $\varepsilon>0$ and write $m_\varepsilon:=\max\{\varphi,
\psi-\varepsilon\}$, a bounded $\theta$-psh function by
\Cref{thm:ch21-stability}\,\cref{itm:ch21-stab:max}.  Two applications of
\Cref{thm:bounded-maximum-principle} are used, on the two disjoint sets
\[
 P:=\{\varphi<\psi-\varepsilon\},\qquad
 Q:=\{\varphi>\psi-\varepsilon\} .
\]
Applying \eqref{eq:bounded-maximum-principle} to the pair
$(\psi-\varepsilon,\varphi)$ gives
$\one_P\MA_\theta(m_\varepsilon)=\one_P\MA_\theta(\psi-\varepsilon)
=\one_P\MA_\theta(\psi)$; applying it to the pair
$(\varphi,\psi-\varepsilon)$ gives
$\one_Q\MA_\theta(m_\varepsilon)=\one_Q\MA_\theta(\varphi)$.  In
particular
\begin{equation}
\label{eq:ch25-comparison-step}
 \begin{aligned}
 \int_P\MA_\theta(\psi)&=\int_P\MA_\theta(m_\varepsilon)\\
 &=V-\int_{X^{\An}\setminus P}\MA_\theta(m_\varepsilon)\\
 &\leq V-\int_Q\MA_\theta(m_\varepsilon)
 =V-\int_Q\MA_\theta(\varphi).
 \end{aligned}
\end{equation}
using in the third step that $Q\subseteq X^{\An}\setminus P$ and that
$\MA_\theta(m_\varepsilon)$ is positive of total mass $V$
(\Cref{thm:bounded-ma-extension}).  Since $\MA_\theta(\varphi)$ also has
total mass $V$, the right-hand side equals
$\int_{X^{\An}\setminus Q}\MA_\theta(\varphi)$, that is
\[
 \int_{\{\varphi<\psi-\varepsilon\}}\MA_\theta(\psi)
 \leq\int_{\{\varphi\leq\psi-\varepsilon\}}\MA_\theta(\varphi).
\]
Now let $\varepsilon\downarrow0$ along a sequence.  The sets
$\{\varphi<\psi-\varepsilon\}$ increase to $\{\varphi<\psi\}$, and so do
the sets $\{\varphi\leq\psi-\varepsilon\}$: indeed each is contained in
$\{\varphi<\psi\}$, each increases as $\varepsilon$ decreases, and every
point with $\varphi<\psi$ satisfies $\varphi\leq\psi-\varepsilon$ for
some $\varepsilon>0$, both functions being real valued.  Countable
additivity of each of the two fixed measures along an increasing sequence of
Borel sets gives \eqref{eq:bounded-comparison-principle}.
\end{proof}

\begin{remark}[Why the shift by \texorpdfstring{$\varepsilon$}{epsilon} is
needed, and what it costs]
\label{rem:ch25-comparison-shift}
Running the argument with $\varepsilon=0$ would require comparing
$\MA_\theta(\max\{\varphi,\psi\})$ with the two measures on the contact set
$\{\varphi=\psi\}$, where \Cref{thm:bounded-maximum-principle} says
nothing; see \Cref{rem:ch25-max-sets}.  The shift removes the contact set at
the price of replacing $\{\varphi<\psi\}$ by a slightly smaller set, and
the limit recovers it.  The same device produces the two variants
\[
 \begin{aligned}
 \int_{\{\varphi<\psi\}}\MA_\theta(\psi)
 &\leq\int_{\{\varphi\leq\psi\}}\MA_\theta(\varphi),\\
 \int_{\{\varphi<\psi\}}\MA_\theta(\psi)
 &\leq\int_{\{\varphi<\psi\}}\MA_\theta(\varphi)
 \leq\int_{\{\varphi\leq\psi\}}\MA_\theta(\varphi).
 \end{aligned}
\]
of which only the sharpest,
\eqref{eq:bounded-comparison-principle}, is recorded above.  No inequality
with $\{\varphi\leq\psi\}$ on the \emph{left} is claimed: the mass that
$\MA_\theta(\psi)$ puts on the contact set is genuinely uncontrolled, as
\Cref{ex:ch25-comparison-numbers} shows.
\end{remark}

\begin{corollary}
\label{cor:ch25-sublevel-comparison}
Let $\varphi$ be bounded $\theta$-psh with $\varphi\leq0$, and let
$s>0$, $0<t<1$.  Then
\begin{equation}
\label{eq:ch25-sublevel-comparison}
 \Capa_\theta\bigl(\{\varphi<-s\}\bigr)
 \leq t^{-n}\,\MA_\theta(\varphi)
   \bigl(\{\varphi<-s+t(1+s)\}\bigr).
\end{equation}
\end{corollary}

\begin{proof}
Apply \eqref{eq:capacity-comparison}, proved as
\Cref{prop:bounded-capacity-domination} below --- whose proof does not use
the present corollary, so there is no circularity --- with $\psi$ the
constant $-s$, which is $\theta$-psh by
\Cref{prop:ch21-constants}\,\cref{itm:ch21-const:const} and satisfies
$\psi\leq0$.  The right-hand side becomes
$t^{-n}\MA_\theta(\varphi)(\{\varphi<-(1-t)s+t\})$, and
$-(1-t)s+t=-s+t(1+s)$.
\end{proof}

This is the form in which the comparison principle enters the energy theory:
\Cref{cor:e1-capacity-tail} is exactly
\eqref{eq:ch25-sublevel-comparison} applied to a truncation, with
$t=1/2$.  \Cref{ex:ch25-comparison-numbers} evaluates both sides of
\eqref{eq:ch25-sublevel-comparison} on a metric graph and measures the loss.

\begin{proposition}
\label{prop:bounded-capacity-domination}
Let $\varphi,\psi$ be bounded $\theta$-psh and suppose $\psi\leq0$.
For $0<t<1$,
\begin{equation}
\label{eq:capacity-comparison}
 \Capa_\theta\{\varphi<\psi\}
 \leq t^{-n}
 \int_{\{\varphi<(1-t)\psi+t\}}\MA_\theta(\varphi).
\end{equation}
If, in addition, $\varphi\leq\psi$ almost everywhere for $\MA_\theta(\psi)$, then
$\varphi\leq\psi$ everywhere.
\end{proposition}

\begin{proof}
\emph{The capacity estimate.}  Fix $u\in\PSH(X,\theta)$ with
$0\leq u\leq1$.  The translation clause in
\Cref{lem:ch25-admissible-calculus}\,\cref{itm:ch25-adm:transl} shows that
these functions have the same Monge--Amp\`ere measures as the members of
$\mathcal U_\theta$.  Thus
\begin{equation}
\label{eq:ch25-dom-sup}
 \Capa_\theta(E)=\sup\bigl\{\MA_\theta(u)(E)\ :\
   u\in\PSH(X,\theta),\ 0\leq u\leq1\bigr\}
\end{equation}
for every Borel $E$.  Fix such a $u$ and put
$\psi_t:=(1-t)\psi+tu$, a bounded $\theta$-psh function by convexity
(\Cref{lem:ch25-admissible-calculus}\,\cref{itm:ch25-adm:nonempty}).  Two
inclusions and one inequality are needed.

First, $\psi_t\geq\psi$, because
$\psi_t-\psi=t(u-\psi)\geq0$: here $u\geq0\geq\psi$, and this is the
only place where the hypothesis $\psi\leq0$ is used.  Hence
$\{\varphi<\psi\}\subseteq\{\varphi<\psi_t\}$.  Second, $u\leq1$ gives
$\psi_t\leq(1-t)\psi+t$, so
$\{\varphi<\psi_t\}\subseteq\{\varphi<(1-t)\psi+t\}$.  Third, expanding the
$n$-fold product of
$\theta+\ddc\psi_t=(1-t)(\theta+\ddc\psi)+t(\theta+\ddc u)$ by
multilinearity and discarding all but one nonnegative term,
\begin{equation}
\label{eq:ch25-dom-lower}
 \MA_\theta(\psi_t)\geq t^n\MA_\theta(u),
\end{equation}
which is the content of \eqref{eq:ma-convex-lower-bound}.  Now apply
\Cref{cor:bounded-comparison-principle} to the pair $(\varphi,\psi_t)$:
\[
 \begin{aligned}
 t^n\MA_\theta(u)\bigl(\{\varphi<\psi\}\bigr)
 &\leq t^n\MA_\theta(u)\bigl(\{\varphi<\psi_t\}\bigr)\\
 &\leq\int_{\{\varphi<\psi_t\}}\MA_\theta(\psi_t)
 \leq\int_{\{\varphi<\psi_t\}}\MA_\theta(\varphi)\\
 &\leq\int_{\{\varphi<(1-t)\psi+t\}}\MA_\theta(\varphi).
 \end{aligned}
\]
Dividing by $t^n$ and taking the supremum over $u$ as in
\eqref{eq:ch25-dom-sup} gives \eqref{eq:capacity-comparison}.

\emph{The domination principle.}  Assume in addition that
$\MA_\theta(\psi)(\{\varphi>\psi\})=0$.  Let $C\geq0$
be such that $\varphi\geq-C$, and set
$c:=\max\{0,\sup_{X^{\An}}\varphi\}$, which is finite.  Fix
$\varepsilon>0$.  We apply \eqref{eq:capacity-comparison} to the pair
\[
 \varphi':=\psi+\varepsilon-c,\qquad
 \psi':=\varphi-c ,
\]
both bounded $\theta$-psh
(\Cref{prop:ch21-constants}\,\cref{itm:ch21-const:transl}), with
$\psi'\leq0$ by the choice of $c$.  Note that
$\{\varphi'<\psi'\}=\{\psi+\varepsilon<\varphi\}$, a set unchanged by the
common translation; this is the point at which translating only one of the
two functions would destroy the hypothesis, and it is why both are shifted
by the same constant.  For $0<t<1$,
\[
 \begin{aligned}
 \Capa_\theta\bigl(\{\psi+\varepsilon<\varphi\}\bigr)
 &\leq t^{-n}
   \int_{\{\varphi'<(1-t)\psi'+t\}}\MA_\theta(\varphi')\\
 &=t^{-n}\,\MA_\theta(\psi)
   \bigl(\{\varphi'<(1-t)\psi'+t\}\bigr),
 \end{aligned}
\]
since $\MA_\theta(\varphi')=\MA_\theta(\psi)$.  Unwinding the translations,
\[
 \{\varphi'<(1-t)\psi'+t\}
 =\bigl\{\psi+\varepsilon<\varphi+t(c+1-\varphi)\bigr\}
 \subseteq\bigl\{\psi+\varepsilon<\varphi+t(c+1+C)\bigr\},
\]
using $-\varphi\leq C$.  Choose $t\in(0,1)$ so small that
$t(c+1+C)\leq\varepsilon/2$; then the set on the right is contained in
$\{\psi+\varepsilon/2<\varphi\}\subseteq\{\psi<\varphi\}$, which is
$\MA_\theta(\psi)$-null by hypothesis.  Therefore
\[
 \Capa_\theta\bigl(\{\psi+\varepsilon<\varphi\}\bigr)=0 .
\]
By \Cref{lem:divisorial-positive-capacity}, a Borel set of capacity zero
contains no divisorial point.  Consequently
$\varphi(x)\leq\psi(x)+\varepsilon$ at every divisorial point $x$.
Letting $\varepsilon\downarrow0$ gives $\varphi\leq\psi$ there, and
\Cref{prop:ch21-divisorial-order} gives the inequality on all of $X^{\An}$.
This final step requires neither openness of the strict set nor an
identification of Borel capacity with outer capacity.
\end{proof}

\begin{warning}{Normalization and boundedness in domination}
No continuity hypothesis on $\varphi$ is needed: divisorial order upgrades
the capacity-zero conclusion to global order.  The
hypothesis $\psi\leq0$ in \eqref{eq:capacity-comparison} is used once, to
guarantee $\psi_t\geq\psi$; in the application it is arranged by
translating \emph{both} potentials by the same constant, which leaves the
strict-inequality set and the almost-everywhere hypothesis intact.
Translating only $\varphi$, so as to normalize $0\geq\varphi\geq-C$,
would replace the hypothesis ``$\varphi\leq\psi$ almost everywhere'' by a
strictly weaker one and would not give the stated conclusion.  Finally,
boundedness of both potentials is essential: $\MA_\theta(\psi)$ is defined
only for bounded $\psi$ at this stage, and the bound $C$ enters the
choice of $t$.
\end{warning}

\begin{corollary}
\label{cor:ch25-domination-uniqueness}
Let $\varphi,\psi$ be bounded $\theta$-psh.
If $\MA_\theta(\psi)$ gives no mass to $\{\varphi>\psi\}$, then
$\varphi\leq\psi$ on all of $X^{\An}$.  If
\[
 \MA_\theta(\psi)\bigl(\{\varphi\neq\psi\}\bigr)
 =\MA_\theta(\varphi)\bigl(\{\varphi\neq\psi\}\bigr)=0 ,
\]
that is, if each of the two measures is carried by the contact set, then
$\varphi=\psi$.
\end{corollary}

\begin{proof}
After a common translation arranging $\psi\leq0$, the first assertion
is the second half of \Cref{prop:bounded-capacity-domination}.  For the second, apply that half
twice.  Since $\{\varphi>\psi\}\subseteq\{\varphi\neq\psi\}$, the
hypothesis on $\MA_\theta(\psi)$ gives $\varphi\leq\psi$; since
$\{\psi>\varphi\}\subseteq\{\varphi\neq\psi\}$, the hypothesis on
$\MA_\theta(\varphi)$ gives $\psi\leq\varphi$, the roles of the two
functions being exchanged and normalized by a common translation.
\end{proof}

\begin{remark}[Why both carriers are assumed]
\label{rem:ch25-both-carriers}
The second hypothesis in
\Cref{cor:ch25-domination-uniqueness} is not a consequence of the first.
For example, on the projective line of
\Cref{ex:ch25-pluripolar-computed}, take $\psi=0$ and
$\varphi=\max\{\varphi_0,-1\}$.  The measure
$\MA_\theta(\psi)=\delta_\zeta$ is carried by the contact set, whereas
$\MA_\theta(\varphi)=\delta_{x_1}$ is carried by the strict lower set.
Thus one carrier assumption alone cannot force equality.  More generally, once $\varphi\leq\psi$ is known, $\{\varphi\neq\psi\}=\{\varphi<\psi\}$,
and the comparison principle \eqref{eq:bounded-comparison-principle} gives
$\int_{\{\varphi<\psi\}}\MA_\theta(\psi)\leq
\int_{\{\varphi<\psi\}}\MA_\theta(\varphi)$; the left-hand side vanishes,
but that says nothing about the right-hand side, which is the quantity one
would need.  Thus this comparison inequality alone does not give the reversed
capacity-null conclusion.  Equality of Monge--Amp\`ere measures requires
a separate uniqueness argument, as \Cref{rem:ch25-no-uniqueness} records.
\end{remark}

\begin{remark}[What is not deduced: uniqueness for the Monge--Amp\`ere
equation]
\label{rem:ch25-no-uniqueness}
It is tempting to deduce from
\Cref{cor:ch25-domination-uniqueness} that two bounded $\theta$-psh
functions with the same Monge--Amp\`ere measure differ by a constant.  That
deduction is not valid, and we do not make it.  The hypothesis of the
domination principle is a statement about the \emph{null sets} of
$\MA_\theta(\psi)$, and the equality
$\MA_\theta(\varphi)=\MA_\theta(\psi)$ gives no information about
$\MA_\theta(\psi)(\{\varphi>\psi\})$.  Uniqueness for the
Monge--Amp\`ere equation is proved in this book by the variational
method, using concavity of the energy along segments,
\Cref{cor:ch20-concavity}, differentiability, and the separate
nondegeneracy input in \Cref{thm:ch29-uniqueness}.  The direct argument
above does not supply that uniqueness proof.  Concavity alone need not
be strict for the merely nef model classes that occur along a segment.
Strictness and nondegeneracy must therefore be checked separately.
\end{remark}

\section{Unbounded potentials and finite energy}
\label{sec:ch25-e1}

The finite-energy class $\cE^1$ is constructed in
\Cref{chap:bfj-energy}, after the present chapter.  Its members are
generally unbounded, so the bounded theorem does not apply to them directly.
The following truncation lemma isolates the tail property that passes from
bounded real-valued quasicontinuity to extended-real-valued
quasicontinuity.  In the present BFJ setting that property follows from the
universal sublevel estimate already proved in this chapter; finite energy
will later improve its rate, not its validity.

\begin{lemma}
\label{lem:ch25-quasicontinuity-from-tail}
Let $\varphi\in\PSH(X,\theta)$ and suppose
\begin{equation}
\label{eq:ch25-tail-hypothesis}
 \Capa_\theta\bigl(\{\varphi<-T\}\bigr)\longrightarrow0
 \qquad\text{as }T\to+\infty .
\end{equation}
Then for every $\varepsilon>0$ there is an open set $G$ with
$\Capa_\theta(G)<\varepsilon$ such that $\varphi$ is finite and
continuous on $X^{\An}\setminus G$.  In particular $\varphi$ is
quasicontinuous as a map to $\mathbf R\cup\{-\infty\}$.
\end{lemma}

\begin{proof}
Normalize $\sup\varphi\leq0$.  Let $\varepsilon>0$.  By
\eqref{eq:ch25-tail-hypothesis} choose $T>0$ with
$\Capa_\theta(\{\varphi<-T\})<\varepsilon/2$; the set
$G_1:=\{\varphi<-T\}$ is open, $\varphi$ being upper semicontinuous.  The
truncation $\varphi^{(T)}:=\max\{\varphi,-T\}$ is bounded $\theta$-psh
(\Cref{thm:ch21-stability}\,\cref{itm:ch21-stab:max}), so by
\Cref{thm:bounded-quasicontinuity} there is an open $G_2$ with
$\Capa_\theta(G_2)<\varepsilon/2$ such that
$\varphi^{(T)}|_{X^{\An}\setminus G_2}$ is continuous.  Put
$G:=G_1\cup G_2$, open with $\Capa_\theta(G)<\varepsilon$ by
\eqref{eq:capacity-subadditivity}.  Off $G$ we have $\varphi\geq-T$,
hence $\varphi=\varphi^{(T)}$; so $\varphi$ is finite there, and its
restriction to $X^{\An}\setminus G\subseteq X^{\An}\setminus G_2$ is
continuous, a restriction of a continuous function to a subspace being
continuous.
\end{proof}

\begin{remark}[Why the tail hypothesis is displayed]
\label{rem:ch25-tail-not-automatic}
Condition \eqref{eq:ch25-tail-hypothesis} holds for every \emph{bounded}
$\theta$-psh function trivially, and for every $\theta$-psh function
whatsoever by \Cref{prop:ch25-sublevel-capacity}, once one normalizes
$\sup\varphi=0$.  So in the present setting
\Cref{lem:ch25-quasicontinuity-from-tail} applies to all of
$\PSH(X,\theta)$.  We nevertheless state it with the hypothesis displayed,
for two reasons.  First, the rate matters: the size of the exceptional set
is governed by the rate in \eqref{eq:ch25-tail-hypothesis}, and
\Cref{cor:e1-capacity-tail} supplies a better rate for $\cE^1$ than the
universal $O(1/T)$ of \eqref{eq:ch25-sublevel-capacity}.  Second, in the
trivially valued theory of \Cref{part:bj} the analogue of
\Cref{prop:ch25-sublevel-capacity} is not available in the same form, and
the tail hypothesis is then a genuine restriction; the lemma is stated so
that it can be quoted there without change.
\end{remark}

\begin{theorem}
\label{thm:ch25-psh-quasicontinuity}
\label{thm:e1-quasicontinuity}
Every $\theta$-psh function is quasicontinuous as a map to
$\mathbf R\cup\{-\infty\}$.  More precisely, outside an open set of
arbitrarily small capacity it is finite and continuous.  In particular this
holds for every finite-energy potential introduced in
\Cref{chap:bfj-energy}.
\end{theorem}

\begin{proof}
Let $\varphi\in\PSH(X,\theta)$ and normalize $\sup\varphi=0$.  By
\Cref{prop:ch25-sublevel-capacity},
\[
 \Capa_\theta(\{\varphi<-T\})
 \leq\frac{C_\theta+nV}{T}\longrightarrow0.
\]
Thus \eqref{eq:ch25-tail-hypothesis} holds, and
\Cref{lem:ch25-quasicontinuity-from-tail} applies.  The finite-energy clause
is immediate because $\cE^1(X,\theta)\subseteq\PSH(X,\theta)$.
\end{proof}

\begin{remark}[The sharper finite-energy tail belongs to Chapter 26]
\label{rem:ch25-no-circularity}
No result of this chapter now depends on the later energy theory.
\Cref{cor:e1-capacity-tail} combines the capacity comparison
\eqref{eq:capacity-comparison} proved here with the energy integrability
estimate \eqref{eq:e1-integrability}; it gives a sharper tail bound for
$\cE^1$ than the universal $O(T^{-1})$ estimate above.  That refinement is
used in the quantitative energy arguments, but it is not needed to prove
quasicontinuity.
\end{remark}

\begin{warning}{The bounded estimates are not asserted for
\texorpdfstring{$\cE^1$}{E1}}
\Cref{thm:bounded-quasicontinuity} does not apply to an unbounded potential,
and the reason is not a technicality.  Its proof uses
\Cref{lem:fundamental-capacity-estimate}, whose constant is
$4MV^{1-2^{-n}}$ with $M$ a uniform bound for the potentials; for a
finite-energy potential no such $M$ exists, and applying the estimate to
the truncations $\max\{\varphi,-T\}$ degrades the constant linearly in
$T$.  Extended-real-valued quasicontinuity instead follows by combining
bounded quasicontinuity of one truncation with the tail estimate
\eqref{eq:ch25-tail-hypothesis}; the universal bound
\eqref{eq:ch25-sublevel-capacity} already supplies that estimate, while
finite energy supplies a sharper one later.  The constants in \Cref{lem:ma-capacity-bound} and
\Cref{lem:fundamental-capacity-estimate} are not applied verbatim with a
finite $M$ to an unbounded potential.  Finite-energy tail and convergence
arguments require their own truncation estimates.
\end{warning}

\section{Capacity on a curve}
\label{sec:ch25-curves}

On a curve, the capacity of a tube over a fixed skeleton descends to
optimization on a finite metric graph.  The examples below admit exact
maximizers; arbitrary Borel subsets need not have a finite-dimensional
optimization formula.  The mechanism is a pushforward identity
which has no counterpart in higher dimension and which rests on the curve
retraction theorem \Cref{prop:ch21-curve-retraction}.

Throughout this section $C$ is a smooth projective connected $K$-curve,
$\theta$ is a closed model form determined on a semistable model
$\mathcal C$, $\Gamma=\Delta_{\mathcal C}$ is the dual graph with its
canonical metric, $p=p_{\mathcal C}:C^{\An}\to\Gamma$ the retraction, and
$\mu_0$ the curvature measure of $\theta$ on $\Gamma$ in the sense of
\eqref{eq:graph-curvature}, a positive measure of total mass $V=\deg\{\theta\}$
supported on the vertices.  Recall from \eqref{eq:ch21-curve-psh-test} that a
model function $\psi$ determined on a regular SNC model
$\mathcal C'\geq\mathcal C$, with skeleton
$\Gamma'\supseteq\Gamma$, is $\theta$-psh precisely when
$\mu_{\theta,\mathcal C'}+
\Delta^{\mathrm w}_{\Gamma'}\psi\geq0$ on $\Gamma'$.  Moreover,
\Cref{prop:ch20-curve-laplacian} and
\eqref{eq:ch20-curve-laplacian-general} give
\begin{equation}
\label{eq:ch25-curve-ma}
 \MA_\theta(\psi)
 =\mu_{\theta,\mathcal C'}+\Delta^{\mathrm w}_{\Gamma'}\psi
\end{equation}
as measures on $\Gamma'\subseteq C^{\An}$.  Here the reference measure is
\eqref{eq:ch21-curve-reference-measure} and the weighted Laplacian is
\eqref{eq:ch21-weighted-graph-laplacian}; on $\Gamma$ we write the
reference measure simply as $\mu_0$.
For a continuous function on the finite weighted graph, the same notation
$\Delta^{\mathrm w}_\Gamma$ denotes the distributional weighted
Laplacian; positivity means that the resulting distribution is a positive
Radon measure.

\begin{proposition}
\label{prop:ch25-curve-pushforward}
For every bounded $\theta$-psh function $u$ on $C^{\An}$,
\begin{equation}
\label{eq:ch25-curve-pushforward}
 p_*\MA_\theta(u)=\MA_\theta(u\circ p).
\end{equation}
Equivalently, $\MA_\theta(u)(p^{-1}(A))=\MA_\theta(u\circ p)(A)$ for every
Borel $A\subseteq\Gamma$.
\end{proposition}

\begin{proof}
\emph{Step 1: model functions.}  Let $\psi$ be a $\theta$-psh model
function.  Choose a regular SNC model $\mathcal C'\geq\mathcal C$
determining it, and write $\Gamma'=\Delta_{\mathcal C'}$.  By
\Cref{prop:ch21-curve-retraction}, $\psi\circ p$ is again a $\theta$-psh
model function, determined on a model with skeleton $\Gamma$, and
\eqref{eq:ch25-curve-ma} computes both measures.  Both are supported on
$\Gamma'$, so it suffices to compute the mass of $\MA_\theta(\psi)$ on
$p^{-1}(\{q\})$ for $q\in\Gamma$.

Write $\Gamma'\setminus\Gamma$ as a disjoint union of connected components;
by the argument in Step~2 of the proof of
\Cref{prop:ch21-curve-retraction}, the closure $T$ of such a component
meets $\Gamma$ in a single point $q(T)$.  Since
$p|_{\Gamma'}$ is the retraction of the graph $\Gamma'$ onto the
subgraph $\Gamma$, which collapses $T\setminus\{q(T)\}$ to $q(T)$, we
have
\[
 p^{-1}(\{q\})\cap\Gamma'
 =\{q\}\cup\bigcup_{T:\,q(T)=q}\bigl(T\setminus\{q\}\bigr).
\]
Now $\mu_0(\Gamma'\setminus\Gamma)=0$ by
\eqref{eq:ch21-curve-no-new-mass}, and summing
\eqref{eq:ch21-weighted-graph-laplacian} over the points of
$T\setminus\{q\}$ gives, as in \eqref{eq:ch21-curve-flux},
\[
 \Delta^{\mathrm w}_{\Gamma'}\psi\bigl(T\setminus\{q\}\bigr)
 =-\sum_{e\subseteq T,\ e\ni q}m_e\partial_e\psi(q),
\]
all interior edges contributing twice with opposite signs.  On the other
hand, at $q$ itself,
\[
 \Delta^{\mathrm w}_{\Gamma'}\psi(\{q\})
 =\sum_{e\subseteq\Gamma,\ e\ni q}m_e\partial_e\psi(q)
  +\sum_{T:\,q(T)=q}\ \sum_{e\subseteq T,\ e\ni q}
       m_e\partial_e\psi(q).
\]
Adding the two displays, the branch terms cancel and
\[
 \begin{aligned}
 \MA_\theta(\psi)\bigl(p^{-1}(\{q\})\bigr)
 &=\mu_0(\{q\})+
   \sum_{e\subseteq\Gamma,\ e\ni q}m_e\partial_e\psi(q)\\
 &=\bigl(\mu_0+\Delta^{\mathrm w}_\Gamma
                 (\psi|_\Gamma)\bigr)(\{q\})\\
 &=\MA_\theta(\psi\circ p)(\{q\}),
 \end{aligned}
\]
the last step because $\psi\circ p$ is constant on each branch and agrees
with $\psi$ on $\Gamma$.  Both measures being atomic and supported on
$\Gamma'$, respectively $\Gamma$, this proves
\eqref{eq:ch25-curve-pushforward} for $\psi$.

\emph{Step 2: bounded $\theta$-psh functions.}  Let $u$ be bounded
$\theta$-psh and choose a decreasing net $\psi_\alpha\searrow u$ of
$\theta$-psh model functions
(\Cref{thm:bfj-decreasing-regularization}).  Then
$\psi_\alpha\circ p\searrow u\circ p$, each $\psi_\alpha\circ p$ is a
$\theta$-psh model function and $u\circ p\in\PSH(C,\theta)$ is bounded,
all by \Cref{prop:ch21-curve-retraction}.  By
\Cref{thm:bounded-ma-extension},
$\MA_\theta(\psi_\alpha)\to\MA_\theta(u)$ and
$\MA_\theta(\psi_\alpha\circ p)\to\MA_\theta(u\circ p)$ weakly.  Pushforward
along the continuous map $p$ preserves weak convergence, since
$\int f\,\rd(p_*\nu)=\int(f\circ p)\,\rd\nu$ and $f\circ p$ is continuous
for continuous $f$ (\Cref{lem:pushforward-radon}).  Passing to the limit in
Step~1 gives \eqref{eq:ch25-curve-pushforward}.  The same passage to the
limit in \eqref{eq:ch25-curve-ma}, tested against continuous piecewise
affine functions on $\Gamma$, also gives the distributional identity
\[
 \MA_\theta(u\circ p)
 =\mu_0+\Delta^{\mathrm w}_\Gamma(u|_\Gamma).
\]
\end{proof}

\begin{corollary}
\label{cor:ch25-curve-capacity-formula}
Let
\[
 \mathcal P_\Gamma
 :=\bigl\{f:\Gamma\to[-1,0]\ :\ f=u|_\Gamma\ \text{for some}\
   u\in\mathcal U_\theta\bigr\}.
\]
Then for every Borel $A\subseteq\Gamma$,
\begin{equation}
\label{eq:ch25-curve-capacity}
 \Capa_\theta\bigl(p^{-1}(A)\bigr)
 =\sup_{f\in\mathcal P_\Gamma}
     \bigl(\mu_0+\Delta^{\mathrm w}_\Gamma f\bigr)(A).
\end{equation}
Moreover $\mathcal P_\Gamma$ is contained in
\[
 \mathcal Q_\Gamma:=\bigl\{f:\Gamma\to[-1,0]\ \text{continuous}\ :\
    \mu_0+\Delta^{\mathrm w}_\Gamma f\geq0\bigr\}
\]
and contains every piecewise affine member of $\mathcal Q_\Gamma$ whose
breakpoints lie at points of $\Gamma$ with rational coordinates and whose
values are rational.  Consequently, if the supremum of
$(\mu_0+\Delta^{\mathrm w}_\Gamma f)(A)$ over
$\mathcal Q_\Gamma$ is attained at such
a piecewise affine $f$, then it equals
$\Capa_\theta(p^{-1}(A))$.
\end{corollary}

\begin{proof}
For $u\in\mathcal U_\theta$ we have, by
\Cref{prop:ch25-curve-pushforward} and
the distributional identity at the end of its proof,
\[
 \MA_\theta(u)\bigl(p^{-1}(A)\bigr)=\MA_\theta(u\circ p)(A)
 =\bigl(\mu_0+\Delta^{\mathrm w}_\Gamma(u|_\Gamma)\bigr)(A),
\]
where we used $u\circ p|_\Gamma=u|_\Gamma$.  Taking the supremum over
$u\in\mathcal U_\theta$ gives ``$\leq$'' in
\eqref{eq:ch25-curve-capacity}; and if $f\in\mathcal P_\Gamma$, say
$f=u|_\Gamma$, then $u\circ p\in\mathcal U_\theta$ as well, so the same
display read backwards gives ``$\geq$''.

For the inclusion $\mathcal P_\Gamma\subseteq\mathcal Q_\Gamma$: if
$f=u|_\Gamma$ with $u\in\mathcal U_\theta$, then $u\circ p$ is
$\theta$-psh.  The model approximation used in Step~2 of
\Cref{prop:ch25-curve-pushforward}, together with
\eqref{eq:ch25-curve-ma}, gives
$\MA_\theta(u\circ p)=\mu_0+\Delta^{\mathrm w}_\Gamma f\geq0$; and $f$ is continuous because
$\theta$-psh functions restrict continuously to dual complexes
(\Cref{thm:bfj-skeletal-equicontinuity}).  Conversely let $f$ be piecewise
affine on $\Gamma$ with
$\mu_0+\Delta^{\mathrm w}_\Gamma f\geq0$, with values in
$[-1,0]$, and with rational breakpoints and values.  Then $f\circ p$ is a
model function: by \Cref{thm:skeleton-subdivision} there is a vertical blowup
of $\mathcal C$ with the same skeleton $\Gamma$, subdivided so that
$f$ is affine on each cell, and $f\circ p$ is the model function of the
corresponding vertical $\mathbf Q$-Cartier divisor
(\Cref{lem:model-function-affine}).  The weighted finite criterion
\eqref{eq:ch21-curve-psh-test}, or equivalently
\Cref{thm:ch21-model-consistency}, makes it $\theta$-psh.  Its values lie in
$[-1,0]$, so $f\circ p\in\mathcal U_\theta$ and
$f=(f\circ p)|_\Gamma\in\mathcal P_\Gamma$.  The final assertion is
immediate from the two inclusions.
\end{proof}

\subsection{The two-vertex graph}
\label{sec:ch25-two-vertex}

We now specialize to the running example.  Let $\Gamma\subset C^{\An}$
be a skeleton whose underlying metric graph is a segment of length
$\ell>0$; we suppress any valence-two vertices and parametrize it by
$x\in[0,\ell]$.  Assume here that every edge weight is one---for example,
all nodes in the chain are residue-rational---and that the reference measure on $\Gamma$ is
supported at the endpoints $x_{E_0}$ and $x_{E_1}$.  This includes a
regular semistable chain on which the reference line bundle has degree zero
on every intermediate component.  Put
\[
 m_0:=\deg(\theta|_{E_0}),\qquad m_1:=\deg(\theta|_{E_1}),
 \qquad \mu_0=m_0\delta_0+m_1\delta_\ell,
 \qquad V=m_0+m_1 .
\]
On a segment the Laplacian of a continuous $f$ is its distributional second
derivative in the interior together with the outgoing slopes at the two
endpoints, so
\begin{equation}
\label{eq:ch25-segment-class}
 \mathcal Q_\Gamma
 =\Bigl\{f:[0,\ell]\to[-1,0]\ \text{convex}\ :\
   f'(0^+)\geq-m_0,\ f'(\ell^-)\leq m_1\Bigr\},
\end{equation}
and for $f\in\mathcal Q_\Gamma$,
\begin{equation}
\label{eq:ch25-segment-measure}
 \mu_0+\Delta_\Gamma f
 =\bigl(m_0+f'(0^+)\bigr)\delta_0
  +f''|_{(0,\ell)}
  +\bigl(m_1-f'(\ell^-)\bigr)\delta_\ell ,
\end{equation}
a positive measure of total mass $m_0+m_1=V$, since
$f''((0,\ell))=f'(\ell^-)-f'(0^+)$.

\begin{proposition}
\label{prop:ch25-segment-capacities}
With the notation above, write $x_c\in\Gamma$ for the point at distance
$c$ from $x_{E_0}$.  Then
\begin{align}
\label{eq:ch25-cap-endpoint}
 \Capa_\theta\bigl(p^{-1}(\{x_0\})\bigr)
 &=m_0+\min\Bigl\{m_1,\ \frac1\ell\Bigr\},\\
\label{eq:ch25-cap-interval}
 \Capa_\theta\bigl(p^{-1}([a,b])\bigr)
 &=\min\Bigl\{m_0,\frac1a\Bigr\}+\min\Bigl\{m_1,\frac1{\ell-b}\Bigr\},
 \qquad 0<a\leq b<\ell,
\end{align}
and in particular
$\Capa_\theta(p^{-1}(\{x_c\}))
 =\min\{m_0,1/c\}+\min\{m_1,1/(\ell-c)\}$ for $0<c<\ell$.
All three suprema are attained by continuous piecewise affine graph
potentials.  When the metric, the reference masses and the indicated
breakpoints are rational, these may be chosen with rational data and hence
are model functions.  In general they are uniform limits of such model
functions.
\end{proposition}

\begin{proof}
By \Cref{cor:ch25-curve-capacity-formula} it suffices to maximize
$(\mu_0+\Delta_\Gamma f)(A)$ over $f\in\mathcal Q_\Gamma$ and to exhibit
a piecewise affine maximizer.

\emph{The endpoint.}  By \eqref{eq:ch25-segment-measure} the mass at $x_0$
is $m_0+f'(0^+)$.  Convexity gives
$f'(0^+)\leq f'(\ell^-)\leq m_1$, and
$f(\ell)-f(0)\geq\ell f'(0^+)$ together with $-1\leq f\leq0$ gives
$f'(0^+)\leq1/\ell$.  Hence
$m_0+f'(0^+)\leq m_0+\min\{m_1,1/\ell\}$.  Equality holds for the affine
function $f(x)=-1+sx$ with $s=\min\{m_1,1/\ell\}$, which lies in
$\mathcal Q_\Gamma$ because $s\ell\leq1$ and $s\geq0\geq-m_0$.

\emph{An interval.}  By \eqref{eq:ch25-segment-measure} the mass on
$[a,b]\subseteq(0,\ell)$ is $f''([a,b])=f'(b^+)-f'(a^-)$.  Convexity and
$-1\leq f\leq0$ give
\[
 f'(a^-)\geq\frac{f(a)-f(0)}{a}\geq-\frac1a,
 \qquad
 f'(b^+)\leq\frac{f(\ell)-f(b)}{\ell-b}\leq\frac1{\ell-b},
\]
while $f'(a^-)\geq f'(0^+)\geq-m_0$ and
$f'(b^+)\leq f'(\ell^-)\leq m_1$.  Hence
$f''([a,b])\leq\min\{m_0,1/a\}+\min\{m_1,1/(\ell-b)\}$.

For equality put $\alpha:=\min\{m_0,1/a\}$ and
$\beta:=\min\{m_1,1/(\ell-b)\}$, and let
\[
 f(x):=\begin{cases}
  -1+\alpha(a-x),&0\leq x\leq a,\\
  -1,&a\leq x\leq b,\\
  -1+\beta(x-b),&b\leq x\leq\ell .
 \end{cases}
\]
This $f$ is continuous and convex, its slopes being
$-\alpha\leq0\leq\beta$ in that order.  Its values lie in $[-1,0]$: the
minimum is $-1$, and the two maxima are $f(0)=-1+\alpha a\leq0$ and
$f(\ell)=-1+\beta(\ell-b)\leq0$, because $\alpha a\leq1$ and
$\beta(\ell-b)\leq1$ by the definition of $\alpha$ and $\beta$.  The
two endpoint constraints in \eqref{eq:ch25-segment-class} read
$-\alpha\geq-m_0$ and $\beta\leq m_1$, again true by definition.
Finally $f''$ consists of the two atoms $\alpha\delta_a$ and
$\beta\delta_b$, so $f''([a,b])=\alpha+\beta$.  The breakpoints and
values are rational whenever $a,b,\ell,m_0,m_1$ are, so
\Cref{cor:ch25-curve-capacity-formula} identifies the supremum with the
capacity directly in that case.

For real reference masses and real indicated breakpoints we keep
$\theta$ fixed.  The skeletal length $\ell$ is rational in the standing
discrete valuation normalization.  Choose rational
$a_j,b_j\in(0,\ell)$ with $a_j\leq b_j$, converging to $a,b$
(take $a_j=b_j$ when $a=b$), and rational nonnegative numbers
$\alpha_j\leq\min\{m_0,1/a_j\}$,
$\beta_j\leq\min\{m_1,1/(\ell-b_j)\}$ converging to
$\alpha,\beta$.  The displayed three-piece formula with these parameters
has rational vertex values, lies in $\mathcal Q_\Gamma$, and converges
uniformly to the maximizer.  For the endpoint use rational
$0\leq s_j\leq\min\{m_1,1/\ell\}$ with $s_j\to s$ and
$f_j(x)=-1+s_jx$.  In both cases
\Cref{cor:ch25-curve-capacity-formula} makes $f_j\circ p$ a
$\theta$-psh model function, without changing the reference masses.
Uniform stability
\Cref{thm:ch21-stability}\,\cref{itm:ch21-stab:unif} shows that the uniform
limit $f\circ p$ is bounded $\theta$-psh.  Its measure is
$\mu_0+\Delta_\Gamma f$, by
\Cref{prop:ch25-curve-pushforward} and \eqref{eq:ch25-curve-ma}; it therefore
attains the upper bound just proved.  The case $a=b$ is included: the two
atoms then coincide, even when their common breakpoint is irrational, and
give the last display of the statement.
\end{proof}

\begin{remark}[Reading the formulas]
\label{rem:ch25-segment-reading}
Three features of \Cref{prop:ch25-segment-capacities} are worth noting.  The
answer depends on the \emph{metric} of the graph, hence on the total modulus
of the annular chain, and not only on its combinatorics; the two terms in
\eqref{eq:ch25-cap-interval} are the two ``fluxes'' that a normalized
potential can push into the interval from the two ends; and the truncation
$\min\{m_i,\cdot\}$ records which of the two constraints---the bound
$-1\leq f\leq0$ or the semipositivity of $\theta$ at the corresponding
vertex---is active.  When $\ell$ is small the bound $1/\ell$ is inactive
and the capacity saturates at $V$; when $\ell$ is large the geometry
takes over.
\end{remark}

\subsection{Thuillier's relative capacity}
\label{sec:ch25-thuillier}
\index{capacity!relative, on a curve}
\index{Thuillier, Amaury!relative capacity}

The capacity \eqref{eq:capacity-definition} is not the only one available on
a curve, and it is not the one for which the sharpest theorems are known.
Potential theory on Berkovich curves has a complete capacity theory of its
own, due to Thuillier \cite{Thu05}, built not from Monge--Amp\`ere masses but
from energies of measures, in exact imitation of the classical logarithmic
capacity.  We state it here, in the notation of this book, for three
reasons: it settles unconditionally, in dimension one, the three questions
that \Cref{prop:ch25-choquet-axioms} and
\Cref{cor:ch25-polar-countable} leave open; it explains what a
``relative'' capacity is, which is what the local Bedford--Taylor theory
uses and what \eqref{eq:capacity-definition} deliberately is not; and its
normalizations are exactly those of \Cref{sec:ch08-harmonic}, so that no
conversion is needed on the non-Archimedean side.

Throughout this subsection $C$ is a connected, proper, smooth, strictly
$K$-analytic curve; $I(C)\subseteq C^{\An}$ denotes Thuillier's set of points of types
$\mathrm{II}$ and $\mathrm{III}$, equivalently the points occurring in
Shilov boundaries of analytic affinoid domains
\cite[\S3.2]{Thu05}.  In the standing discretely valued algebraic
setting this is the quasimonomial locus of
\Cref{rem:quasimonomial-vs-divisorial}; the analytic definition is used
in the broader source statements below; $\mathrm{SH}(\Omega)$ denotes the
cone of Thuillier-subharmonic functions on an open $\Omega\subseteq C^{\An}$
in the sense of \Cref{sec:ch19-wanner} (they are upper semicontinuous with
values in $[-\infty,\infty)$ and are not identically $-\infty$ on any
connected component); and $\ddc$ is the operator of
\Cref{sec:ch08-harmonic}, with the base-change-compatible edge weights;
on residue-rational simply semistable skeleta it is
\eqref{eq:graph-laplacian}.  Following \cite{Thu05} we call
$E\subseteq C^{\An}$ \emph{polar} if it admits a countable cover by
connected open sets $\Omega$ on each of which there is
$u\in\mathrm{SH}(\Omega)$ with $E\cap\Omega\subseteq\{u=-\infty\}$; the
countability requirement is not cosmetic, and
\Cref{rem:ch25-thuillier-vs-bfj} records the counterexample that forces it.

\begin{definition}[Relative capacity; {\cite[\S3.5, \S3.6.1]{Thu05}}]
\label{def:ch25-thuillier-capacity}
Let $\Omega\subseteq C^{\An}$ be open with $\partial\Omega$ a nonempty
subset of $I(C)$.  For a positive Radon measure $\mu$ on $\Omega$ use its Dirichlet
potential $g_\mu$ in the sense of \cite[\S3.4.3]{Thu05}, and set
\[
 E_\Omega(\mu):=-\int_\Omega g_\mu\,\rd\mu\in[0,+\infty] ,
\]
the \emph{energy} of $\mu$ relative to $\Omega$.  Finite-energy
measures have the subharmonic Dirichlet representatives characterized in
\cite[Proposition~3.5.10]{Thu05}; no classical finite-valued solution
for an arbitrary measure is assumed.  For $K\subseteq\Omega$
compact the \emph{relative capacity} is
\begin{equation}
\label{eq:ch25-thuillier-capacity}
 C(K,\Omega):=\Bigl(\inf\bigl\{E_\Omega(\mu):
   \mu\in\operatorname{Prob}(K)\bigr\}\Bigr)^{-1}
 =\Bigl(\sup\bigl\{\mu(K):\mu\geq0,\ E_\Omega(\mu)=1\bigr\}\Bigr)^{2},
\end{equation}
with the conventions $(+\infty)^{-1}=0$ and $\sup\varnothing=0$.  It is
extended to arbitrary $E\subseteq\Omega$ by inner regularity,
$C(E,\Omega):=\sup\{C(K,\Omega):K\subseteq E\ \text{compact}\}$, and the
\emph{outer relative capacity} is
\begin{equation}
\label{eq:ch25-thuillier-outer}
 C^*(E,\Omega):=\inf\bigl\{C(U,\Omega):
    E\subseteq U\subseteq\Omega,\ U\ \text{open}\bigr\}.
\end{equation}
Finally, for $E\subseteq\Omega$ the \emph{relative extremal function} is
the upper semicontinuous regularization $u^*_E$ of the upper envelope of
\[
 \mathcal F(E):=\bigl\{v\in\mathrm{SH}(\Omega):v\leq0\ \text{and}\
    v|_E\leq-1\bigr\}.
\]
\end{definition}

\refstepcounter{theorem}\label{thm:ch25-thuillier-capacity}
\begin{quote}
\textbf{Foundational input \thetheorem\ (Thuillier's relative capacity;
source boundary).}
Let $C$ be a connected proper smooth strictly $K$-analytic curve and let
$\Omega\subseteq C^{\An}$ be open with $\partial\Omega$ a nonempty
subset of $I(C)$.  Then, with the notation of
\Cref{def:ch25-thuillier-capacity}:
\begin{enumerate}
\item\label{itm:ch25-thu:extremal}
 for every compact $K\subseteq\Omega$ one has
 $C(K,\Omega)=\int_K\ddc u^*_K<\infty$, and for every $K$-affinoid
 domain $V\subseteq\Omega$ one has
 $C(V,\Omega)=C(\partial V,\Omega)=\int_\Omega\ddc u^*_V$;
\item\label{itm:ch25-thu:descend}
 $C(\cdot,\Omega)$ descends on compacts: if $K_1\supseteq K_2\supseteq
 \cdots$ are compact with intersection $K$, then
 $C(K_j,\Omega)\searrow C(K,\Omega)$;
\item\label{itm:ch25-thu:strong}
 $C(\cdot,\Omega)$ is strongly subadditive on compacts,
 \[
  C(K_1\cup K_2,\Omega)+C(K_1\cap K_2,\Omega)
  \leq C(K_1,\Omega)+C(K_2,\Omega);
 \]
\item\label{itm:ch25-thu:outer}
 $C^*(\cdot,\Omega)$ is monotone, agrees with $C(\cdot,\Omega)$ on
 compact sets, is strongly subadditive on arbitrary subsets, and is
 continuous along increasing sequences of \emph{arbitrary} subsets; thus it
 is a Choquet capacity in the sense of \Cref{sec:capacity-preview}, with
 strong subadditivity available in addition;
\item\label{itm:ch25-thu:capacitable}
 every analytic subset of $\Omega$, in particular every Borel subset
 $B$, is capacitable: $C^*(B,\Omega)=C(B,\Omega)$;
\item\label{itm:ch25-thu:rel-compact}
 for every relatively compact $E\subseteq\Omega$ one has
 $C^*(E,\Omega)=\int_\Omega\ddc u^*_E$.
\end{enumerate}
Items \cref{itm:ch25-thu:extremal}--\cref{itm:ch25-thu:rel-compact} are
imported in full from \cite[Propositions~3.6.3, 3.6.5, 3.6.6, 3.6.8,
3.6.10 and Th\'eor\`eme~3.6.9]{Thu05}; no proof is claimed here.
\end{quote}

Item \cref{itm:ch25-thu:extremal} is the point of contact with the present
chapter: it says that the energy-theoretic capacity is computed by the
Monge--Amp\`ere mass of a single extremal potential, which is what
\eqref{eq:capacity-definition} takes as its definition.  In the complex
theory the corresponding identification of the Bedford--Taylor capacity with
the mass of the relative extremal function is \cite[Proposition~4.2]{BT82}.

\refstepcounter{theorem}\label{thm:ch25-thuillier-polar}
\begin{quote}
\textbf{Foundational input \thetheorem\ (polar sets and capacity zero on a
curve; source boundary).}
Let $C$ and $\Omega$ be as in \Cref{thm:ch25-thuillier-capacity} and
assume in addition that $\Omega$ is connected.  For a subset
$E\subseteq\Omega$ the following are equivalent:
\begin{enumerate}
\item\label{itm:ch25-thupol:polar} $E$ is polar;
\item\label{itm:ch25-thupol:cap} $C^*(E,\Omega)=0$;
\item\label{itm:ch25-thupol:extremal} $u^*_E\equiv0$;
\item\label{itm:ch25-thupol:single} there is $u\in\mathrm{SH}(\Omega)$ with
 $u\leq0$ and $E\subseteq\{u=-\infty\}$.
\end{enumerate}
Moreover every measure of finite energy gives measure zero to every polar
set, and a point $x\in C^{\An}$ is polar if and only if it is of type
$\mathrm{I}$ in Thuillier's classification; in
particular no quasimonomial point is polar.  Here type
$\mathrm{I}$ has Thuillier's general-field meaning: the completed residue
field embeds over $K$ into a completed algebraic closure of $K$
\cite[\S2.1.3, Proposition~3.4.10]{Thu05}.  Under algebraically closed
$K$ this agrees with the type-$\mathrm{I}$ classification in
\Cref{thm:berk-a1-classification}.  The hypothesis that
$\Omega$ be connected may be weakened to $\pi_0(\Omega)$ countable, and
cannot be dropped.
The equivalences and the final assertions are imported from
\cite[Th\'eor\`eme~3.5.11, Proposition~3.4.10 and
Th\'eor\`eme~3.6.11]{Thu05}; no proof is claimed here.
\end{quote}

\begin{remark}[Two capacities, and what does and does not transfer]
\label{rem:ch25-thuillier-vs-bfj}
\Cref{thm:ch25-thuillier-capacity} and \Cref{thm:ch25-thuillier-polar} are
statements about $C(\cdot,\Omega)$, not about $\Capa_\theta$, and the two
set functions are different objects even when $n=1$.  Four
differences must be kept in view.

\emph{Relative versus global.}  $C(\cdot,\Omega)$ is attached to an open
$\Omega$ with nonempty boundary inside $I(C)$; the boundary is what
makes the Dirichlet problem solvable and the energy finite.  There is no
reference form: the normalization is carried by $\partial\Omega$, where
competitors are required to vanish.  By contrast
\eqref{eq:capacity-definition} is global on $X^{\An}$ and normalized by
$\theta$ and by the constraint $-1\leq u\leq0$.  Consistently,
$\Capa_\theta(X^{\An})=V$ is finite
(\Cref{cor:ch25-negligible-empty-interior}) whereas
$\sup\{C(V,\Omega):V\subseteq\Omega\ \text{affinoid}\}=+\infty$
\cite[Lemme~3.6.28]{Thu05}.

\emph{Energy versus mass.}  \eqref{eq:ch25-thuillier-capacity} is an
inverse minimal energy; \eqref{eq:capacity-definition} is a supremum of
masses.  They are reconciled on compacts by
\Cref{thm:ch25-thuillier-capacity}\,\cref{itm:ch25-thu:extremal}, but only
after the relative extremal function has been constructed, which is a
theorem about $\Omega$ and has no analogue here.

\emph{What is therefore not claimed.}  Nothing in
\Cref{thm:ch25-thuillier-capacity} is asserted for $\Capa_\theta$.  In
particular strong subadditivity of $\Capa_\theta$, and with it Choquet
capacitability, remain unproved in this book even for curves; the warning
after \Cref{prop:ch25-choquet-axioms} stands as written.  Likewise the
converse ``negligible $\Rightarrow$ pluripolar'' is not deduced for
$\Capa_\theta$ from
\Cref{thm:ch25-thuillier-polar}\,\cref{itm:ch25-thupol:cap}.

\emph{What is nevertheless learnt.}  Two facts do transfer, because they are
statements about sets and subharmonic functions rather than about a
particular capacity.  First,
\Cref{thm:ch25-thuillier-polar} identifies the polar subsets of a curve as
exactly the sets swallowed by one subharmonic function, and locates them
inside the type~$\mathrm{I}$ points; this is a sharpening of the
qualitative statement \Cref{lem:divisorial-positive-capacity} in dimension
one, where the dichotomy is complete rather than one-sided.  Second, the
countability clause in the definition of ``polar'' is unavoidable: over a
residue field of cardinality greater than $\aleph_0$, the set of type~I
points reducing to the distinct residues of the Gauss point of
$\mathbf A^{1,\An}_K$ is closed in
$\Omega=\mathbf A^{1,\An}_K\setminus\{\zeta_{0,1}\}$, lies in the
$-\infty$ locus of a subharmonic function on each component, and yet has
$C^*(\cdot,\Omega)=+\infty$ \cite[Remarque~3.6.12]{Thu05}.  This is the
non-Archimedean phenomenon that has no complex counterpart, and it is
another instance of the discipline of \Cref{sec:ch25-nets}: an uncountable
family of components is not a sequence.
\end{remark}

\begin{remark}[Normalizations: none to convert, except against the complex
model]
\label{rem:ch25-thuillier-normalization}
Thuillier's $\ddc$ is characterized by compatibility with base change
together with the requirement that on the skeleton of a simply semistable
formal model it be the integral-affine Laplacian; on every skeleton it is
the weighted outgoing-slope operator in the general-field setting.
On a residue-rational simply semistable skeleton it reduces to
\eqref{eq:graph-laplacian}.  His Poincar\'e--Lelong formula reads
$\ddc\log|f|=\delta_{\operatorname{div}(f)}$
\cite[Proposition~3.3.15]{Thu05}, with no factor $\tfrac12$.  This is
exactly the convention of \Cref{sec:ch08-harmonic} and
\Cref{thm:ch08-slope-formula}, and $\log$ is the natural logarithm on both
sides.  No rescaling of any formula quoted above is required, and the
lengths on skeleta are the modulus lengths of
\Cref{def:ch08-path-metric}.

One factor does have to be declared.  Thuillier presents his Dirichlet form
$(\varphi,\psi)\mapsto-\int\varphi\,\ddc\psi$ as the non-Archimedean
analogue of
$\frac{\rmi}{\pi}\int_M\partial\varphi\wedge\bar\partial\psi$
on a compact Riemann surface, that is, of the complex operator
$\frac{\rmi}{\pi}\partial\bar\partial$.  The book's complex operator is
$\ddc=\frac{\rmi}{2\pi}\partial\bar\partial$ \eqref{eq:appJ-ddc}, for which
$\ddc\log|f|=\tfrac12[\operatorname{div}(f)]$
\eqref{eq:appJ-complex-pl}.  Thuillier's complex model is thus twice ours:
\[
 \ddc_{\text{\cite{Thu05}}}
 =\frac{\rmi}{\pi}\partial\bar\partial
 =2\,\ddc_{\eqref{eq:appJ-ddc}}
 \qquad\text{on a complex curve.}
\]
The discrepancy is confined to the archimedean side of the analogy: it
affects no non-Archimedean formula in this book, but it must be inserted by
hand whenever a constant is carried from his archimedean comparison
statements into \Cref{app:normalizations}.  We have not rescaled anything
silently; every display above is in his non-Archimedean normalization, which
coincides with ours.
\end{remark}

\begin{example}[The relative capacity of the unit disc, computed]
\label{ex:ch25-thuillier-disc}
Take $C=\mathbf P^1_K$ with $K$ algebraically closed and nontrivially
valued, fix $r>1$, and let
$\Omega=D(0,r)=\{|T|<r\}\subseteq\mathbf A^{1,\An}_K$.  Its topological
boundary in $C^{\An}$ is the single point $\zeta_{0,r}$, which is
quasimonomial, so \Cref{def:ch25-thuillier-capacity} applies.  Let
$E=E(0,1)=\{|T|\leq1\}$ be the closed unit disc, a compact subset of
$\Omega$.

Parametrize the segment $[\zeta_{0,r},\zeta_{0,1}]$ by arclength
$s\in[0,\log r]$; its total length is $\log r$ by
\Cref{def:ch08-path-metric}, since $t(\zeta_{0,\rho})=-\log\rho$.  Let $\tau$ retract $\Omega$ onto the radial segment
$[\zeta_{0,1},\zeta_{0,r})$, collapsing the closed unit disc to
$\zeta_{0,1}$ and each attached branch to its point on the segment.
If $s(\tau(x))=\log(r/\rho)$ for
$\tau(x)=\zeta_{0,\rho}$, define
\[
 u(x):=-\frac{s(\tau(x))}{\log r}.
\]
This defines $u$ on every branch as well as the radial segment; it
is continuous on $\Omega$, takes values in $[-1,0]$, equals $-1$ on
$E$, tends to $0$ at $\zeta_{0,r}$, and is harmonic on
$\Omega\setminus E$ because it is affine on the segment with no interior
break point and constant on the fibres of the retraction.  Its only
Laplacian atom in $\Omega$ is positive---see the next display---so $u$ is
subharmonic \cite[Th\'eor\`eme~3.4.12]{Thu05}.  It is therefore the relative
extremal function of $E$ in $\Omega$, by the uniqueness clause of
\cite[Th\'eor\`eme~3.6.13]{Thu05}: $u^*_K$ is the unique function
subharmonic on $\Omega$, harmonic on $\Omega\setminus K$, vanishing on
$\partial\Omega$ and equal to $-1$ quasi-everywhere on $K$, and $u$
has all four properties.  Hence $u^*_E=u$.  Its Laplacian
\eqref{eq:graph-laplacian} has a single atom in $\Omega$, at
$\zeta_{0,1}$, where the only nonzero outgoing slope is the one along the
segment, equal to $1/\log r$; the residue directions into $E(0,1)$
contribute $0$ because $u$ is constant there.  Therefore
\begin{equation}
\label{eq:ch25-thuillier-disc}
 C\bigl(E(0,1),D(0,r)\bigr)=\int_E\ddc u^*_E=\frac1{\log r},
\end{equation}
by \Cref{thm:ch25-thuillier-capacity}\,\cref{itm:ch25-thu:extremal}, which is
the value recorded in \cite[Remarque~3.6.12]{Thu05}.  The same function is
the relative extremal function of the single point $\zeta_{0,1}$, by the
same uniqueness clause: $u$ is harmonic on
$\Omega\setminus\{\zeta_{0,1}\}$ as well, being locally constant on
$E(0,1)\setminus\{\zeta_{0,1}\}$, and equals $-1$ there.  Hence
\[
 C\bigl(\{\zeta_{0,1}\},D(0,r)\bigr)=\frac1{\log r}
 =C\bigl(E(0,1),D(0,r)\bigr).
\]
This is the instance $g(\zeta_{0,1})=\log r$ of Thuillier's formula
$C(\{x\},\Omega)=g(x)^{-1}$, where $g\geq0$ is the equilibrium
potential of $x$ in $\Omega$, continuous, vanishing on
$C^{\An}\setminus\Omega$ and with $\ddc g|_\Omega=-\delta_x$
\cite[Lemme~3.6.28]{Thu05}: here $g=-(\log r)u$ is the function
equal to $\log r$ on $E(0,1)$ and affine to $0$ along the segment, whose
outgoing-slope sum at $\zeta_{0,1}$ is $-1$.  Letting $x\to\partial\Omega$
makes $g(x)\to0$ and $C(\{x\},\Omega)\to+\infty$, which is why the
relative capacity is unbounded on $\Omega$ while $\Capa_\theta$ is
bounded by $V$.

Two contrasts with \Cref{prop:ch25-segment-capacities} are worth stating
explicitly.  There the ambient object is a proper curve with a reference
form $\theta$ and the answer is capped by $V=\deg\{\theta\}$; here there is no
reference form, the cap is replaced by the distance to $\partial\Omega$,
and the answer diverges as that distance goes to zero.  And there the
computation is a convex optimization over the whole graph
\eqref{eq:ch25-curve-capacity}; here it is the solution of one Dirichlet
problem.  These formulas belong to different normalizations and ambient domains;
no numerical comparison between them is inferred here.
\end{example}

\refstepcounter{theorem}\label{thm:ch25-thuillier-classical-points}
\begin{quote}
\textbf{Foundational input \thetheorem\ (Thuillier's determination by
classical points; source boundary).}
Let $C$ be a smooth strictly $K$-analytic curve and put
\[
 C_0:=\bigl\{x\in C^{\An}:[\kappa(x):K]<\infty\bigr\}.
\]
Then a subharmonic function on $C^{\An}$ is uniquely determined by its
restriction to $C_0$.  The proof is a capacity statement: for every open
$\Omega$ with $\partial\Omega$ a nonempty subset of $I(C)$ and every
$K$-affinoid domain $V\subseteq\Omega$,
\[
 C(V\cap C_0,\Omega)=C(V,\Omega),
\]
so that $\Omega\cap C_0$ contains compacts of arbitrarily large relative
capacity, and $C_0$ is dense for the fine topology generated by the
subharmonic functions.
The determination statement and its capacity mechanism are imported from
\cite[Th\'eor\`eme~3.6.26 and Proposition~3.6.29]{Thu05}; no proof is
claimed here.
\end{quote}

\begin{remark}[Which determination statement this is]
\label{rem:ch25-thuillier-determination}
\Cref{thm:ch25-thuillier-classical-points} runs in the opposite direction to
the determination statements used elsewhere in \Cref{part:bfj}.  A
BFJ-psh function is determined by its restriction to the quasimonomial
points essentially by construction, since $\PSH(X,\theta)$ is built from
model functions and every model function factors through a dual complex; that
is the content of \Cref{thm:bfj-skeletal-equicontinuity} and of the
regularization theorem \Cref{thm:bfj-decreasing-regularization}.  Thuillier's
theorem asserts determination by the type~$\mathrm{I}$ points---the points
that carry \emph{no} skeleton and that
\Cref{thm:ch25-thuillier-polar} identifies as the only candidates for
polarity.  It is a genuinely different and harder statement, and this book
proves no analogue of it in dimension $\geq2$; the capacity input it uses,
$C(V\cap C_0,\Omega)=C(V,\Omega)$, has no counterpart for
$\Capa_\theta$.
\end{remark}

\section{Worked examples}
\label{sec:ch25-worked}

\begin{example}[A point can have full capacity]
\label{ex:ch25-two-vertex-full-capacity}
Return to the two-vertex graph of \Cref{ex:two-vertex-ma}, take edge
length one, and let
$\theta=\delta_{x_0}+\delta_{x_1}$, so $V=2$. An admissible psh function
has vertex values $s_0,s_1\in[-1,0]$, and
\[
 \MA_\theta(u)(\{x_0\})=1+s_1-s_0.
\]
The positivity inequalities are automatic when $|s_1-s_0|\leq1$. Taking
$s_0=-1,s_1=0$ gives mass $2$ at $x_0$, hence
\[
                 \Capa_\theta(\{x_0\})=2=V.
\]
Thus capacity is not a disguised topological size: a single divisorial point
may have full capacity. This is exactly why measures supported on a finite
skeleton are visible to the theory.
\end{example}

The computation is consistent with
\eqref{eq:ch25-cap-endpoint}: with $m_0=m_1=1$ and $\ell=1$,
\[
 \Capa_\theta\bigl(p^{-1}(\{x_0\})\bigr)=1+\min\{1,1\}=2=V,
\]
and since $\{x_0\}\subseteq p^{-1}(\{x_0\})$ while the maximizer places all
of its mass at the single point $x_0$, the capacity of the point itself is
also $2$.  The next example shows that the coincidence is an artefact of
the short edge.

\begin{example}[Capacities on a segment of length four]
\label{ex:ch25-segment-capacities}
Let $\mathcal C$ be a regular semistable model whose special fibre has a
chain $E_0-F_1-F_2-F_3-E_1$ of five components and four residue-rational
ordinary nodes.
After suppressing the three valence-two vertices, its metric graph is
$\Gamma=[0,4]$.  Choose a semipositive reference form whose component
degrees are $1$ on $E_0,E_1$ and $0$ on $F_1,F_2,F_3$; thus
$m_0=m_1=1$ and $V=2$.  \Cref{prop:ch25-segment-capacities}
gives the following table, in which $x_c$ is the point of $\Gamma$ at
distance $c$ from $x_{E_0}$.

\begin{center}
\small
\begin{tabular}{l|l|l}
\toprule
\textbf{Set }$A\subseteq\Gamma$&$\Capa_\theta(p^{-1}(A))$
 &\textbf{maximizer}\\
\midrule
$\{x_0\}$&$1+\tfrac14=\tfrac54$&$f(x)=-1+x/4$\\
$\{x_{1/2}\}$&$1+\tfrac27=\tfrac97$&kink at $x=1/2$\\
$\{x_1\}$&$1+\tfrac13=\tfrac43$&kink at $x=1$\\
$\{x_2\}$&$\tfrac12+\tfrac12=1$&kink at $x=2$\\
$\{x_3\}$&$\tfrac13+1=\tfrac43$&kink at $x=3$\\
$[1,2]$&$1+\tfrac12=\tfrac32$&flat on $[1,2]$\\
$[2,3]$&$\tfrac12+1=\tfrac32$&flat on $[2,3]$\\
$[1,3]$&$1+1=2=V$&flat on $[1,3]$\\
$\Gamma$&$2=V$&$f\equiv0$\\
\bottomrule
\end{tabular}
\end{center}

Every entry is exact, the maximizers being piecewise affine with integer or
half-integer breakpoints.  For instance the entry for $[1,3]$ is realized
by
\[
 f(x)=\begin{cases}-x,&0\leq x\leq1,\\ -1,&1\leq x\leq3,\\
  x-4,&3\leq x\leq4,\end{cases}
\]
whose Monge--Amp\`ere measure is
$\MA_\theta(f\circ p)=\delta_{x_1}+\delta_{x_3}$; both endpoint terms in
\eqref{eq:ch25-segment-measure} vanish, since
$m_0+f'(0^+)=1-1=0$ and $m_1-f'(4^-)=1-1=0$.

Two consequences are worth recording.  First, the capacity of a divisorial
point has no uniform lower bound depending only on $n$ and $V$ across
all these polarized curves: replacing the four-edge chain
by a chain of total length $\ell$ and taking $c=\ell/2$ gives
$\Capa_\theta(p^{-1}(\{x_{\ell/2}\}))=4/\ell$ as soon as
$\ell\geq2/m_0$ and $\ell\geq2/m_1$, which tends to $0$ as
$\ell\to\infty$.  Along even integral $\ell$, the midpoint is the
divisorial vertex of an intermediate component, and the maximizing kink
puts the whole displayed tube mass at that vertex; hence its singleton has
the same capacity $4/\ell$.  Thus
\Cref{lem:divisorial-positive-capacity} is not
a quantitative bound in those parameters alone; this varying-geometry
example makes no claim about an infimum on one fixed curve.  The caveat in
\Cref{rem:ch25-divisorial-explicit} is real.  Second, a proper closed subset
can already have full capacity: $\Capa_\theta(p^{-1}([1,3]))=V$ while
$p^{-1}([1,3])\neq C^{\An}$.
\end{example}

\begin{example}[Subadditivity is strict]
\label{ex:ch25-subadditivity-strict}
In the situation of \Cref{ex:ch25-segment-capacities},
\[
 \Capa_\theta\bigl(p^{-1}([1,3])\bigr)=2,
 \qquad
 \Capa_\theta\bigl(p^{-1}([1,2])\bigr)
 +\Capa_\theta\bigl(p^{-1}([2,3])\bigr)=\tfrac32+\tfrac32=3 ,
\]
so the inequality \eqref{eq:capacity-subadditivity} is strict, with a loss of
$1$.  This is the phenomenon flagged in the proof of
\Cref{prop:capacity-elementary}: no single admissible potential is optimal
for both halves.  The potential achieving $3/2$ on $[1,2]$ has slopes
$(-1,0,\tfrac12)$ on the three pieces $[0,1],[1,2],[2,4]$, and the one
achieving $3/2$ on $[2,3]$ has slopes
$(-\tfrac12,0,1)$ on $[0,2],[2,3],[3,4]$; they are different functions,
and no convex combination of them does better than $2$ on $[1,3]$, since
$2=V$ is the total mass.

The same numbers show that \emph{strong} subadditivity, if true, is not far
from sharp here: with $A=p^{-1}([1,2])$ and $B=p^{-1}([2,3])$,
\[
 \Capa_\theta(A\cup B)+\Capa_\theta(A\cap B)=2+1=3
 =\Capa_\theta(A)+\Capa_\theta(B),
\]
using $\Capa_\theta(p^{-1}(\{x_2\}))=1$ from the table.  We emphasize that
this is one computed instance and not a proof; strong subadditivity is not
established in this book.
\end{example}

\begin{example}[A capacity estimate on a segment]
\label{ex:ch25-fundamental-numbers}
Keep $\Gamma=[0,4]$, $m_0=m_1=1$, $V=2$, $n=1$.  Take
\[
 \varphi(x)=-1+\tfrac x4,\qquad \psi\equiv0,
\]
so $-1\leq\varphi\leq\psi\leq0$ and $M=1$.  By
\eqref{eq:ch25-segment-measure},
\[
 \MA_\theta(\varphi)=\tfrac54\delta_{x_0}+\tfrac34\delta_{x_4},
 \qquad
 \MA_\theta(\varphi/2)=\tfrac98\delta_{x_0}+\tfrac78\delta_{x_4}.
\]
The left-hand side of \eqref{eq:fundamental-capacity-estimate}, with
$\psi-\varphi=1-x/4$, takes the value $1$ for $u=0$ and the value
$\tfrac54$ for $u=\varphi$.  The right-hand side is
\[
 4MV^{1-2^{-1}}\Bigl(\int(\psi-\varphi)\MA_\theta(\varphi/2)\Bigr)^{1/2}
 =4\sqrt2\cdot\Bigl(\tfrac98\Bigr)^{1/2}=6 .
\]
The estimate therefore holds with a loss of a factor between $4.8$ and
$6$.  The shape of the bound---the exponent $2^{-n}=1/2$ and the linear
dependence on $M$---is what the applications use; the constant is not
optimized anywhere in this book, and this example shows the order of
magnitude of the slack.
\end{example}

\begin{example}[A pluripolar set and its capacity, computed exactly]
\label{ex:ch25-pluripolar-computed}
Let $C=\mathbf P^1_K$ with coordinate $T$ and let $\theta$ be the
curvature of the standard model metric on $\cO(1)$, so that by
\eqref{eq:p1-standard-measure}
\[
 \theta=\delta_{\zeta_{0,1}},\qquad V=1,\qquad n=1,
\]
where $\zeta_{0,1}$ is the Gauss point of the unit disc.  Let
\[
 \varphi_0(x):=\log\frac{|T(x)|}{\max\{1,|T(x)|\}} .
\]
This is the logarithmic norm of the section $Z_1$ of $\cO(1)$ in the standard model
metric \eqref{eq:p1-standard-weight}; it satisfies $\varphi_0\leq0$,
$\varphi_0(\zeta_{0,1})=0$, and $\varphi_0=0$ outside the open unit disc.
On the segment from $\zeta_{0,1}$ to the rational point $T=0$,
parametrized by $s\geq0$ through $x_s=\zeta_{0,\rme^{-s}}$, one has
$\varphi_0(x_s)=-s$.  The Poincar\'e--Lelong formula, with the section's
logarithmic norm relative to the reference metric, gives
$\ddc\varphi_0=\delta_0-\delta_{\zeta_{0,1}}$, whence
\[
 \theta+\ddc\varphi_0=\delta_0,
\]
which is the Poincar\'e--Lelong formula \eqref{eq:p1-poincare-lelong} in the
present normalization.  Membership
$\varphi_0\in\PSH(\mathbf P^1,\theta)$ also follows directly from the
decreasing nef model truncations of \Cref{ex:ch23-pole-model-check};
$\sup\varphi_0=0$, and its polar set is the single rational point
$\{T=0\}$, which is therefore pluripolar.

\emph{The capacity of the sublevel sets.}  We claim that for every
$t>0$,
\begin{equation}
\label{eq:ch25-p1-capacity}
 \Capa_\theta\bigl(\{\varphi_0\leq-t\}\bigr)=\min\Bigl\{1,\frac1t\Bigr\} .
\end{equation}
The set in question is the closed ball $\{|T|\leq\rme^{-t}\}$.

For the upper bound, note first that the dual complex of the standard model
$\mathbf P^1_R$ is the single point $\{\zeta_{0,1}\}$, so that
$\theta^n=\theta=\delta_{\zeta_{0,1}}$ and, for any $u$ with
$\sup u=0$, $\int(-u)\theta=-u(\zeta_{0,1})$.  Applying this to
$u=\varphi_0$ gives $\int(-\varphi_0)\theta=0$, and the constant
$C_\theta$ of \Cref{lem:ch25-normalized-bound} may be taken to be $0$,
the diameter of a point being $0$.  Hence
\Cref{prop:ch25-sublevel-capacity} reads
\[
 \Capa_\theta\bigl(\{\varphi_0<-t\}\bigr)\leq\frac{0+1\cdot1}{t}
 =\frac1t ,
\]
and the same bound holds for the closed sublevel set by applying the estimate
at any $t'<t$ and letting $t'\uparrow t$.  Together with
$\Capa_\theta\leq V=1$ this gives ``$\leq$'' in
\eqref{eq:ch25-p1-capacity}.

For the lower bound, let $t\geq1$ and put
$u_t:=\max\{\varphi_0/t,-1\}$, which lies in $\mathcal U_\theta$ because
$\varphi_0/t$ is a convex combination of $\varphi_0$ and $0$.  On the
segment, $u_t(x_s)=\max\{-s/t,-1\}$, affine of slope $-1/t$ for
$s\leq t$ and constant $-1$ for $s\geq t$.  On attached branches it
has the value at the radial attachment point; it is $0$ outside the open
unit disc.  For rational $t\geq1$ the model graph Laplacian has a single negative contribution at
$\zeta_{0,1}$, of size $-1/t$, and a single positive kink at $x_t$, of
size $1/t$, so
\begin{equation}
\label{eq:ch25-p1-truncated-ma}
 \MA_\theta(u_t)=\Bigl(1-\frac1t\Bigr)\delta_{\zeta_{0,1}}
   +\frac1t\,\delta_{x_t},
\end{equation}
a positive measure of mass $1$.  For real $t\geq1$, rational
$t_j\geq1$ tending to $t$ give uniform convergence $u_{t_j}\to u_t$
(the scalar functions $\max\{-s/t_j,-1\}$ converge uniformly for
$s\in[0,+\infty]$).  Weak continuity from
\Cref{prop:ch24-continuous-uniform} and continuity of $t\mapsto x_t$
give the same displayed measure for real $t$.
Since $\varphi_0(x_t)=-t$, the point $x_t$ lies in
$\{\varphi_0\leq-t\}$, whence
$\Capa_\theta(\{\varphi_0\leq-t\})\geq1/t$.  For $0<t\leq1$ use
$u_1$, whose measure is $\delta_{x_1}$ and whose atom lies in
$\{\varphi_0\leq-1\}\subseteq\{\varphi_0\leq-t\}$; this gives the lower
bound $1=V$.  This proves \eqref{eq:ch25-p1-capacity}.

The same value holds for the \emph{open} sublevel set: monotonicity gives
$\Capa_\theta(\{\varphi_0<-t\})\leq\Capa_\theta(\{\varphi_0\leq-t\})$, and
for $t'>t$ the inclusion
$\{\varphi_0\leq-t'\}\subseteq\{\varphi_0<-t\}$ gives the reverse bound
$\min\{1,1/t'\}$; letting $t'\downarrow t$ yields
\begin{equation}
\label{eq:ch25-p1-capacity-open}
 \Capa_\theta\bigl(\{\varphi_0<-t\}\bigr)=\min\Bigl\{1,\frac1t\Bigr\},
 \qquad t>0 .
\end{equation}

\emph{Consequences.}  Letting $t\to\infty$ in
\eqref{eq:ch25-p1-capacity} gives
$\Capa_\theta(\{T=0\})=0$: the pluripolar set is negligible, in accordance
with \Cref{cor:ch25-polar-negligible}.  More is true: the rate
$O(1/t)$ of \eqref{eq:ch25-sublevel-capacity} is attained, and so is its
constant $C_\theta+nV=1$.  Thus \Cref{prop:ch25-sublevel-capacity} is
sharp, and \Cref{rem:ch25-two-polar-proofs} should be read accordingly.

\emph{A nonpolar set of positive capacity.}  By contrast, the single point
$x_s=\zeta_{0,\rme^{-s}}$ is quasimonomial, hence not pluripolar
(\Cref{prop:ch21-polar-structure}\,\cref{itm:ch21-pol:skeleta}), and
\eqref{eq:ch25-p1-truncated-ma} with $t=s\geq1$ shows
$\Capa_\theta(\{x_s\})\geq1/s>0$.  So the negligible sets are strictly
smaller than the sets with empty interior, and the capacity distinguishes
points that the topology does not.
\end{example}

\begin{example}[Convergence in capacity without uniform convergence]
\label{ex:ch25-capacity-not-uniform}
Keep the setting of \Cref{ex:ch25-pluripolar-computed} and put
\[
 w_j:=\max\Bigl\{\frac{\varphi_0}{j},-1\Bigr\}\in\mathcal U_\theta,
 \qquad j\geq1 .
\]
The sequence is increasing and $w_j\leq0$.  For $0<\delta<1$,
\[
 \{|w_j-0|>\delta\}=\{w_j<-\delta\}=\{\varphi_0<-\delta j\},
\]
so by \eqref{eq:ch25-p1-capacity-open},
\[
 \Capa_\theta\bigl(\{|w_j|>\delta\}\bigr)
 =\min\Bigl\{1,\frac1{\delta j}\Bigr\}
 =\frac1{\delta j}\quad\text{as soon as }\delta j\geq1 .
\]
Hence $w_j\to0$ in capacity, with the explicit rate
$(\delta j)^{-1}$.  On the other hand $\|w_j\|_\infty=1$ for every
$j$, because $w_j=-1$ on the ball $\{|T|\leq\rme^{-j}\}$; so the
convergence is not uniform, and it is not even pointwise, since
$w_j(0)=-1$ for all $j$.  This proves the first half of
\Cref{prop:ch25-convergence-comparison}\,\cref{itm:ch25-cc:none}.

Note that the exceptional sets $\{|w_j|>\delta\}$ shrink to the single
pluripolar point $\{T=0\}$; the example is, in this sense, the minimal one.
\end{example}

\begin{example}[Pointwise convergence without convergence in capacity]
\label{ex:ch25-pointwise-not-capacity}
Keep $C=\mathbf P^1_K$ and $\theta=\delta_{\zeta_{0,1}}$, $V=1$.  The
residue field $k$ has characteristic zero, hence is infinite; choose
elements $a_1,a_2,\ldots\in R=K^\circ$ whose residues in $k$ are pairwise
distinct, and set
\[
 \varphi_j(x):=\log\frac{|T(x)-a_j|}{\max\{1,|T(x)|\}},
 \qquad
 u_j:=\max\{\varphi_j,-1\}\in\mathcal U_\theta .
\]
As in \Cref{ex:ch25-pluripolar-computed}, $\varphi_j$ is $\theta$-psh
with $\theta+\ddc\varphi_j=\delta_{a_j}$, $\sup\varphi_j=0$, and
$\varphi_j$ vanishes outside the open ball
$B_j:=\{|T-a_j|<1\}$.  Indeed if $|T(x)|\leq1$ and $x\notin B_j$ then
$|T(x)-a_j|=1$ by the ultrametric inequality and the distinctness of
residues, while if $|T(x)|>1$ then $|T(x)-a_j|=|T(x)|$.

\emph{Pointwise convergence.}  The balls $B_j$ are pairwise disjoint,
because their points have pairwise distinct residues.  Hence every
$x\in(\mathbf P^1)^{\An}$ lies in at most one $B_j$, so $u_j(x)=0$ for
all but at most one $j$, and $u_j\to0$ pointwise \emph{everywhere}.

\emph{No convergence in capacity.}  The computation
\eqref{eq:ch25-p1-capacity-open}, carried out at the centre $a_j$ instead of
$0$---the two situations differ by an automorphism of the pair
$(\mathbf P^1,\theta)$ given by the translation $T\mapsto T-a_j$, which
preserves the standard model and its Gauss point---gives
\[
 \Capa_\theta\bigl(\{|u_j|>\tfrac12\}\bigr)
 =\Capa_\theta\bigl(\{\varphi_j<-\tfrac12\}\bigr)
 =\min\{1,2\}=1=V
\]
for every $j$.  Thus $\Capa_\theta(\{|u_j-0|>1/2\})=V$ does not tend to
$0$, and $u_j$ does not converge to $0$ in capacity.  This proves the
second half of
\Cref{prop:ch25-convergence-comparison}\,\cref{itm:ch25-cc:none}.

\emph{What the example separates.}  The functions $u_j$ are
$\theta$-psh, uniformly bounded, and converge pointwise everywhere to a
$\theta$-psh limit; they even converge in $L^1(\theta)$, since
$\int u_j\,\theta=u_j(\zeta_{0,1})=0$ for every $j$.  None of this
implies convergence in capacity.  What fails is monotonicity: the sequence
is neither increasing nor decreasing, and
\Cref{prop:decreasing-capacity-convergence} is therefore unavailable.  The
example also shows that the hypothesis of
\Cref{lem:ch25-joint-limit}---convergence in capacity of the test
functions---cannot be weakened to pointwise convergence.
\end{example}

\begin{example}[Comparison, maximum and domination, with numbers]
\label{ex:ch25-comparison-numbers}
Return to $\Gamma=[0,4]$ with $m_0=m_1=1$, $V=2$, and consider the two
affine potentials
\[
 \varphi(x)=-1+\tfrac x4,\qquad \psi(x)=-\tfrac x4,
\]
both in $\mathcal Q_\Gamma$ and hence, after composing with $p$, in
$\mathcal U_\theta$.  By \eqref{eq:ch25-segment-measure},
\[
 \MA_\theta(\varphi)=\tfrac54\delta_{x_0}+\tfrac34\delta_{x_4},
 \qquad
 \MA_\theta(\psi)=\tfrac34\delta_{x_0}+\tfrac54\delta_{x_4}.
\]

\emph{The comparison principle.}  One has $\varphi<\psi$ exactly for
$x<2$, so $\{\varphi<\psi\}=p^{-1}([0,2))$ and
\[
 \int_{\{\varphi<\psi\}}\MA_\theta(\psi)=\tfrac34
 \ \leq\ \tfrac54=\int_{\{\varphi<\psi\}}\MA_\theta(\varphi),
\]
in accordance with \eqref{eq:bounded-comparison-principle}, with a strict
gap of $1/2$.

\emph{The maximum principle and the contact set.}  The maximum
$\max\{\varphi,\psi\}$ has a kink at $x=2$, where its value is
$-1/2$, and
\[
 \MA_\theta(\max\{\varphi,\psi\})
 =\tfrac34\delta_{x_0}+\tfrac12\delta_{x_2}+\tfrac34\delta_{x_4},
\]
of total mass $2$.  On $\{\varphi>\psi\}=p^{-1}((2,4])$ this measure is
$\tfrac34\delta_{x_4}$, which is exactly
$\one_{\{\varphi>\psi\}}\MA_\theta(\varphi)$, as
\eqref{eq:bounded-maximum-principle} requires; symmetrically on
$\{\psi>\varphi\}$.  On the contact set
$\{\varphi=\psi\}=p^{-1}(\{x_2\})$, however, the maximum carries an atom of
mass $\tfrac12$ while both $\MA_\theta(\varphi)$ and
$\MA_\theta(\psi)$ vanish there.  This is the phenomenon of
\Cref{rem:ch25-max-sets}, and it is why the proof of
\Cref{cor:bounded-comparison-principle} shifts one potential by
$\varepsilon$.

\emph{The capacity comparison, and its loss.}  By
\Cref{prop:ch25-segment-capacities} applied to $A=[0,2)$, whose mass under
\eqref{eq:ch25-segment-measure} is $m_0+f'(2^-)$,
\[
 \Capa_\theta\bigl(\{\varphi<\psi\}\bigr)
 =1+\min\Bigl\{1,\tfrac12\Bigr\}=\tfrac32 .
\]
On the other hand, for $0<t<1$ and $x\in[0,4]$,
\[
 \varphi(x)<(1-t)\psi(x)+t
 \iff \frac{(2-t)x}{4}<1+t
 \iff x<x_t:=\frac{4(1+t)}{2-t},
\]
so that $\{\varphi<(1-t)\psi+t\}=p^{-1}\bigl([0,4]\cap[0,x_t)\bigr)$.  Now
$x_t\leq4$ exactly when $t\leq\tfrac12$, so the right-hand side of
\eqref{eq:capacity-comparison} equals $t^{-1}\cdot\tfrac54$ for
$t\leq\tfrac12$, the atom at $x_4$ being excluded, and
$t^{-1}\cdot2$ for $t>\tfrac12$.  Minimizing over $t\in(0,1)$ gives
the bound $2$, approached as $t\to1^-$.  Thus
\[
 \tfrac32=\Capa_\theta\{\varphi<\psi\}\ \leq\ 2 ,
\]
a loss of a factor $4/3$.  The factor $t^{-n}$ is what costs the
estimate its sharpness, and no attempt is made in this book to remove it.

\emph{The domination principle, and an independent check.}  Keep
$\psi(x)=-x/4$ and let $\chi$ be \emph{any} bounded $\theta$-psh
function on $C^{\An}$.  Since $\MA_\theta(\psi)$ is carried by the two
points $x_0,x_4$, the hypothesis of
\Cref{prop:bounded-capacity-domination} reads
\[
 \chi(x_0)\leq\psi(x_0)=0
 \qquad\text{and}\qquad
 \chi(x_4)\leq\psi(x_4)=-1 ,
\]
a condition at two points; the conclusion is that $\chi\leq\psi$ on all of
$C^{\An}$.  On this example the conclusion can be verified by hand, which
is a useful check on the signs.  Indeed $\chi|_\Gamma$ is convex by the curve psh Laplacian criterion, and $\psi|_\Gamma$ is affine, so
$\chi|_\Gamma-\psi|_\Gamma$ is convex on $[0,4]$ and hence attains its
maximum at an endpoint, where it is $\leq0$; thus $\chi\leq\psi$ on
$\Gamma$.  Off $\Gamma$ use the retraction inequality
\eqref{eq:ch21-retraction-inequality}: $\chi\leq\chi\circ p$, and
$\chi\circ p\leq\psi\circ p=\psi$ because the inequality holds on
$\Gamma$ and $\psi$ is already pulled back from $\Gamma$.  So
$\chi\leq\psi$ everywhere, as the domination principle predicts.

\emph{What happens when the hypothesis fails.}  Take $\chi=\varphi$, so
that $\chi(x_0)=-1\leq0$ but $\chi(x_4)=0>-1=\psi(x_4)$.  The
hypothesis fails at the atom $x_4$, which carries mass $5/4$, and the
conclusion fails too: $\varphi>\psi$ on $p^{-1}((2,4])$.  Thus the
two-point condition is not merely sufficient; on this graph it is also
necessary.
\end{example}

\section{Scope: what each hypothesis is doing}
\label{sec:ch25-scope}

\subsection{Hypothesis use, statement by statement}
\label{sec:ch25-hypothesis-use}

\begin{center}
\small
\begin{tabular}{p{.30\textwidth}|p{.31\textwidth}|p{.31\textwidth}}
\toprule
\textbf{Statement}&\textbf{Uses}&\textbf{Does not use}\\
\midrule
\Cref{prop:capacity-elementary},
\Cref{cor:ch25-outer-subadditivity}
 &positivity and mass $V$ of $\MA_\theta$; countable additivity of a
  measure
 &discreteness of $v_K$; $\operatorname{char}k=0$\\
\Cref{prop:ch25-choquet-axioms}
 &compactness of $X^{\An}$; outer regularity of the definition
 &strong subadditivity; capacitability\\
\Cref{lem:divisorial-positive-capacity}
 &ampleness of $\{\theta\}$
  (\Cref{prop:ch20-positive-representative}); Nakai--Moishezon on
  $\mathcal X_s$; density of divisorial points
 &$\operatorname{char}k=0$, except through the existence of the model\\
\Cref{lem:ma-capacity-bound}
 &semipositivity of $\theta$; multilinearity
 &any regularization\\
\Cref{lem:ch25-uniform-integrability},
\Cref{lem:ch25-normalized-bound}
 &integration by parts \eqref{eq:model-ibp};
  \Cref{thm:bfj-skeletal-equicontinuity};
  \Cref{prop:ch21-sup-on-skeleton}
 &boundedness of the tested potential $u$\\
\Cref{lem:fundamental-capacity-estimate}
 &\eqref{eq:model-dirichlet-negativity} and the Hodge-type negativity
  behind it (\Cref{thm:ch20-mixed-hodge}); boundedness by $M$
 &quasicontinuity; capacity\\
\Cref{thm:bounded-quasicontinuity},
\Cref{cor:psh-sequential-regularization}
 &\Cref{thm:bfj-decreasing-regularization}, hence
  $\operatorname{char}k=0$; weighted decreasing-net continuity in
  \Cref{thm:bounded-ma-extension}
 &metrizability of $X^{\An}$; any countable cofinal family\\
\Cref{thm:bounded-maximum-principle}
 &the continuous maximum identity in
  \Cref{prop:model-maximum-principle}; \Cref{lem:ch25-joint-limit}
 &continuity of envelopes\\
\Cref{cor:bounded-comparison-principle},
\Cref{prop:bounded-capacity-domination}
 &total mass $V$; \eqref{eq:ma-convex-lower-bound};
  \Cref{lem:divisorial-positive-capacity}
 &finite energy; orthogonality; continuity of potentials\\
\Cref{prop:ch25-curve-pushforward}
 &\Cref{prop:ch21-curve-retraction}, valid only for curves
 &anything in dimension $\geq2$\\
\bottomrule
\end{tabular}
\end{center}

\subsection{Later extensions}
\label{sec:ch25-later}

\emph{(E1) Arbitrary complete nontrivially valued fields.}
The definition \eqref{eq:capacity-definition} and the elementary
set-function properties make sense whenever the required positive bounded
operator of mass $V$ is available.  The integral estimates additionally
use the stated mixed-product, integration-by-parts, Dirichlet-negativity
and skeletal inputs.  These are not consequences of positivity and mass
alone.  The current model construction and positive-representative
argument belong to the standing discrete valuation framework.
\cite{BE21} supplies broader continuous-metric intersection and volume
results, not by itself a transplantation of every singular bounded-psh
statement in this chapter.

\emph{(E2) Positive residue characteristic.}
The elementary capacity calculus uses no multiplier ideal once its bounded
operator is supplied.  The present proofs of the estimates and divisorial
positivity retain the imported intersection, SNC and positivity framework;
they are not asserted over every positive-characteristic field merely
because they do not mention multiplier ideals.  Quasicontinuity,
sequential regularization and bounded locality also use decreasing model
regularization.  Later envelope-continuity results cited in
\Cref{sec:ch23-later} do not alone provide all these inputs.
An extension requires the appropriate reconstruction and regularization,
the bounded operator with weighted continuity, positive representatives,
and the remaining geometric hypotheses.  No unconditional
positive-characteristic extension is asserted here.

\emph{(E3) Singular $X$.}  Smoothness of $X$ is used here only through
the availability of SNC models and the equicontinuity estimate
\Cref{thm:bfj-skeletal-equicontinuity}, that is, in the normalized bound and regularization arguments.  The elementary
set-function calculus uses only the positive bounded operator, but the
standing construction of that operator and the divisorial positivity
argument retain their geometric model hypotheses.

\emph{(E4) The finite-energy capacity.}  A different set function, adapted
to $\cE^1$, is used in \Cref{chap:bfj-finite-energy-measures} and, in the
trivially valued theory, in \Cref{chap:bj-strong-topologies}.  None of the
estimates of this chapter is asserted for it, and the symbol
$\Capa_\theta$ is never used for it.

\emph{(E5) Curves: a stronger theory exists, for a different capacity.}
In dimension one the relative capacity \eqref{eq:ch25-thuillier-capacity} of
\cite{Thu05} satisfies strong subadditivity, Choquet capacitability, and the
characterization of null sets as exactly the polar sets
(\Cref{thm:ch25-thuillier-capacity},
\Cref{thm:ch25-thuillier-polar})---the three items this chapter declines for
$\Capa_\theta$.  It is not an extension of $\Capa_\theta$ but a different
functional, relative to an open set rather than global and defined by
energies rather than masses, and \Cref{rem:ch25-thuillier-vs-bfj} lists what
does and does not transfer.  We record it here because a reader who needs
capacitability on a curve should use \cite{Thu05} and not attempt to repair
the arguments of \Cref{sec:ch25-capacity}.

\section{Dictionary}
\label{sec:ch25-dictionary}

\subsection{Sets, functions, and the statements that control them}
\label{sec:ch25-dict-objects}

\begin{center}
\small
\begin{tabular}{p{.29\textwidth}|p{.32\textwidth}|p{.31\textwidth}}
\toprule
\textbf{Analytic notion}&\textbf{Capacity-theoretic content}
 &\textbf{Established at}\\
\midrule
set carrying no mass for any normalized bounded potential
 &$\Capa_\theta(E)=0$; negligible
 &\Cref{def:bfj-capacity},
  \Cref{prop:ch25-negligible-sigma-ideal}\\
polar set $\{v=-\infty\}$ of a $\theta$-psh function
 &negligible, with rate $O(1/t)$ on sublevels
 &\Cref{prop:ch25-sublevel-capacity},
  \Cref{cor:ch25-polar-negligible}\\
countable union of polar sets
 &negligible
 &\Cref{thm:ch25-pluripolar-negligible}\\
divisorial point
 &positive capacity, computable as an intersection number
 &\Cref{lem:divisorial-positive-capacity},
  \eqref{eq:ch25-divisorial-atom}\\
open set of small capacity
 &exceptional set for quasicontinuity
 &\Cref{thm:bounded-quasicontinuity}\\
decreasing model net
 &convergence in capacity
 &\Cref{prop:decreasing-capacity-convergence}\\
decreasing model \emph{sequence}
 &available for every $\theta$-psh function
 &\Cref{cor:psh-sequential-regularization}\\
strict inequality locus $\{\varphi>\psi\}$
 &locality of $\MA_\theta$
 &\eqref{eq:bounded-maximum-principle}\\
contact locus $\{\varphi=\psi\}$
 &no general equality with either operand; extra mass may appear
 &\Cref{rem:ch25-max-sets},
  \Cref{ex:ch25-comparison-numbers}\\
sublevel locus $\{\varphi<\psi\}$
 &comparison and capacity comparison
 &\eqref{eq:bounded-comparison-principle},
  \eqref{eq:capacity-comparison}\\
tube $p^{-1}(A)$ over a metric graph
 &convex optimization over admissible graph potentials
 &\Cref{cor:ch25-curve-capacity-formula}\\
regularization defect $\{\varphi<\varphi^*\}$
 &contains no quasimonomial point; negligibility not proved
 &\Cref{prop:ch25-regularization-defect}\\
\bottomrule
\end{tabular}
\end{center}

\subsection{Reading the dictionary backwards}
\label{sec:ch25-dict-backward}

Four entries are one-way, and it is worth saying which.  \emph{From
negligible to pluripolar} there is no arrow: only the forward implication is
proved, and the missing ingredient is described after
\Cref{cor:ch25-polar-countable}.  \emph{From convergence in capacity to
pointwise convergence} there is an arrow only for sequences and only after
passing to a subsequence
(\Cref{prop:ch25-convergence-comparison}\,\cref{itm:ch25-cc:seq}), and no analogous subnet theorem is proved here
(\Cref{rem:ch25-nets-in-capacity}).  \emph{From pointwise
convergence to convergence in capacity} there is no arrow even for uniformly
bounded $\theta$-psh sequences
(\Cref{ex:ch25-pointwise-not-capacity}).  \emph{From equality of
Monge--Amp\`ere measures to equality of potentials} there is no arrow within
this chapter (\Cref{rem:ch25-no-uniqueness}); the arrow exists, but it is
variational and belongs to
\Cref{chap:bfj-differentiability,chap:bfj-finite-energy-measures}.

\subsection{Numbers to remember}
\label{sec:ch25-dict-numbers}

The following values were computed in full and are the test suite for the
rest of \Cref{part:bfj}.

\begin{center}
\small
\begin{tabular}{p{.36\textwidth}|p{.24\textwidth}|p{.32\textwidth}}
\toprule
\textbf{Quantity}&\textbf{Value}&\textbf{Where}\\
\midrule
$\Capa_\theta(\{x_0\})$, segment of length $1$, $m_0=m_1=1$
 &$2=V$
 &\Cref{ex:ch25-two-vertex-full-capacity}\\
$\Capa_\theta(p^{-1}(\{x_0\}))$, length $4$
 &$5/4$
 &\Cref{ex:ch25-segment-capacities}\\
$\Capa_\theta(p^{-1}(\{x_1\}))$, $\Capa_\theta(p^{-1}(\{x_2\}))$
 &$4/3$, $1$
 &\Cref{ex:ch25-segment-capacities}\\
$\Capa_\theta(p^{-1}([1,2]))+\Capa_\theta(p^{-1}([2,3]))$
 versus $\Capa_\theta(p^{-1}([1,3]))$
 &$3$ versus $2$
 &\Cref{ex:ch25-subadditivity-strict}\\
capacity of a midpoint, segment of length $\ell$
 &$4/\ell$ for $\ell\geq2$
 &\Cref{ex:ch25-segment-capacities}\\
$\Capa_\theta(\{\varphi_0\leq-t\})$ on $\mathbf P^1$
 &$\min\{1,1/t\}$
 &\eqref{eq:ch25-p1-capacity}\\
sharp constant in \eqref{eq:ch25-sublevel-capacity} on $\mathbf P^1$
 &$C_\theta+nV=1$
 &\Cref{ex:ch25-pluripolar-computed}\\
slack in \eqref{eq:fundamental-capacity-estimate} on the segment
 &factor $6$ versus $5/4$
 &\Cref{ex:ch25-fundamental-numbers}\\
slack in \eqref{eq:capacity-comparison} on the segment
 &$2$ versus $3/2$
 &\Cref{ex:ch25-comparison-numbers}\\
extra mass of $\MA_\theta(\max\{\varphi,\psi\})$ on the contact set
 &$1/2$
 &\Cref{ex:ch25-comparison-numbers}\\
$\Capa_\theta(\{|u_j|>1/2\})$ for the residue-disc potentials
 &$1=V$, for every $j$
 &\Cref{ex:ch25-pointwise-not-capacity}\\
\bottomrule
\end{tabular}
\end{center}

\section*{Notes and further reading}

The Monge--Amp\`ere capacity \eqref{eq:capacity-definition}, the two
estimates of \Cref{sec:ch25-estimates}, quasicontinuity
\Cref{thm:bounded-quasicontinuity} and the sequential reduction
\Cref{cor:psh-sequential-regularization} follow
\cite[\S4]{BFJ15}; the maximum principle
\Cref{thm:bounded-maximum-principle} and the comparison principle
\Cref{cor:bounded-comparison-principle} are
\cite[Theorem~5.1 and Corollaries~5.2--5.3]{BFJ15}, and
\eqref{eq:capacity-comparison} is the bounded case of
\cite[Lemma~8.3]{BFJ15}.  The model input
\Cref{prop:ch20-positive-representative}, used in
\Cref{lem:divisorial-positive-capacity}, is \cite[\S5]{BFJ16}.
The extended-real-valued statement
\Cref{thm:ch25-psh-quasicontinuity} is also
\cite[Proposition~4.6]{BFJ15}.  Its truncation proof here uses
\Cref{prop:ch25-sublevel-capacity} and
\Cref{thm:bounded-quasicontinuity}, independently of the later
finite-energy theory.  The bounded domination conclusion is derived here
from capacity comparison and divisorial order; its broader bounded scope
is not attributed verbatim to \cite[Lemma~8.4]{BFJ15}.

\paragraph{The complex ancestry.}  The construction is a transposition of
Bedford and Taylor's capacity for plurisubharmonic functions \cite{BT82},
where \eqref{eq:capacity-definition} is taken over psh functions with values
in $[-1,0]$ on a bounded domain, and of its compact K\"ahler counterpart
due to Guedj and Zeriahi \cite{GZ05}.  Three points of the transposition are
worth recording.  First, the Cauchy--Schwarz bootstrap
\Cref{lem:fundamental-capacity-estimate} is the exact analogue of the
Chern--Levine--Nirenberg iteration used by Ko\l odziej \cite{Kol98} and by
Guedj--Zeriahi \cite{GZ05,GZ07} to compare integrals against different
Monge--Amp\`ere measures; the exponent $2^{-n}$ is the same, and its
origin---one Cauchy--Schwarz per factor---is the same.  Second, the complex
theory has at its disposal a converse to
\Cref{cor:ch25-polar-negligible}, namely that a set of zero capacity is
pluripolar, proved by way of extremal functions and Josefson's theorem; that
converse is not available here, for the structural reason recorded after
\Cref{cor:ch25-polar-countable}.  Third, in the complex setting the sets
$\{\varphi>\psi\}$ are already reasonably tame, whereas here the
absence of an assumed countably cofinal model system leads to the
systematic use of nets, and the
extraction of a sequence in
\Cref{thm:bounded-quasicontinuity} is a theorem rather than a convenience;
this is the point at which the non-Archimedean argument genuinely differs
from its model.

\paragraph{Potential theory on curves.}  The whole of
\Cref{sec:ch25-thuillier} is Thuillier's \cite{Thu05}, and no proof there is
ours except the verification carried out in
\Cref{ex:ch25-thuillier-disc}.  His Chapter~3 constructs, for a smooth
strictly $K$-analytic curve, the sheaves of harmonic and subharmonic
functions, the operator $\ddc$ and its currents, the Hilbert space
$\mathrm W^1$ of finite Dirichlet energy---taken up in
\Cref{sec:ch26-thuillier}---and then, in \S3.6, the relative capacity
\eqref{eq:ch25-thuillier-capacity}, its outer regularization
\eqref{eq:ch25-thuillier-outer}, the equilibrium potentials, and the
capacitary characterization of polar sets.  He follows Klimek's presentation
of the complex theory \cite{Kl91}, simplified: in particular
\emph{no quasicontinuity of subharmonic functions is used anywhere}.  That
is the sharpest structural contrast with the present chapter, whose central
technical result is quasicontinuity
(\Cref{thm:bounded-quasicontinuity}) and which needs it because the
regularization available in dimension $\geq2$ is by nets of model
functions rather than by solving a Dirichlet problem.  Where the general
theory here specializes on a curve to something Thuillier proves
unconditionally, the fact is recorded at the point of contact:
\Cref{rem:ch25-curve-choquet-unconditional} for the Choquet axioms,
\Cref{rem:ch25-curve-converse} for the converse to
\Cref{cor:ch25-polar-negligible}, \Cref{rem:ch25-thuillier-vs-bfj} for the
strengthening of \Cref{lem:divisorial-positive-capacity} to the type
dichotomy, and \Cref{rem:ch25-thuillier-determination} for the determination
of subharmonic functions by classical points.  In the scope where $K$ is
algebraically closed, complete, and nontrivially valued and $C$ is the
analytification of a smooth proper algebraic curve, the identification of
continuous Thuillier-subharmonic functions with continuous current-positive
functions is \Cref{thm:wanner-curve-comparison}; their further identification
with the first ordinary-smooth CLD decreasing-limit stage is
\Cref{thm:p063-current-psh-ordinary-regularization}.  These identifications
are not being asserted under the broader standing hypotheses of this
subsection.  The normalizations are audited in
\Cref{rem:ch25-thuillier-normalization}, the
only discrepancy being a factor $2$ in the \emph{archimedean} model of the
Dirichlet form.

\paragraph{Negligible sets.}  Bedford and Taylor's theorem that negligible
sets---in their sense, regularization defects---are pluripolar
\cite{BT82} has no proof here.  What replaces it is
\Cref{prop:ch25-regularization-defect}: the defect misses every skeleton, a
statement with no complex analogue, obtained from
\Cref{thm:ch23-upper-envelope}.  The two statements are not comparable, and
the reader should not substitute one for the other.  Cegrell's work on
domains of definition of the complex Monge--Amp\`ere operator
\cite{Ceg98,Ceg04} is the source of the systematic distinction between the
bounded and the finite-energy regimes which
\Cref{sec:ch25-bounded-regime} imposes here.

\paragraph{Where the chapter is used.}  Quasicontinuity is used in
\Cref{chap:bfj-energy} to integrate psh functions against Monge--Amp\`ere
measures of other psh functions, and in
\Cref{chap:bfj-finite-energy-measures} to identify the domain of the energy
pairing.  The comparison principle in its capacity form
\eqref{eq:capacity-comparison} is used in
\Cref{cor:e1-capacity-tail} and in the Ko\l odziej-type estimate
\Cref{lem:bfj-kolodziej-sublevel} of \Cref{chap:bfj-continuity}.  The
maximum principle \eqref{eq:bounded-maximum-principle} is the basis of the
non-pluripolar product \Cref{def:nonpluripolar-ma}.  The abstract calculus
imported from \Cref{sec:capacity-preview} is used again, for a different
capacity, in \Cref{chap:bj-strong-topologies}.

\paragraph{What has not been claimed.}  No converse from negligible to
pluripolar; no strong subadditivity and hence no capacitability; no
negligibility of regularization defects; no verbatim use of a bounded
$M$-estimate on an unbounded potential; and no uniqueness statement for the
Monge--Amp\`ere equation.  Each of these is flagged at the point where a
reader might expect it.
